% Generated mechanically from the frozen stacks-more-morphisms.tex target.
% Do not edit this generated file; edit the bound locale source instead.

% OK, start here.
%

\setcounter{chapter}{105}
\chapter{叠态射的进一步讨论}



\phantomsection
\label{stacks-more-morphisms-section-phantom}


\section{引言}
\label{stacks-more-morphisms-section-introduction}

\noindent
本章继续研究代数叠态射的性质。对于具有可表示对角态射的
拟分离代数叠，可参见 \cite{LM-B}。




\section{约定与语言上的简化}
\label{stacks-more-morphisms-section-conventions}

\noindent
我们继续采用《叠的性质》第
\ref{stacks-properties-section-conventions} 节中引入的约定与
语言上的简化说法。






\section{增厚}
\label{stacks-more-morphisms-section-thickenings}

\noindent
下面的术语或许并非完全标准，但使用起来很方便。
若 $\mathcal{Y}$ 是代数叠 $\mathcal{X}$ 的闭子叠，
则态射 $\mathcal{Y} \to \mathcal{X}$ 是可表示的。

\begin{definition}
\label{stacks-more-morphisms-definition-thickening}
增厚。
\begin{enumerate}
\item 若 $\mathcal{X}$ 是代数叠 $\mathcal{X}'$ 的闭子叠，且相应的
拓扑空间相等，则称代数叠 $\mathcal{X}'$ 是代数叠 $\mathcal{X}$ 的一个
{\it 增厚}。
\item 给定两个增厚 $\mathcal{X} \subset \mathcal{X}'$ 与
$\mathcal{Y} \subset \mathcal{Y}'$，所谓{\it 增厚的态射}，是指代数叠
的态射 $f' : \mathcal{X}' \to \mathcal{Y}'$，使得
$f'|_\mathcal{X}$ 经过闭子叠 $\mathcal{Y}$ 分解。在这种情形下，令
$f = f'|_\mathcal{X} : \mathcal{X} \to \mathcal{Y}$，并称
$(f, f') : (\mathcal{X} \subset \mathcal{X}') \to
(\mathcal{Y} \subset \mathcal{Y}')$ 是一个增厚的态射。
\item 设 $\mathcal{Z}$ 为代数叠。类似地定义
{\it $\mathcal{Z}$ 上的增厚}以及
{\it $\mathcal{Z}$ 上增厚的态射}。
这意味着代数叠 $\mathcal{X}'$ 和 $\mathcal{Y}'$ 都配备到
$\mathcal{Z}$ 的结构态射，且 $f'$ 位于一个适当的代数叠
$2$-交换图中。
\end{enumerate}
\end{definition}

\noindent
设 $\mathcal{X} \subset \mathcal{X}'$ 为代数叠的增厚。
设 $U'$ 为概形，并设 $U' \to \mathcal{X}'$ 为满光滑态射。
令 $U = \mathcal{X} \times_{\mathcal{X}'} U'$，便得到增厚的态射
$$
(U \subset U') \longrightarrow (\mathcal{X} \subset \mathcal{X}')
$$
且 $U \to \mathcal{X}$ 是满光滑态射。我们经常可以从增厚
$U \subset U'$ 的相应性质推出增厚
$\mathcal{X} \subset \mathcal{X}'$ 的性质。
有时，作为语言上的简化，若态射 $\mathcal{X} \to \mathcal{X}'$
是闭浸入并诱导双射 $|\mathcal{X}| \to |\mathcal{X}'|$，
也称它为一个增厚。

\begin{lemma}
\label{stacks-more-morphisms-lemma-thickening}
设 $i : \mathcal{X} \to \mathcal{X}'$ 为代数叠的态射。
下列条件等价：
\begin{enumerate}
\item $i$ 是代数叠的增厚（采用上面的简化说法）；
\item $i$ 可由代数空间表示，并且是《叠的性质》第
\ref{stacks-properties-section-properties-morphisms} 节意义下的增厚。
\end{enumerate}
在这种情形下，$i$ 是闭浸入及万有同胚。
\end{lemma}

\begin{proof}
由《空间态射的进一步讨论》中的引理
\ref{spaces-more-morphisms-lemma-descending-property-thickening} 和
\ref{spaces-more-morphisms-lemma-base-change-thickening}，
代数空间的态射为（一阶）增厚这一性质 $P$ 在基底上是 fpqc 局部的，
且在基变换下稳定。因此，《叠的性质》第
\ref{stacks-properties-section-properties-morphisms} 节的讨论确实适用。
至此，(1) 与 (2) 的等价性来自 $P = P_1 + P_2$，其中 $P_1$
是成为闭浸入的性质，而 $P_2$ 是成为满态射的性质。
% P12-ERR-0062: see ../control/ERRATA_CANDIDATES.jsonl
（严格说来，读者还应参阅《空间态射的进一步讨论》中的定义
\ref{spaces-more-morphisms-definition-thickening}、
《叠的性质》中的定义 \ref{stacks-properties-definition-immersion}
及其后的讨论、《空间的态射》中的引理
\ref{spaces-morphisms-lemma-surjective-representable}，以及
《叠的性质》第 \ref{stacks-properties-section-surjective} 节，
以确认所有这些概念彼此吻合。）最后一个断言由前述内容立即可得。
\end{proof}











\noindent
下文将不再另加说明地使用这个引理。利用该引理所用的同一组参考文献，
即《空间态射的进一步讨论》中的引理
\ref{spaces-more-morphisms-lemma-descending-property-thickening} 和
\ref{spaces-more-morphisms-lemma-base-change-thickening}，
可以如下定义一阶增厚。

\begin{definition}
\label{stacks-more-morphisms-definition-first-order-thickening}
若 $\mathcal{X}$ 是代数叠 $\mathcal{X}'$ 的闭子叠，且
$\mathcal{X} \to \mathcal{X}'$ 是《叠的性质》第
\ref{stacks-properties-section-properties-morphisms} 节意义下的一阶增厚，
则称代数叠 $\mathcal{X}'$ 是代数叠 $\mathcal{X}$ 的一个
{\it 一阶增厚}。
\end{definition}

\noindent
若 $(U \subset U') \to (\mathcal{X} \subset \mathcal{X}')$
是如上的概形光滑覆盖，那么这恰好意味着 $U \subset U'$
是一阶增厚。下面陈述必需的几个引理。

\begin{lemma}
\label{stacks-more-morphisms-lemma-base-change-thickening}
设 $\mathcal{Y} \subset \mathcal{Y}'$ 为代数叠的增厚。
设 $\mathcal{X}' \to \mathcal{Y}'$ 为代数叠的态射，并令
$\mathcal{X} = \mathcal{Y} \times_{\mathcal{Y}'} \mathcal{X}'$。
则
$(\mathcal{X} \subset \mathcal{X}') \to (\mathcal{Y} \subset \mathcal{Y}')$
是增厚的态射。若 $\mathcal{Y} \subset \mathcal{Y}'$ 是一阶增厚，
则 $\mathcal{X} \subset \mathcal{X}'$ 也是一阶增厚。
\end{lemma}

\begin{proof}
参见上面的讨论、《叠的性质》第
\ref{stacks-properties-section-properties-morphisms} 节，以及
《空间态射的进一步讨论》中的引理
\ref{spaces-more-morphisms-lemma-base-change-thickening}。
\end{proof}

\begin{lemma}
\label{stacks-more-morphisms-lemma-composition-thickening}
若 $\mathcal{X} \subset \mathcal{X}'$ 与
$\mathcal{X}' \subset \mathcal{X}''$ 都是代数叠的增厚，
则 $\mathcal{X} \subset \mathcal{X}''$ 也是代数叠的增厚。
\end{lemma}

\begin{proof}
参见上面的讨论、《叠的性质》第
\ref{stacks-properties-section-properties-morphisms} 节，以及
《空间态射的进一步讨论》中的引理
\ref{spaces-more-morphisms-lemma-composition-thickening}
\end{proof}

\begin{example}
\label{stacks-more-morphisms-example-reduction-thickening}
设 $\mathcal{X}'$ 为代数叠。则 $\mathcal{X}'$ 是其约化
$\mathcal{X}'_{red}$ 的一个增厚，参见《叠的性质》中的定义
\ref{stacks-properties-definition-reduced-induced-stack}。
此外，若 $\mathcal{X} \subset \mathcal{X}'$ 是代数叠的增厚，则
$\mathcal{X}'_{red} = \mathcal{X}_{red} \subset \mathcal{X}$。
换言之，$\mathcal{X} = \mathcal{X}'_{red}$ 当且仅当
$\mathcal{X}$ 是约化代数叠。
\end{example}

\begin{lemma}
\label{stacks-more-morphisms-lemma-reduced-diagonal}
设 $(f, f') : (\mathcal{X} \subset \mathcal{X}') \to
(\mathcal{Y} \subset \mathcal{Y}')$ 为代数叠增厚的态射。则
$\mathcal{X} \times_\mathcal{Y} \mathcal{X} \to
\mathcal{X}' \times_{\mathcal{Y}'} \mathcal{X}'$
是增厚，并且典范图表
$$
\xymatrix{
\mathcal{X} \ar[r]_-\Delta \ar[d] &
\mathcal{X} \times_\mathcal{Y} \mathcal{X} \ar[d] \\
\mathcal{X}' \ar[r]^-{\Delta'} &
\mathcal{X}' \times_{\mathcal{Y}'} \mathcal{X}'
}
$$
是笛卡儿图表。
\end{lemma}

\begin{proof}
由于 $\mathcal{X} \to \mathcal{Y}'$ 经过闭子叠
$\mathcal{Y}$ 分解，我们有
$\mathcal{X} \times_\mathcal{Y} \mathcal{X} =
\mathcal{X} \times_{\mathcal{Y}'} \mathcal{X}$。
因而
$\mathcal{X} \times_\mathcal{Y} \mathcal{X} \to
\mathcal{X}' \times_{\mathcal{Y}'} \mathcal{X}'$
同构于复合
$$
\mathcal{X} \times_{\mathcal{Y}'} \mathcal{X} \to
\mathcal{X} \times_{\mathcal{Y}'} \mathcal{X}' \to
\mathcal{X}' \times_{\mathcal{Y}'} \mathcal{X}'
$$
其中两个态射都是增厚的基变换，故都是增厚
（引理 \ref{stacks-more-morphisms-lemma-base-change-thickening}）。因此其复合也是增厚
（引理 \ref{stacks-more-morphisms-lemma-composition-thickening}）。
由于 $\mathcal{X} \to \mathcal{X}'$ 是单态射，将《叠的性质》中的引理
\ref{stacks-properties-lemma-monomorphism-diagonal}
应用于 $\mathcal{X} \to \mathcal{X}' \to \mathcal{Y}'$，
便得到引理的最后一个断言。
\end{proof}

\begin{lemma}
\label{stacks-more-morphisms-lemma-thickening-diagonals}
设 $(f, f') : (\mathcal{X} \subset \mathcal{X}') \to
(\mathcal{Y} \subset \mathcal{Y}')$ 为代数叠增厚的态射。
令 $\Delta : \mathcal{X} \to \mathcal{X} \times_\mathcal{Y} \mathcal{X}$ 与
$\Delta' : \mathcal{X}' \to \mathcal{X}' \times_{\mathcal{Y}'} \mathcal{X}'$
为相应的对角态射。则下列各性质由 $\Delta$ 满足，当且仅当由
$\Delta'$ 满足：
(a) 可由概形表示，(b) 仿射，(c) 满，(d) 拟紧，
(e) 万有闭，(f) 整，(g) 拟分离，(h) 分离，
(i) 万有单射，(j) 万有开，(k) 局部拟有限，
(l) 有限，(m) 非分歧，(n) 单态射，(o) 浸入，
(p) 闭浸入，以及 (q) 固有。
\end{lemma}

\begin{proof}
注意到
$$
(\Delta, \Delta') :
(\mathcal{X} \subset \mathcal{X}')
\longrightarrow
(\mathcal{X} \times_\mathcal{Y} \mathcal{X} \subset
\mathcal{X}' \times_{\mathcal{Y}'} \mathcal{X}')
$$
是增厚的态射（引理 \ref{stacks-more-morphisms-lemma-reduced-diagonal}）。此外，由
《叠的态射》中的引理 \ref{stacks-morphisms-lemma-properties-diagonal}，
$\Delta$ 与 $\Delta'$ 均可由代数空间表示。因此，结合《叠的性质》第
\ref{stacks-properties-section-properties-morphisms} 节的讨论，
再应用《空间态射的进一步讨论》中的引理
\ref{spaces-more-morphisms-lemma-thicken-property-morphisms}，
便得到情形 (a)、(b)、(c)、(d)、(e)、(f)、(g)、(h)、(i) 与 (j)。

\medskip\noindent
引理 \ref{stacks-more-morphisms-lemma-reduced-diagonal} 告诉我们
$\mathcal{X} = (\mathcal{X} \times_\mathcal{Y} \mathcal{X})
\times_{(\mathcal{X}' \times_{\mathcal{Y}'} \mathcal{X}')} \mathcal{X}'$。
此外，由前述《叠的态射》中的引理
\ref{stacks-morphisms-lemma-properties-diagonal}，
$\Delta$ 与 $\Delta'$ 都是局部有限型的。
% P12-ERR-0063: see ../control/ERRATA_CANDIDATES.jsonl
因此，对情形 (k)、(l)、(m)、(n)、(o)、(p) 与 (q)，
应用《空间态射的进一步讨论》中的引理
\ref{spaces-more-morphisms-lemma-properties-that-extend-over-thickenings}
即可得到结论。
\end{proof}

\noindent
作为推论，得到下面这个令人满意的结果。

\begin{lemma}
\label{stacks-more-morphisms-lemma-thickening-properties}
\begin{reference}
\cite[Theorem 2.2.5]{Conrad-moduli}
\end{reference}
设 $\mathcal{X} \subset \mathcal{X}'$ 为代数叠的增厚。则
\begin{enumerate}
\item $\mathcal{X}$ 是代数空间，当且仅当 $\mathcal{X}'$ 是代数空间；
\item $\mathcal{X}$ 是概形，当且仅当 $\mathcal{X}'$ 是概形；
\item $\mathcal{X}$ 是 DM，当且仅当 $\mathcal{X}'$ 是 DM；
\item $\mathcal{X}$ 是拟-DM，当且仅当 $\mathcal{X}'$ 是拟-DM；
\item $\mathcal{X}$ 是分离的，当且仅当 $\mathcal{X}'$ 是分离的；
\item $\mathcal{X}$ 是拟分离的，当且仅当 $\mathcal{X}'$ 是拟分离的；并且
% P12-ERR-0064: see ../control/ERRATA_CANDIDATES.jsonl
\item 在此处补充更多内容。
\end{enumerate}
\end{lemma}

\begin{proof}
在每一种情形中，我们都把问题化归为关于对角态射的问题，
再把引理 \ref{stacks-more-morphisms-lemma-thickening-diagonals} 应用于增厚的态射
$$
(\mathcal{X} \subset \mathcal{X}') \to
\left(\Spec(\mathbf{Z}) \subset \Spec(\mathbf{Z})\right)
$$
这里先通过《代数叠》中的定义 \ref{algebraic-definition-viewed-as}，
把 $\mathcal{X} \subset \mathcal{X}'$ 看成
$\Spec(\mathbf{Z})$ 上代数叠的增厚。

\medskip\noindent
情形 (1)。代数叠是代数空间，当且仅当其对角态射为单态射，
参见《叠的态射》中的引理 \ref{stacks-morphisms-lemma-hierarchy}
（这也可直接由《代数叠》中的命题
\ref{algebraic-proposition-algebraic-stack-no-automorphisms} 得到）。

\medskip\noindent
情形 (2)。由 (1)，可设 $\mathcal{X}$ 与 $\mathcal{X}'$
均为代数空间，随后应用《空间态射的进一步讨论》中的引理
\ref{spaces-more-morphisms-lemma-thickening-scheme}。

\medskip\noindent
情形 (3)--(6)。每种情形都对应于对角态射的一个条件，
参见《叠的态射》中的定义
\ref{stacks-morphisms-definition-separated} 和
\ref{stacks-morphisms-definition-absolute-separated}。
\end{proof}









\section{增厚的态射}
\label{stacks-more-morphisms-section-morphisms-thickenings}

\noindent
若 $(f, f') : (\mathcal{X} \subset \mathcal{X}') \to
(\mathcal{Y} \subset \mathcal{Y}')$ 是代数叠增厚的态射，
那么态射 $f$ 的性质往往由 $f'$ 继承。这里有若干种变体。

\begin{lemma}
\label{stacks-more-morphisms-lemma-thicken-property-morphisms}
设 $(f, f') : (\mathcal{X} \subset \mathcal{X}') \to
(\mathcal{Y} \subset \mathcal{Y}')$
为代数叠增厚的态射。则
\begin{enumerate}
\item $f$ 是仿射态射，当且仅当 $f'$ 是仿射态射；
\item $f$ 是满态射，当且仅当 $f'$ 是满态射；
% P12-ERR-0065: see ../control/ERRATA_CANDIDATES.jsonl
\item $f$ 是拟紧的，当且仅当 $f'$ 是拟紧的；
\item $f$ 是万有闭的，当且仅当 $f'$ 是万有闭的；
\item $f$ 是整态射，当且仅当 $f'$ 是整态射；
\item $f$ 是万有单射的，当且仅当 $f'$ 是万有单射的；
\item $f$ 是万有开的，当且仅当 $f'$ 是万有开的；
\item $f$ 是拟-DM 的，当且仅当 $f'$ 是拟-DM 的；
\item $f$ 是 DM 的，当且仅当 $f'$ 是 DM 的；
\item $f$ 是（拟）分离的，当且仅当 $f'$ 是（拟）分离的；
\item $f$ 是可表示的，当且仅当 $f'$ 是可表示的；
\item $f$ 可由代数空间表示，当且仅当 $f'$ 可由代数空间表示；
% P12-ERR-0066: see ../control/ERRATA_CANDIDATES.jsonl
\item 在此处补充更多内容。
\end{enumerate}
\end{lemma}

\begin{proof}
由引理 \ref{stacks-more-morphisms-lemma-thickening}，态射
$\mathcal{X} \to \mathcal{X}'$ 与 $\mathcal{Y} \to \mathcal{Y}'$
都是万有同胚。因此，关于
$|f| : |\mathcal{X}| \to |\mathcal{Y}|$ 的任何条件，
都等价于 $|f'| : |\mathcal{X}'| \to |\mathcal{Y}'|$ 上的相应条件；
在用任意态射 $\mathcal{Z}' \to \mathcal{Y}'$ 作基变换后，
同样成立。这证明了 (2)、(3)、(4)、(6)、(7)。

\medskip\noindent
在情形 (8)、(9)、(10)、(12) 中，可以像引理
\ref{stacks-more-morphisms-lemma-thickening-diagonals} 那样，把 $f$ 与 $f'$ 上的条件
转化为对角态射 $\Delta$ 与 $\Delta'$ 上的条件。参见
《叠的态射》中的定义 \ref{stacks-morphisms-definition-separated}
和引理 \ref{stacks-morphisms-lemma-hierarchy}。
因此，这些情形由引理 \ref{stacks-more-morphisms-lemma-thickening-diagonals} 得出。

\medskip\noindent
(11) 的证明。若 $f'$ 可表示，则 $f$ 也可表示，因为对概形 $T$
及态射 $T \to \mathcal{Y}$，有
$\mathcal{X} \times_\mathcal{Y} T =
\mathcal{X} \times_{\mathcal{X}'} (\mathcal{X}' \times_{\mathcal{Y}'} T)$，
而 $\mathcal{X} \to \mathcal{X}'$ 是闭浸入（因而可表示）。
反过来，设 $f$ 可表示，并设 $T' \to \mathcal{Y}'$ 为态射，
其中 $T'$ 是概形。则
$$
\mathcal{X} \times_{\mathcal{Y}}
(\mathcal{Y} \times_{\mathcal{Y}'} T') =
\mathcal{X} \times_{\mathcal{X}'}
(\mathcal{X}' \times_{\mathcal{Y}'} T') \to
\mathcal{X}' \times_{\mathcal{Y}'} T'
$$
是一个增厚（由引理 \ref{stacks-more-morphisms-lemma-base-change-thickening}），
且其源是概形。故由引理 \ref{stacks-more-morphisms-lemma-thickening-properties}，
其靶也是概形。

\medskip\noindent
在情形 (1) 与 (5) 中，若 $f$ 或 $f'$ 中任一个具有所述性质，
则由 (11)，$f$ 与 $f'$ 都可表示。在这种情形下，选取代数空间 $V'$
及满光滑态射 $V' \to \mathcal{Y}'$。令
$V = \mathcal{Y} \times_{\mathcal{Y}'} V'$、
$U' = \mathcal{X}' \times_{\mathcal{Y}'} V'$，以及
$U = \mathcal{X} \times_{\mathcal{Y}'} V'$。
于是，由《叠的性质》中的引理
\ref{stacks-properties-lemma-check-property-covering} 所给原则，
所需结论可由代数空间增厚的态射
$(U \subset U') \to (V \subset V')$ 的相应结论得到。
关于代数空间情形的相应结果，参见《空间态射的进一步讨论》中的引理
\ref{spaces-more-morphisms-lemma-thicken-property-morphisms}。
\end{proof}






\section{代数叠的无穷小形变}
\label{stacks-more-morphisms-section-deform}

\noindent
本节对应于《空间态射的进一步讨论》第
\ref{spaces-more-morphisms-section-deform} 节。

\begin{lemma}
\label{stacks-more-morphisms-lemma-flatness-morphism-thickenings-fp-over-ft}
考虑代数叠增厚的交换图表
$$
\xymatrix{
(\mathcal{X} \subset \mathcal{X}') \ar[rr]_{(f, f')} \ar[rd] & &
(\mathcal{Y} \subset \mathcal{Y}') \ar[ld] \\
& (\mathcal{B} \subset \mathcal{B}')
}
$$
假设
\begin{enumerate}
\item $\mathcal{Y}' \to \mathcal{B}'$ 是局部有限型的；
\item $\mathcal{X}' \to \mathcal{B}'$ 是平坦且局部有限表示的；
\item $f$ 是平坦的；并且
\item $\mathcal{X} = \mathcal{B} \times_{\mathcal{B}'} \mathcal{X}'$ 且
$\mathcal{Y} = \mathcal{B} \times_{\mathcal{B}'} \mathcal{Y}'$。
\end{enumerate}
则 $f'$ 是平坦的，并且对 $|f'|$ 的像中的每个
$y' \in |\mathcal{Y}'|$，态射 $\mathcal{Y}' \to \mathcal{B}'$
在 $y'$ 处平坦。
\end{lemma}

\begin{proof}
选取代数空间 $U'$ 及满光滑态射 $U' \to \mathcal{B}'$。
选取代数空间 $V'$ 及满光滑态射
$V' \to U' \times_{\mathcal{B}'} \mathcal{Y}'$。
选取代数空间 $W'$ 及满光滑态射
$W' \to V' \times_{\mathcal{Y}'} \mathcal{X}'$。
令 $U, V, W$ 分别为 $U', V', W'$ 沿
$\mathcal{B} \to \mathcal{B}'$ 的基变换。
于是，$f'$ 的平坦性等价于 $W' \to V'$ 的平坦性，
而已知 $W \to V$ 平坦。因此，可以把代数空间情形的引理
应用于代数空间增厚的图表
$$
\xymatrix{
(W \subset W') \ar[rr] \ar[rd] & & (V \subset V') \ar[ld] \\
& (U \subset U')
}
$$
参见《空间态射的进一步讨论》中的引理
\ref{spaces-more-morphisms-lemma-flatness-morphism-thickenings-fp-over-ft}。
关于 $\mathcal{Y}'/\mathcal{B}'$ 在 $|f'|$ 像中各点处平坦的断言，
也以同样方式得出。
\end{proof}

\begin{lemma}
\label{stacks-more-morphisms-lemma-deform-property-fp-over-ft}
考虑代数叠增厚的交换图表
$$
\xymatrix{
(\mathcal{X} \subset \mathcal{X}') \ar[rr]_{(f, f')} \ar[rd] & &
(\mathcal{Y} \subset \mathcal{Y}') \ar[ld] \\
& (\mathcal{B} \subset \mathcal{B}')
}
$$
% P12-ERR-0067: see ../control/ERRATA_CANDIDATES.jsonl
假设 $\mathcal{Y}' \to \mathcal{B}'$ 是局部有限型的，
$\mathcal{X}' \to \mathcal{B}'$ 是平坦且局部有限表示的，
$\mathcal{X} = \mathcal{B} \times_{\mathcal{B}'} \mathcal{X}'$，并且
$\mathcal{Y} = \mathcal{B} \times_{\mathcal{B}'} \mathcal{Y}'$。则
\begin{enumerate}
\item $f$ 平坦，当且仅当 $f'$ 平坦；
\label{stacks-more-morphisms-item-flat-fp-over-ft}
\item $f$ 是同构，当且仅当 $f'$ 是同构；
\label{stacks-more-morphisms-item-isomorphism-fp-over-ft}
\item $f$ 是开浸入，当且仅当 $f'$ 是开浸入；
\label{stacks-more-morphisms-item-open-immersion-fp-over-ft}
\item $f$ 是单态射，当且仅当 $f'$ 是单态射；
\label{stacks-more-morphisms-item-monomorphism-fp-over-ft}
\item $f$ 是局部拟有限的，当且仅当 $f'$ 是局部拟有限的；
\label{stacks-more-morphisms-item-quasi-finite-fp-over-ft}
\item $f$ 是 syntomic 的，当且仅当 $f'$ 是 syntomic 的；
\label{stacks-more-morphisms-item-syntomic-fp-over-ft}
\item $f$ 是光滑的，当且仅当 $f'$ 是光滑的；
\label{stacks-more-morphisms-item-smooth-fp-over-ft}
\item $f$ 是非分歧的，当且仅当 $f'$ 是非分歧的；
\label{stacks-more-morphisms-item-unramified-fp-over-ft}
\item $f$ 是 \'etale 的，当且仅当 $f'$ 是 \'etale 的；
\label{stacks-more-morphisms-item-etale-fp-over-ft}
\item $f$ 是有限的，当且仅当 $f'$ 是有限的；并且
\label{stacks-more-morphisms-item-finite-fp-over-ft}
% P12-ERR-0068: see ../control/ERRATA_CANDIDATES.jsonl
\item 在此处补充更多内容。
\end{enumerate}
\end{lemma}

\begin{proof}
在情形 (\ref{stacks-more-morphisms-item-flat-fp-over-ft}) 中，结论由引理
\ref{stacks-more-morphisms-lemma-flatness-morphism-thickenings-fp-over-ft} 得出。

\medskip\noindent
在情形 (\ref{stacks-more-morphisms-item-syntomic-fp-over-ft})、
(\ref{stacks-more-morphisms-item-smooth-fp-over-ft}) 中，可以采用引理
\ref{stacks-more-morphisms-lemma-flatness-morphism-thickenings-fp-over-ft} 证明中所用的方法。
具体而言，选取代数空间 $U'$ 及满光滑态射
$U' \to \mathcal{B}'$。选取代数空间 $V'$ 及满光滑态射
$V' \to U' \times_{\mathcal{B}'} \mathcal{Y}'$。
选取代数空间 $W'$ 及满光滑态射
$W' \to V' \times_{\mathcal{Y}'} \mathcal{X}'$。
令 $U, V, W$ 分别为 $U', V', W'$ 沿
$\mathcal{B} \to \mathcal{B}'$ 的基变换。
% P12-ERR-0069: see ../control/ERRATA_CANDIDATES.jsonl
于是，$f$（相应地，\ $f'$）具有该性质，等价于
$W' \to V'$（相应地，\ $W \to V$）具有该性质。
因此，可以把代数空间情形的引理应用于代数空间增厚的图表
$$
\xymatrix{
(W \subset W') \ar[rr] \ar[rd] & & (V \subset V') \ar[ld] \\
& (U \subset U')
}
$$
参见《空间态射的进一步讨论》中的引理
\ref{spaces-more-morphisms-lemma-deform-property-fp-over-ft}。

\medskip\noindent
在情形 (\ref{stacks-more-morphisms-item-unramified-fp-over-ft}) 与
(\ref{stacks-more-morphisms-item-etale-fp-over-ft}) 中，首先由对 $f$ 或 $f'$ 的假设看出，
$f$ 与 $f'$ 都是代数叠的 DM 态射，参见引理
\ref{stacks-more-morphisms-lemma-thicken-property-morphisms}。然后选取代数空间 $U'$
及满光滑态射 $U' \to \mathcal{B}'$。
选取代数空间 $V'$ 及满光滑态射
$V' \to U' \times_{\mathcal{B}'} \mathcal{Y}'$。
选取代数空间 $W'$ 及满 \'etale(!) 态射
$W' \to V' \times_{\mathcal{Y}'} \mathcal{X}'$。
令 $U, V, W$ 分别为 $U', V', W'$ 沿
$\mathcal{B} \to \mathcal{B}'$ 的基变换。
于是 $W \to V \times_\mathcal{Y} \mathcal{X}$ 也满且 \'etale。
% P12-ERR-0069: see ../control/ERRATA_CANDIDATES.jsonl
因此，$f$（相应地，\ $f'$）具有该性质，等价于
$W' \to V'$（相应地，\ $W \to V$）具有该性质。
故可以把代数空间情形的引理应用于代数空间增厚的图表
$$
\xymatrix{
(W \subset W') \ar[rr] \ar[rd] & & (V \subset V') \ar[ld] \\
& (U \subset U')
}
$$
参见《空间态射的进一步讨论》中的引理
\ref{spaces-more-morphisms-lemma-deform-property-fp-over-ft}。

\medskip\noindent
在情形 (\ref{stacks-more-morphisms-item-isomorphism-fp-over-ft})、
(\ref{stacks-more-morphisms-item-open-immersion-fp-over-ft})、
(\ref{stacks-more-morphisms-item-monomorphism-fp-over-ft})、
(\ref{stacks-more-morphisms-item-finite-fp-over-ft}) 中，先由引理
\ref{stacks-more-morphisms-lemma-thicken-property-morphisms} 得到 $f$ 与 $f'$
都可由代数空间表示。因此，可以选取代数空间 $U'$ 及满光滑态射
$U' \to \mathcal{B}'$，再选取代数空间 $V'$ 及满光滑态射
$V' \to U' \times_{\mathcal{B}'} \mathcal{Y}'$；于是
$W' = V' \times_{\mathcal{Y}'} \mathcal{X}'$ 是代数空间。
令 $U, V, W$ 分别为 $U', V', W'$ 沿
$\mathcal{B} \to \mathcal{B}'$ 的基变换。
此时同样有 $W = V \times_\mathcal{Y} \mathcal{X}$。
接下来只需证明：$W' \to V'$ 是同构，\ 相应地为开浸入，\ 相应地为单态射，
\ 相应地为有限态射，
当且仅当 $W \to V$ 具有同一性质。参见《叠的性质》中的引理
\ref{stacks-properties-lemma-check-property-covering}。
由此应用上述代数空间的结果便得到结论。

\medskip\noindent
在情形 (\ref{stacks-more-morphisms-item-quasi-finite-fp-over-ft}) 中，首先由
《叠的态射》中的引理
\ref{stacks-morphisms-lemma-finite-type-permanence} 看到
$f$ 与 $f'$ 都是局部有限型的。另一方面，由引理
\ref{stacks-more-morphisms-lemma-thicken-property-morphisms}，$f$ 是拟-DM 的，
当且仅当 $f'$ 是拟-DM 的。要检验 $f$ 或 $f'$ 是否局部拟有限
（《叠的态射》中的定义
\ref{stacks-morphisms-definition-locally-quasi-finite}），
最后还需检验一个关于底层拓扑空间的条件；由证明第一段的讨论，
该条件对 $f$ 成立，当且仅当对 $f'$ 成立。
\end{proof}









\section{仿射概形的提升}
\label{stacks-more-morphisms-section-lifting-affines}

\noindent
考虑实线图表
$$
\xymatrix{
W \ar[d] \ar@{..>}[r] & W' \ar@{..>}[d] \\
\mathcal{X} \ar[r] & \mathcal{X}'
}
$$
其中 $\mathcal{X} \subset \mathcal{X}'$ 是代数叠的增厚，
$W$ 是仿射概形，且 $W \to \mathcal{X}$ 光滑。本节所研究的问题是：
能否找到 $W'$ 及虚线箭头，使得该方形为笛卡儿方形，
且 $W' \to \mathcal{X}'$ 光滑。一般情形下我们不知道答案，
但若 $\mathcal{X} \subset \mathcal{X}'$ 是一阶增厚，
将证明答案是肯定的。

\medskip\noindent
为研究这个问题，引入如下范畴。

\begin{remark}[提升范畴]
\label{stacks-more-morphisms-remark-gerbe-of-lifts}
考虑图表
$$
\xymatrix{
W \ar[d]_x \\
\mathcal{X} \ar[r] & \mathcal{X}'
}
$$
其中 $\mathcal{X} \subset \mathcal{X}'$ 是代数叠的增厚，
$W$ 是代数空间，且 $W \to \mathcal{X}$ 光滑。
如下构造范畴 $\mathcal{C}$ 及函子
$$
p : \mathcal{C} \longrightarrow W_{spaces, \etale}
$$
（记号见《空间的性质》中的定义
\ref{spaces-properties-definition-spaces-etale-site}）。
$\mathcal{C}$ 的一个对象是构成交换图表
\begin{equation}
\label{stacks-more-morphisms-equation-object}
\vcenter{
\xymatrix{
U \ar[d]_a \ar[r]_i & U' \ar[dd]^{x'} \\
W \ar[d]_x & \\
\mathcal{X} \ar[r] & \mathcal{X}'
}
}
\end{equation}
的系统 $(U, U', a, i, x', \alpha)$，其交换性由 $2$-态射
$\alpha : x \circ a \to x' \circ i$ 见证，并且
$U$ 与 $U'$ 是代数空间，$a : U \to W$ 是 \'etale 态射，
$x' : U' \to \mathcal{X}'$ 是光滑态射，且
$U = \mathcal{X} \times_{\mathcal{X}'} U'$。
特别地，$U \subset U'$ 是一个增厚。
一个态射
$$
(U, U', a, i, x', \alpha) \to (V, V', b, j, y', \beta)
$$
由 $(f, f', \gamma)$ 给出，其中 $f : U \to V$ 是 $W$ 上的态射，
$f' : U' \to V'$ 是一个态射，其在 $U$ 上的限制给出 $f$，并且
$\gamma : x' \circ f' \to y'$ 是见证下图右侧三角形交换性的
$2$-态射：
% P12-ERR-0070: see ../control/ERRATA_CANDIDATES.jsonl
% P12-ERR-0071: see ../control/ERRATA_CANDIDATES.jsonl
\begin{equation}
\label{stacks-more-morphisms-equation-morphism}
\vcenter{
\xymatrix{
& V \ar[ld]_f \ar[ldd]^b \ar[rr]_j & & V' \ar[ld]_{f'} \ar[lddd]^{y'} \\
U \ar[d]_a \ar[rr]_i & & U' \ar[dd]_{x'} \\
W \ar[d]_x & \\
\mathcal{X} \ar[rr] & & \mathcal{X}'
}
}
\end{equation}
最后，要求 $\gamma$ 与 $\alpha$、$\beta$ 相容：按照《范畴论》第
\ref{categories-section-formal-cat-cat} 节及第
\ref{categories-section-2-categories} 节中的 $2$-范畴演算，这写成
$$
\beta = (\gamma \star \text{id}_j) \circ (\alpha \star \text{id}_f)
$$
（更简洁地说：$\beta = j^*\gamma \circ f^*\alpha$）。
另一种表述是：对象是带有若干附加性质的交换图表
(\ref{stacks-more-morphisms-equation-object})，态射则是交换图表
(\ref{stacks-more-morphisms-equation-morphism})；这些图表位于《叠的性质》中的注
\ref{stacks-properties-remark-representable-over} 所引入的范畴
$\textit{Spaces}/\mathcal{X}'$ 中。由此清楚可见，$\mathcal{C}$ 是范畴，
而把 $(U, U', a, i, x', \alpha)$ 送到 $a : U \to W$ 的规则
$p : \mathcal{C} \to W_{spaces, \etale}$ 是函子。
\end{remark}

\begin{lemma}
\label{stacks-more-morphisms-lemma-morphisms-lifts-etale}
对任意态射 (\ref{stacks-more-morphisms-equation-morphism})，映射
$f' : V' \to U'$ 都是 \'etale 的。
\end{lemma}

\begin{proof}
事实上，$f : V \to U$ 作为 $W_{spaces, \etale}$ 中的态射是
\'etale 的；又因为 $U' \to \mathcal{X}'$ 与
$V' \to \mathcal{X}'$ 光滑，并且
$U = \mathcal{X} \times_{\mathcal{X}'} U'$、
$V = \mathcal{X} \times_{\mathcal{X}'} V'$，
所以可以应用引理 \ref{stacks-more-morphisms-lemma-deform-property-fp-over-ft}。
\end{proof}

\begin{lemma}
\label{stacks-more-morphisms-lemma-gerbe-of-lifts-fibred}
注 \ref{stacks-more-morphisms-remark-gerbe-of-lifts} 中构造的范畴
$p : \mathcal{C} \to W_{spaces, \etale}$ 是群胚纤维化的。
\end{lemma}

\begin{proof}
我们断言 $p$ 的各纤维范畴都是群胚。
若 (\ref{stacks-more-morphisms-equation-morphism}) 中的 $(f, f', \gamma')$ 是态射，
且 $f : U \to V$ 为同构，则由引理
\ref{stacks-more-morphisms-lemma-deform-property-fp-over-ft}，$f'$ 也是同构，
因而 $(f, f', \gamma')$ 是同构。

\medskip\noindent
考虑 $W_{spaces, \etale}$ 中的态射 $f : V \to U$，
以及 $\mathcal{C}$ 中位于 $U$ 上的对象
$\xi = (U, U', a, i, x', \alpha)$。我们构造位于 $V$ 上的“拉回”
$f^*\xi$。令 $b = a \circ f$。令 $f' : V' \to U'$ 为
其在 $V$ 上的限制是 $f$ 的 \'etale 态射（《空间态射的进一步讨论》中的
引理 \ref{spaces-more-morphisms-lemma-topological-invariance}）。
以 $j : V \to V'$ 记相应的增厚。令 $y' = x' \circ f'$，并令
$\gamma = \text{id} : x' \circ f' \to y'$。置
$$
\beta = \alpha \star \text{id}_f :
x \circ b = x \circ a \circ f \to
x' \circ i \circ f = x' \circ f' \circ j = y' \circ j
$$
显然，$(f, f', \gamma) : (V, V', b, j, y', \beta) \to
(U, U', a, i, x', \alpha)$ 是 (\ref{stacks-more-morphisms-equation-morphism}) 中的态射。
如此构造的态射 $(f, f', \gamma)$ 是强笛卡儿的
（《范畴论》中的定义 \ref{categories-definition-cartesian-over-C}）。
略去详细证明，但其根本原因如下。给定 $\mathcal{C}$ 中的态射
$(g, g', \epsilon) : (Y, Y', c, k, z', \delta) \to
(U, U', a, i, x', \alpha)$，若 $g$ 对某个 $h : Y \to V$ 有
$g = f \circ h$，则由《空间态射的进一步讨论》中的引理
\ref{spaces-more-morphisms-lemma-topological-invariance}，
得到唯一分解 $g' = f' \circ h'$；此后可构造所需的 $\zeta$，使得
$(h, h', \zeta) :  (Y, Y', c, k, z', \delta) \to
(V, V', b, j, y', \beta)$ 是 $\mathcal{C}$ 的态射，并满足
$(g, g', \epsilon) = (f, f', \gamma) \circ (h, h', \zeta)$。

\medskip\noindent
因此，$p : \mathcal{C} \to W_\etale$ 是纤维化范畴
（《范畴论》中的定义 \ref{categories-definition-fibred-category}）。
结合上面已经看到的各纤维范畴都是群胚这一事实，再由《范畴论》中的引理
\ref{categories-lemma-fibred-groupoids}，得到
$p : \mathcal{C} \to W_\etale$ 是群胚纤维化的。
\end{proof}

\begin{lemma}
\label{stacks-more-morphisms-lemma-gerbe-of-lifts-stack}
注 \ref{stacks-more-morphisms-remark-gerbe-of-lifts} 中构造的范畴
$p : \mathcal{C} \to W_{spaces, \etale}$ 是一个群胚叠。
\end{lemma}

\begin{proof}
由引理 \ref{stacks-more-morphisms-lemma-gerbe-of-lifts-fibred}，可见《叠》中的定义
\ref{stacks-definition-stack-in-groupoids} 的第一个条件成立。
按照惯例，我们检验对象的下降，而把态射的下降留给读者。
因此，设在 $W_{spaces, \etale}$ 中给定 $a : U \to W$、
$W_{spaces, \etale}$ 中的覆盖 $\{U_k \to U\}_{k \in K}$、
$\mathcal{C}$ 中位于 $U_k$ 上的对象
$\xi_k = (U_k, U'_k, a_k, i_k, x'_k, \alpha_k)$，以及满足余圈条件的
限制之间的态射
$$
\varphi_{kk'} = (f_{kk'}, f'_{kk'}, \gamma_{kk'}) :
\xi_k|_{U_k \times_U U_{k'}} \to
\xi_{k'}|_{U_k \times_U U_{k'}}
$$
为了证明有效性，可以先细化覆盖。因此，可以假设每个 $U_k$
都是概形（如果愿意，甚至可假设为仿射概形）。写成
% P12-ERR-0072: see ../control/ERRATA_CANDIDATES.jsonl
$$
\xi_k|_{U_k \times_U U_{k'}} =
(U_k \times_U U_{k'}, U'_{kk'}, a_{kk'}, x'_{kk'}, \alpha_{kk'})
$$
于是，把 $\mathcal{C}$ 中态射
$\xi_k|_{U_k \times_U U_{k'}} \to \xi_k$ 的第二个分量取出，
便得到 \'etale 态射
$s_{kk'} : U'_{kk'} \to U'_k$
（由引理 \ref{stacks-more-morphisms-lemma-morphisms-lifts-etale}）。
类似地，考察复合的第二个分量
$$
\xi_k|_{U_k \times_U U_{k'}} \xrightarrow{\varphi_{kk'}}
\xi_{k'}|_{U_k \times_U U_{k'}} \to \xi_{k'}
$$
便得到 \'etale 态射 $t_{kk'} : U'_{kk'} \to U'_{k'}$。
我们断言
$$
j :
\coprod\nolimits_{(k, k') \in K \times K} U'_{kk'}
\xrightarrow{(\coprod s_{kk'}, \coprod t_{kk'})}
(\coprod\nolimits_{k \in K} U'_k) \times (\coprod\nolimits_{k \in K} U'_k)
$$
是一个 \'etale 等价关系。首先，已经看到所示态射的两个分量
$s, t$ 都是 \'etale 的。用
$(\coprod U_k) \times (\coprod U_k) \to
(\coprod U'_k) \times (\coprod U'_k)$
对态射 $j$ 作基变换，得到单态射
$$
\coprod\nolimits_{(k, k') \in K \times K} U_k \times_U U_{k'}
\longrightarrow
(\coprod\nolimits_{k \in K} U_k) \times (\coprod\nolimits_{k \in K} U_k)
$$
因此，由《态射的进一步讨论》中的引理
\ref{more-morphisms-lemma-properties-that-extend-over-thickenings}，
$j$ 是单态射。最后，关系 $j$ 的对称性来自
$\varphi_{kk'}^{-1}$ 是 $\varphi_{k'k}$ 的“翻转”（参见《叠》中的注
\ref{stacks-remarks-definition-descent-datum}），而传递性来自余圈条件
（略去细节）。于是，$\coprod U'_k$ 关于 $j$ 的商是代数空间 $U'$
（《空间》中的定理 \ref{spaces-theorem-presentation}）。
上面已经说明存在增厚 $i : U \to U'$，因为 $j$ 在
$\coprod U_k$ 上的限制给出
$(\coprod U_k) \times_U (\coprod U_k)$。
最后，暂时把 $1$-态射 $x'_k : U'_k \to \mathcal{X}'$
视为叠 $\mathcal{X}'$ 在 $U'_k$ 上的对象，则可见这些对象配备了
相对于 \'etale 覆盖 $\{U'_k \to U'\}$ 的下降数据；该数据由
$\mathcal{C}$ 中态射 $\varphi_{kk'}$ 的第三个分量
$\gamma_{kk'}$ 给出。由于 $\mathcal{X}'$ 是叠，该下降数据有效；
翻译回原语言，得到 $1$-态射 $x' : U' \to \mathcal{X}'$，
使各复合 $U'_k \to U' \to \mathcal{X}'$ 配备到 $x'_k$ 的同构，
并与 $\gamma_{kk'}$ 相容。这意味着态射
$\alpha_k : x \circ a_k \to x'_k \circ i_k$ 粘合成态射
$\alpha : x \circ a \to x' \circ i$。于是
$\xi = (U, U', a, i, x', \alpha)$ 就是所求的位于 $U$ 上的对象。
\end{proof}

\begin{lemma}
\label{stacks-more-morphisms-lemma-etale-local-lifts}
设 $\mathcal{X} \subset \mathcal{X}'$ 是代数叠的一个增厚。
设 $W$ 是代数空间，且 $W \to \mathcal{X}$ 是光滑态射。
存在一个 \'etale 覆盖 $\{W_i \to W\}_{i \in I}$，并且对每个 $i$，
存在笛卡儿图表
$$
\xymatrix{
W_i \ar[r] \ar[d] & W_i' \ar[d] \\
\mathcal{X} \ar[r] & \mathcal{X}'
}
$$
其中 $W_i' \to \mathcal{X}'$ 是光滑的。
\end{lemma}

\begin{proof}
选取一个概形 $U'$ 以及一个满光滑态射 $U' \to \mathcal{X}'$。
像往常一样，置 $U = \mathcal{X} \times_{\mathcal{X}'} U'$。于是
$U \to \mathcal{X}$ 是满光滑态射。因此，基变换
$$
V = W \times_{\mathcal{X}} U \longrightarrow W
$$
是代数空间的满光滑态射。由《空间上的拓扑》中的引理
\ref{spaces-topologies-lemma-etale-dominates-smooth}，
可以找到一个 \'etale 覆盖 $\{W_i \to W\}$，使得
$W_i \to W$ 经由 $V \to W$ 分解。再用仿射概形覆盖 $W_i$
（《空间的性质》中的引理
\ref{spaces-properties-lemma-cover-by-union-affines}），
便可假设每个 $W_i$ 都是仿射的。我们可以并且确实以 $W_i$ 代替
$W$，从而归结到下一段讨论的情形。

\medskip\noindent
假设 $W$ 是仿射的，并且给定态射 $W \to \mathcal{X}$ 经由 $U$
分解。图示为
$$
W \xrightarrow{i} U \to \mathcal{X}
$$
由于 $W$ 和 $U$ 在 $\mathcal{X}$ 上都是光滑的，可见 $i$ 是局部有限型的
（《叠的态射》中的引理
\ref{stacks-morphisms-lemma-finite-type-permanence}）。
以 $\mathbf{A}^n_U$ 代替 $U$ 后，可以假设 $i$ 是浸入；参见
《态射》中的引理 \ref{morphisms-lemma-quasi-affine-finite-type-over-S}。
由《叠的态射》中的引理
\ref{stacks-morphisms-lemma-lci-permanence}，态射 $i$ 是局部完全交。
因此，$i$ 是 Koszul-正则浸入（定义见《除子》中的定义
\ref{divisors-definition-regular-immersion}）；这由《态射的进一步讨论》中的引理
\ref{more-morphisms-lemma-lci} 得出。

\medskip\noindent
我们仍可用仿射开覆盖代替 $W$。对每个点 $w \in W$，可以选取
仿射开集 $U'_w \subset U'$，使得若 $U_w \subset U$ 是对应的仿射开集，
则 $w \in i^{-1}(U_w)$，并且 $i^{-1}(U_w) \to U_w$ 是由
Koszul-正则序列
$f_1, \ldots, f_r \in \Gamma(U_w, \mathcal{O}_{U_w})$
截出的闭浸入。这由 Koszul-正则浸入的定义和《除子》中的引理
\ref{divisors-lemma-regular-ideal-sheaf-scheme} 得出。
置 $W_w = i^{-1}(U_w)$；这是 $w \in W$ 的一个仿射开邻域。
选取 $f_1, \ldots, f_r$ 的提升
$f'_1, \ldots, f'_r \in \Gamma(U'_w, \mathcal{O}_{U'_w})$。
这是可行的，因为 $U_w \to U'_w$ 是仿射概形的闭浸入。
令 $W'_w \subset U'_w$ 为由 $f'_1, \ldots, f'_r$ 截出的闭子概形。
我们断言 $W'_w \to \mathcal{X}'$ 是光滑的。该断言完成证明，因为依构造
$W_w = \mathcal{X} \times_{\mathcal{X}'} W'_w$。

\medskip\noindent
为检验该断言，只须对每个在 $\mathcal{X}'$ 上光滑的仿射概形 $X'$，
检验基变换 $W'_w \times_{\mathcal{X}'} X' \to X'$ 是光滑的。
选取一个 \'etale 态射
$$
Y' \to U'_w \times_{\mathcal{X}'} X'
$$
其中 $Y'$ 是仿射的。由于这些态射的像覆盖
$U'_w \times_{\mathcal{X}'} X'$，只须证明 $Y'$ 中由
$f'_1, \ldots, f'_r$ 截出的闭子概形 $Z'$ 在 $X'$ 上光滑。图示为
$$
\xymatrix{
Z' \ar[r] \ar[d] & Y' \ar[d] \\
W'_w \times_{\mathcal{X}'} X' \ar[d] \ar[r] &
U'_w \times_{\mathcal{X}'} X' \ar[d] \ar[r] & X' \\
W'_w = V(f'_1, \ldots, f'_r) \ar[r] & U'_w
}
$$
置 $X = \mathcal{X} \times_{\mathcal{X}'} X'$、
$Y = X \times_{X'} Y' = \mathcal{X} \times_{\mathcal{X}'} Y'$，以及
$Z = Y \times_{Y'} Z' = X \times_{X'} Z' =
\mathcal{X} \times_{\mathcal{X}'} Z'$。
于是 $(Z \subset Z') \to (Y \subset Y') \subset (X \subset X')$
是仿射概形增厚之间的（笛卡儿）态射，并且已知 $Z \to X$ 和
$Y' \to X'$ 都是光滑的。最后，由于 $Y' \to U'_w$ 光滑从而平坦，
函数序列 $f'_1, \ldots, f'_r$ 由《代数的进一步讨论》中的引理
\ref{more-algebra-lemma-koszul-regular-flat-base-change}
映成 $\Gamma(Y', \mathcal{O}_{Y'})$ 中的 Koszul-正则序列。
由《代数的进一步讨论》中的引理
\ref{more-algebra-lemma-cut-by-koszul}
（并利用 Koszul-正则序列是拟正则序列这一事实；参见《代数的进一步讨论》中的引理
\ref{more-algebra-lemma-regular-koszul-regular}、
\ref{more-algebra-lemma-koszul-regular-H1-regular} 和
\ref{more-algebra-lemma-H1-regular-quasi-regular}），
得出 $Z' \to X'$ 如所需是光滑的。
\end{proof}

\begin{lemma}
\label{stacks-more-morphisms-lemma-etale-local-lifts-isomorphic}
设 $\mathcal{X} \subset \mathcal{X}'$ 是代数叠的一个增厚。考虑交换图表
$$
\xymatrix{
W'' \ar[d]_{x''} & W \ar[l] \ar[r] \ar[d]_x & W' \ar[d]^{x'} \\
\mathcal{X}' & \mathcal{X} \ar[l] \ar[r] & \mathcal{X}'
}
$$
其中两个方形都是笛卡儿的，$W', W, W''$ 是代数空间，且竖直箭头
都是光滑的。那么存在
\begin{enumerate}
\item 一个 \'etale 覆盖 $\{f'_k : W'_k \to W'\}_{k \in K}$，
\item \'etale 态射 $f''_k : W'_k \to W''$，以及
\item $2$-态射 $\gamma_k : x'' \circ f''_k \to x' \circ f'_k$，
\end{enumerate}
使得 (a) $(f'_k)^{-1}(W) = (f''_k)^{-1}(W)$，(b)
$f'_k|_{(f'_k)^{-1}(W)} = f''_k|_{(f''_k)^{-1}(W)}$，并且
(c) 把 $\gamma_k$ 拉回到 (a) 中的闭子概形所得的态射，
与初始图表在 $W$ 上的交换性给出的 $2$-态射一致。
\end{lemma}

\begin{proof}
把给定的增厚记作 $i : W \to W'$ 和 $i'' : W \to W''$。
引理陈述中图表的交换性意味着存在一个 $2$-态射
% P12-ERR-0073: see ../control/ERRATA_CANDIDATES.jsonl
$\delta : x' \circ i' \to x'' \circ i''$。
这就是陈述的 (c) 中所指的 $2$-态射。考虑代数空间
$$
I' = W' \times_{x', \mathcal{X}', x''} W''
$$
以及投影 $p' : I' \to W'$ 和 $q' : I' \to W''$。
注意，存在一个“万有的”$2$-态射
$\gamma : x' \circ p' \to x'' \circ q'$（稍后将用到它）。
$\delta$ 的选择定义了态射
$$
\xymatrix{
W \ar[rr]_\delta & & I' \ar[ld]^{p'} \ar[rd]_{q'} \\
& W' & & W''
}
$$
使得复合 $W \to I' \to W'$ 和 $W \to I' \to W''$ 分别为
$i : W \to W'$ 和
% P12-ERR-0074: see ../control/ERRATA_CANDIDATES.jsonl
$i' : W \to W''$。由于 $x''$ 是光滑的，态射 $p' : I' \to W'$
作为 $x''$ 的基变换是光滑的。

\medskip\noindent
假设可以找到一个 \'etale 覆盖 $\{f'_k : W'_k \to W'\}$ 和态射
$\delta_k : W'_k \to I'$，使得 $\delta_k$ 在
% P12-ERR-0075: see ../control/ERRATA_CANDIDATES.jsonl
$W_k = (f'_k)^{-1}$ 上的限制等于 $\delta \circ f_k$，
其中 $f_k = f'_k|_{W_k}$。图示为
$$
\xymatrix{
W_k \ar[r]^{f_k} \ar[d] & W \ar[r]^\delta & I' \ar[d]^{p'} \\
W'_k \ar[rr]^{f'_k} \ar[rru]^{\delta_k} & & W'
}
$$
换言之，我们希望在可能以一个 \'etale 覆盖代替 $W'$ 后，
把 $p'$ 的给定截面 $\delta : W \to I'$ 延拓为 $W'$ 上的截面。

\medskip\noindent
若此事成立，则可置 $f''_k = q' \circ \delta_k$ 和
$\gamma_k = \gamma \star \text{id}_{\delta_k}$（更简洁地说，
$\gamma_k = \delta_k^*\gamma$）。也就是说，剩下只须证明
% P12-ERR-0076: see ../control/ERRATA_CANDIDATES.jsonl
态射 $f''_k$ 是 \'etale 的。依构造，态射 $x' \circ p'$ 与
$x'' \circ q'$ 是 $2$-同构的。因此，$x'' \circ f''_k$ 与
$x' \circ f'_k$ 是 $2$-同构的。由此，复合
$$
W'_k \xrightarrow{f''_k} W'' \xrightarrow{x''} \mathcal{X}'
$$
是光滑的，因为 $x' \circ f'_k$ 是光滑的。由于 $f_k$ 是 \'etale 的，
由引理 \ref{stacks-more-morphisms-lemma-deform-property-fp-over-ft} 得出 $f''_k$ 是 \'etale 的。

\medskip\noindent
如果该增厚是一阶增厚，那么可以选取任何使 $W_k'$ 为仿射的
% P12-ERR-0077: see ../control/ERRATA_CANDIDATES.jsonl
\'etale 覆盖 $\{W'_k \to W'\}$。事实上，由于 $p'$ 是光滑的，
无穷小提升判据表明 $p'$ 是形式光滑的（《空间态射的进一步讨论》中的引理
\ref{spaces-more-morphisms-lemma-smooth-formally-smooth}）。
由于 $W_k$ 是仿射的，并且 $W_k \to W'_k$ 是一阶增厚
（它是 $\mathcal{X} \to \mathcal{X}'$ 的基变换；参见引理
\ref{stacks-more-morphisms-lemma-base-change-thickening}），便得到所需的 $\delta_k$。

\medskip\noindent
在一般情形下，覆盖和态射 $\delta_k$ 的存在性由《空间态射的进一步讨论》中的引理
\ref{spaces-more-morphisms-lemma-smooth-strong-lift} 得出。
\end{proof}

\begin{lemma}
\label{stacks-more-morphisms-lemma-gerbe-of-lifts}
注 \ref{stacks-more-morphisms-remark-gerbe-of-lifts} 中构造的范畴
$p : \mathcal{C} \to W_{spaces, \etale}$ 是一个胚。
\end{lemma}

\begin{proof}
在引理 \ref{stacks-more-morphisms-lemma-gerbe-of-lifts-stack} 中已经看到，它是一个群胚叠。
因此只须检验《叠》中的定义 \ref{stacks-definition-gerbe}
的条件 (2) 和 (3)。条件 (2) 由引理
\ref{stacks-more-morphisms-lemma-etale-local-lifts} 得出。条件 (3) 由引理
\ref{stacks-more-morphisms-lemma-etale-local-lifts-isomorphic} 得出。
\end{proof}

\begin{lemma}
\label{stacks-more-morphisms-lemma-gerbe-of-lifts-first-order}
在注 \ref{stacks-more-morphisms-remark-gerbe-of-lifts} 中，假设
$\mathcal{X} \subset \mathcal{X}'$ 是一阶增厚。那么
\begin{enumerate}
\item 注 \ref{stacks-more-morphisms-remark-gerbe-of-lifts} 中构造的胚
$p : \mathcal{C} \to W_{spaces, \etale}$ 的对象的自同构层都是阿贝尔的；并且
\item 《叠》中的引理 \ref{stacks-lemma-gerbe-abelian-auts}
所构造的群层 $\mathcal{G}$ 是拟凝聚 $\mathcal{O}_W$-模。
\end{enumerate}
\end{lemma}

\begin{proof}
我们将同时证明两个断言。具体而言，给定对象
$\xi = (U, U', a, i, x', \alpha)$，我们将在
$U_{spaces, \etale}$ 上赋予 $\mathit{Aut}(\xi)$ 一个拟凝聚
$\mathcal{O}_U$-模结构，并证明该结构与拉回相容。
由层的粘合（《格址》第 \ref{sites-section-glueing-sheaves} 节）、
《叠》中的引理 \ref{stacks-lemma-gerbe-abelian-auts} 的证明中
把 $\mathcal{G}$ 构造为自同构层 $\mathit{Aut}(\xi)$ 的粘合，
以及只须在取一个 \'etale 覆盖之后检验模的拟凝聚性这一事实
（《空间的性质》中的引理
\ref{spaces-properties-lemma-characterize-quasi-coherent}），
这便足够了。

\medskip\noindent
我们将采用引理 \ref{stacks-more-morphisms-lemma-etale-local-lifts-isomorphic}
的证明中所用的相同方法来描述层 $\mathit{Aut}(\xi)$。
考虑代数空间
$$
I' = U' \times_{x', \mathcal{X}', x'} U'
$$
及其投影 $p' : I' \to U'$ 和 $q' : I' \to U'$。
在 $I'$ 上有一个万有 $2$-态射
$\gamma : x' \circ p' \to x' \circ q'$。
恒等态射 $x' \to x'$ 定义了对角态射
$$
\xymatrix{
U' \ar[rr]_{\Delta'} & & I' \ar[ld]^{p'} \ar[rd]_{q'} \\
& U' & & U'
}
$$
使复合 $U' \to I' \to U'$ 和 $U' \to I' \to U'$
都是恒等态射。我们用
$U, I, p, q, \Delta$ 表示 $U', I', p', q', \Delta'$ 到
$\mathcal{X}$ 的基变换。
% P12-ERR-0078: see ../control/ERRATA_CANDIDATES.jsonl
由于 $W' \to \mathcal{X}'$ 是光滑的，可见 $p' : I' \to U'$
作为基变换是光滑的。

\medskip\noindent
$\mathit{Aut}(\xi)$ 在 $U$ 上的一个截面是态射
$\delta' : U' \to I'$，它满足 $\delta'|_U = \Delta$ 且
$p' \circ \delta' = \text{id}_{U'}$。明确地说，
$(\text{id}_U, q' \circ \delta', (\delta')^*\gamma) : \xi \to \xi$
是相应自同构的公式。更一般地，若 $f : V \to U$ 是 \'etale 态射，
则存在增厚 $j : V \to V'$ 和 \'etale 态射 $f' : V' \to U'$，
后者在 $V$ 上的限制为 $f$；而 $f^*\xi$ 对应于
$(V, V', a \circ f, j, x' \circ f', f^*\alpha)$，参见引理
\ref{stacks-more-morphisms-lemma-gerbe-of-lifts-fibred} 的证明。
% P12-ERR-0079: see ../control/ERRATA_CANDIDATES.jsonl
$\mathit{Aut}(\xi)$ 在 $V$ 上的一个截面是态射
$\delta' : V' \to I'$，它满足 $\delta'|_V = \Delta \circ f$ 且
$p' \circ \delta' = f'$\footnote{相应自同构的一个公式是
$(\text{id}_V, h', (\delta')^*\gamma)$。这里
$h' : V' \to V'$ 是唯一的（同构）态射，满足
$h'|_V = \text{id}_V$，并使
$$
\xymatrix{
V' \ar[r]_{h'} \ar[rd]_{q' \circ \delta'} & V' \ar[d]^{f'} \\
& U'
}
$$
交换。$h'$ 的唯一性和存在性来自 \'etale 格址的拓扑不变性；
% P12-ERR-0080: see ../control/ERRATA_CANDIDATES.jsonl
参见《空间态射的进一步讨论》中的定理
\ref{spaces-more-morphisms-theorem-topological-invariance}。
读者可能觉得，我们反而应该考察态射
$\delta'' : V' \to V' \times_{\mathcal{X}'} V'$，使得
$\delta'' \circ j = \Delta_{V'/\mathcal{X}'}$ 且
$\text{pr}_1 \circ \delta'' = \text{id}_{V'}$。
这样做也可以：由于 $V' \times_{\mathcal{X}'} V' \to I'$ 是 \'etale 的，
同一拓扑不变性表明，把 $\delta''$ 送到
$\delta' = (V' \times_{\mathcal{X}'} V' \to I') \circ \delta''$，
给出了两组态射之间的双射。}。

\medskip\noindent
由此可知，$\mathit{Aut}(\xi)$ 作为集合层，与《空间态射的进一步讨论》中的注
\ref{spaces-more-morphisms-remark-another-special-case}
为增厚 $(U \subset U')$ 和 $(I \subset I')$ 所定义的层一致；
它们通过 $\text{id}_{U'}$ 和 $p'$ 位于 $(U \subset U')$ 上。
对角态射 $\Delta'$ 是该层的一个截面；用《空间态射的进一步讨论》中的引理
\ref{spaces-more-morphisms-lemma-action-sheaf} 作用于这个截面，得到同构
\begin{equation}
\label{stacks-more-morphisms-equation-isomorphism}
\SheafHom_{\mathcal{O}_U}(\Delta^*\Omega_{I/U}, \mathcal{C}_{U/U'})
\longrightarrow
\mathit{Aut}(\xi)
\end{equation}
在 $U_{spaces, \etale}$ 上成立。
% P12-ERR-0081: see ../control/ERRATA_CANDIDATES.jsonl
还有三件事需要检验：
\begin{enumerate}
\item (\ref{stacks-more-morphisms-equation-isomorphism}) 的构造与 \'etale 局部化交换；
\item $\SheafHom_{\mathcal{O}_U}(\Delta^*\Omega_{I/U}, \mathcal{C}_{U/U'})$
是 $U$ 上的拟凝聚模；
\item $\mathit{Aut}(\xi)$ 中的复合对应于这个拟凝聚模中截面的加法。
\end{enumerate}
我们将依次检验这些事项。

\medskip\noindent
为说明 (1)，必须证明：若 $f : V \to U$ 是 \'etale 的，则在 $U$ 上
用 $\xi$ 构造的 (\ref{stacks-more-morphisms-equation-isomorphism}) 限制为映射
(\ref{stacks-more-morphisms-equation-isomorphism})
$$
\SheafHom_{\mathcal{O}_V}(
\Delta_V^*\Omega_{V \times_\mathcal{X} V/V}, \mathcal{C}_{V/V'}) \to
\mathit{Aut}(\xi|_V)
$$
其中后一个映射是在 $V$ 上用 $\xi|_V$ 构造的，并定义于
$V_{spaces, \etale}$ 上。
这由上面脚注中的讨论和《空间态射的进一步讨论》中的引理
\ref{spaces-more-morphisms-lemma-action-by-derivations-etale-localization}
得出。

\medskip\noindent
(2) 的证明。由于 $p'$ 是光滑的，态射 $I \to U$ 是光滑的，
因而相对微分模 $\Omega_{I/U}$ 是有限局部自由的
（《空间态射的进一步讨论》中的引理
\ref{spaces-more-morphisms-lemma-smooth-omega-finite-locally-free}）。
另一方面，$\mathcal{C}_{U/U'}$ 是拟凝聚的
（《空间态射的进一步讨论》中的定义
\ref{spaces-more-morphisms-definition-conormal-sheaf}）。
由《空间的性质》中的引理
\ref{spaces-properties-lemma-properties-quasi-coherent} 即得结论。

\medskip\noindent
(3) 的证明。存在态射 $c' : I' \times_{p', U', q'} I' \to I'$，
使得 $(U', I', p', q', c')$ 是以 $\Delta'$ 为恒等元的代数空间群胚。
例如参见《代数叠》中的引理
\ref{algebraic-lemma-map-space-into-stack}。
$\mathit{Aut}(\xi)$ 中的复合由态射 $c'$ 如下诱导。假设有两个态射
$$
\delta'_1, \delta'_2 : U' \longrightarrow I'
$$
如上对应于 $\mathit{Aut}(\xi)$ 在 $U$ 上的截面；换言之，
% P12-ERR-0082: see ../control/ERRATA_CANDIDATES.jsonl
有 $\delta'_i|U = \Delta_U$ 且
$p' \circ \delta'_i = \text{id}_{U'}$。那么 $\mathit{Aut}(\xi)$ 中的复合为
$$
\delta'_1 \circ \delta'_2 = c'(\delta'_1 \circ q' \circ \delta'_2, \delta'_2)
$$
我们略去详细检验\footnote{读者可以立即看出，必须用
$q' \circ \delta'_2$ 预复合 $\delta'_1$，才能得到纤维积
$I' \times_{p', U', q'} I'$ 的一个良定义的 $U'$-值点。}。
因此，正处于《空间中群胚的进一步讨论》第
\ref{spaces-more-groupoids-section-groupoid-sections} 节所描述的情形，
而所需结论由《空间中群胚的进一步讨论》中的引理
\ref{spaces-more-groupoids-lemma-composition-is-addition} 得出。
\end{proof}

\begin{proposition}[Emerton]
\label{stacks-more-morphisms-proposition-affine-smooth-lift-to-first-order}
\begin{reference}
Matthew Emerton 于 2016 年 4 月 27 日的电子邮件。
\end{reference}
设 $\mathcal{X} \subset \mathcal{X}'$ 是代数叠的一阶增厚。
设 $W$ 是仿射概形，且 $W \to \mathcal{X}$ 是光滑态射。那么存在笛卡儿图表
$$
\xymatrix{
W \ar[d] \ar[r] & W' \ar[d] \\
\mathcal{X} \ar[r] & \mathcal{X}'
}
$$
其中 $W' \to \mathcal{X}'$ 是光滑的，并且 $W'$ 是仿射的。
\end{proposition}

\begin{proof}
考虑注 \ref{stacks-more-morphisms-remark-gerbe-of-lifts} 中引入的范畴
$p : \mathcal{C} \to W_{spaces, \etale}$。该命题断言存在
$\mathcal{C}$ 中位于 $W$ 上的对象。具体而言，如果有这样的对象
$(W, W', a, i, y', \alpha)$，则
$W = \mathcal{X} \times_{\mathcal{X}'} W'$。
因而 $W \to W'$ 是代数空间的增厚，所以由《空间态射的进一步讨论》中的引理
\ref{spaces-more-morphisms-lemma-thickening-scheme}
和《态射的进一步讨论》中的引理
\ref{more-morphisms-lemma-thickening-affine-scheme}，$W'$ 是仿射的。

\medskip\noindent
引理 \ref{stacks-more-morphisms-lemma-gerbe-of-lifts} 告诉我们，$\mathcal{C}$ 是
$W_{spaces, \etale}$ 上的胚。这意味着可以在 \'etale 局部找到解，
并且这些局部解在 \'etale 局部同构；这一部分不需要增厚是一阶的假设。
由引理 \ref{stacks-more-morphisms-lemma-gerbe-of-lifts-first-order}，该胚的对象的自同构层
是阿贝尔的，并且粘合成 $W_{spaces, \etale}$ 上的拟凝聚模
$\mathcal{G}$。我们将检验《格址上的上同调》中的引理
\ref{sites-cohomology-lemma-existence} 的条件 (1) 和 (2)，
以推出存在 $\mathcal{C}$ 中位于 $W$ 上的对象。
条件 (1) 成立：其中每个 $W_i$ 都是仿射的 \'etale 覆盖
$\{W_i \to W\}$ 在所有覆盖组成的集合中共终。
对这样的覆盖，$W_i$ 和 $W_i \times_W W_j$ 都是仿射的，而
$H^1(W_i, \mathcal{G})$ 和 $H^1(W_i \times_W W_j, \mathcal{G})$
都为零：例如，由《空间的上同调》中的命题
\ref{spaces-cohomology-proposition-vanishing}，
仿射代数空间上拟凝聚模的上同调为零。
最后，条件 (2) 是对于拟凝聚层 $\mathcal{G}$ 有
$H^2(W, \mathcal{G}) = 0$；这同样由《空间的上同调》中的命题
\ref{spaces-cohomology-proposition-vanishing} 得出。证明完毕。
\end{proof}

\section{无穷小形变}
\label{stacks-more-morphisms-section-inf}

\noindent
我们继续《Artin 公理》第 \ref{artin-section-inf} 节中的讨论。

\begin{lemma}
\label{stacks-more-morphisms-lemma-inf-quasi-coherent}
设 $\mathcal{X}$ 是概形 $S$ 上的代数叠。假设
$\mathcal{I}_\mathcal{X} \to \mathcal{X}$ 是局部有限表示的。
设 $A \to B$ 是平坦的 $S$-代数同态。设 $x$ 是
$\mathcal{X}$ 在 $A$ 上的对象，并置 $y = x|_B$。那么
$\text{Inf}_x(M) \otimes_A B = \text{Inf}_y(M \otimes_A B)$。
\end{lemma}

\begin{proof}
回忆，$\text{Inf}_x(M)$ 是 $x$ 到 $A[M]$ 的平凡形变的自同构组成的集合，
这些自同构在 $A$ 上诱导 $x$ 的恒等自同构。该平凡形变是经由
$\Spec(A[M]) \to \Spec(A)$ 把 $x$ 拉回到 $\Spec(A[M])$ 所得的。
令 $G \to \Spec(A)$ 为 $x$ 的自同构群代数空间
% P12-ERR-0083: see ../control/ERRATA_CANDIDATES.jsonl
（它存在是因为 $\mathcal{X}$ 是代数空间）。令
$e : \Spec(A) \to G$ 为单位元。《空间态射的进一步讨论》第
\ref{spaces-more-morphisms-section-action-by-derivations} 节的讨论给出
$$
\text{Inf}_x(M) = \Hom_A(e^*\Omega_{G/A}, M)
$$
同理，
$$
\text{Inf}_y(M \otimes_A B) = \Hom_B(e_B^*\Omega_{G_B/B}, M \otimes_A B)
$$
由于依假设 $G \to \Spec(A)$ 是局部有限表示的，可见
$\Omega_{G/A}$ 是局部有限表示的；参见《空间态射的进一步讨论》中的引理
\ref{spaces-more-morphisms-lemma-finite-presentation-differentials}。
因此，$e^*\Omega_{G/A}$ 是有限表示 $A$-模。此外，由《空间态射的进一步讨论》中的引理
\ref{spaces-more-morphisms-lemma-base-change-differentials}，
$\Omega_{G_B/B}$ 是 $\Omega_{G/A}$ 的拉回。因此
$e_B^*\Omega_{G_B/B} = e^*\Omega_{G/A} \otimes_A B$。
% P12-ERR-0084: see ../control/ERRATA_CANDIDATES.jsonl
由《代数的进一步讨论》中的引理
\ref{more-algebra-lemma-pseudo-coherence-and-base-change-ext} 即得结论。
\end{proof}

\begin{lemma}
\label{stacks-more-morphisms-lemma-sheaf-of-infinitesimal-lifts}
设 $\mathcal{X}$ 是基概形 $S$ 上的代数叠。假设
$\mathcal{I}_\mathcal{X} \to \mathcal{X}$ 是局部有限表示的。
设 $(A' \to A, x)$ 是一个形变情形。那么函子
$$
F : B' \longmapsto
\{\text{提升 }x|_{B' \otimes_{A'} A}\text{ 至 } B'\}/\text{同构}
$$
是《拓扑》中的定义 \ref{topologies-definition-big-small-fppf}
所定义的格址 $(\textit{Aff}/\Spec(A'))_{fppf}$ 上的层。
\end{lemma}

\begin{proof}
设 $\{T'_i \to T'\}_{i = 1, \ldots n}$ 是 $A'$ 上仿射概形的一个
标准 fppf 覆盖。写成 $T' = \Spec(B')$。像往常一样，记
$$
T'_{i_0 \ldots i_p} =
T'_{i_0} \times_{T'} \ldots \times_{T'} T'_{i_p} = \Spec(B'_{i_0 \ldots i_p})
$$
其中该环是适当的张量积。置 $B = B' \otimes_{A'} A$ 以及
$B_{i_0 \ldots i_p} = B'_{i_0 \ldots i_p} \otimes_{A'} A$。
记 $y = x|_B$ 和 $y_{i_0 \ldots i_p} = x|_{B_{i_0 \ldots i_p}}$。
设 $\gamma_i \in F(B'_i)$，并假设 $\gamma_{i_0}$ 与
$\gamma_{i_1}$ 在 $F(B'_{i_0i_1})$ 中的像相同。必须找到唯一的
$\gamma \in F(B')$，它在 $F(B'_i)$ 中映到 $\gamma_i$。

\medskip\noindent
在同构类 $\gamma_i$ 中选取 $\textit{Lift}(y_i, B'_i)$ 的一个具体对象
$y'_i$。选取范畴 $\textit{Lift}(y_{i_0i_1}, B'_{i_0i_1})$ 中的同构
$\varphi_{i_0i_1} : y'_{i_0}|_{B'_{i_0i_1}} \to
y'_{i_1}|_{B'_{i_0i_1}}$。如果这些映射
$\varphi_{i_0i_1}$ 满足余圈条件，那么由于 $\mathcal{X}$ 是
fppf 拓扑中的叠，便得到所需的对象 $\gamma$。余圈条件是复合
$$
y'_{i_0}|_{B'_{i_0i_1i_2}}
\xrightarrow{\varphi_{i_0i_1}|_{B'_{i_0i_1i_2}}}
y'_{i_1}|_{B'_{i_0i_1i_2}}
\xrightarrow{\varphi_{i_1i_2}|_{B'_{i_0i_1i_2}}}
y'_{i_2}|_{B'_{i_0i_1i_2}}
\xrightarrow{\varphi_{i_2i_0}|_{B'_{i_0i_1i_2}}}
y'_{i_0}|_{B'_{i_0i_1i_2}}
$$
是恒等态射。若非如此，则这些映射给出元素
$$
\delta_{i_0i_1i_2} \in
\text{Inf}_{y_{i_0i_1i_2}}(J_{i_0i_1i_2}) =
\text{Inf}_y(J) \otimes_B B_{i_0i_1i_2}
$$
这里 $J = \Ker(B' \to B)$，而
$J_{i_0 \ldots i_p} = \Ker(B'_{i_0 \ldots i_p} \to
B_{i_0 \ldots i_p})$。所示等式由引理
\ref{stacks-more-morphisms-lemma-inf-quasi-coherent} 得出；在应用该引理时取
% P12-ERR-0085: see ../control/ERRATA_CANDIDATES.jsonl
$B' \to B'_{i_0 \ldots i_p}$ 以及 $y$ 和 $y_{i_0 \ldots i_p}$。
映射 $B' \to B'_{i_0 \ldots i_p}$ 的平坦性还保证
$J_{i_0 \ldots i_p} = J \otimes_{B'} B'_{i_0 \ldots i_p}$。
一个略去的计算表明，$\delta_{i_0i_1i_2}$ 给出
{\v C}ech 复形中的一个 $2$-余圈
$$
\prod \text{Inf}_y(J) \otimes_B B_{i_0} \to
\prod \text{Inf}_y(J) \otimes_B B_{i_0i_1} \to
\prod \text{Inf}_y(J) \otimes_B B_{i_0i_1i_2} \to \ldots
$$
由《下降》中的引理 \ref{descent-lemma-standard-covering-Cech-quasi-coherent}，
该复形在正次数上无上同调，并且 $H^0 = \text{Inf}_y(J)$。
由于 $\text{Inf}_{y_{i_0i_1}}(J_{i_0i_1})$ 作用在态射上
（《Artin 公理》中的注 \ref{artin-remark-automorphisms}），
这意味着可以修改 $\varphi_{i_0i_1}$ 的选取，从而归结到
$\delta_{i_0i_1i_2} = 0$ 的情形。

\medskip\noindent
唯一性。还须证明，至多存在一个对所有 $i$ 都限制为 $\gamma_i$ 的
$\gamma$。假设有 $\textit{Lift}(y, B')$ 的对象 $y', z'$，以及
$\textit{Lift}(y_i, B'_i)$ 中的同构
$\psi_i : y'|_{B'_i} \to z'|_{B'_i}$。于是可以考察
$$
\psi_{i_1}^{-1} \circ \psi_{i_0} \in
\text{Inf}_{y_{i_0i_1}}(J_{i_0i_1}) =
\text{Inf}_y(J) \otimes_B B_{i_0i_1}
$$
与上面相同地论证可知，在 $B'$ 上寻找 $y'$ 与 $z'$ 之间同构的障碍，
是上述 {\v C}ech 复形的 $H^1$ 中的一个元素，而该群为零。
\end{proof}

\begin{lemma}
\label{stacks-more-morphisms-lemma-T-quasi-coherent}
设 $\mathcal{X}$ 是概形 $S$ 上的代数叠，并且其结构态射
$\mathcal{X} \to S$ 是局部有限表示的。设 $A \to B$ 是平坦的
$S$-代数同态。设 $x$ 是 $\mathcal{X}$ 在 $A$ 上的对象。那么
% P12-ERR-0086: see ../control/ERRATA_CANDIDATES.jsonl
$T_x(M) \otimes_A B = T_y(M \otimes_A B)$。
\end{lemma}

\begin{proof}
选取一个概形 $U$ 以及一个满光滑态射 $U \to \mathcal{X}$。
首先把该引理归结到 $x$ 可提升至 $U$ 的情形。回忆，$T_x(M)$ 是
$x$ 到 $A[M]$ 的提升的同构类组成的集合。因此，引理
\ref{stacks-more-morphisms-lemma-sheaf-of-infinitesimal-lifts}\footnote{该引理适用：
$\Delta : \mathcal{X} \to \mathcal{X} \times_S \mathcal{X}$
由《叠的态射》中的引理
\ref{stacks-morphisms-lemma-finite-presentation-permanence}
以及 $\mathcal{X} \to S$ 是局部有限表示的假设可知是局部有限表示的。
因此，$\mathcal{I}_\mathcal{X} \to \mathcal{X}$ 作为
$\Delta$ 的基变换是局部有限表示的。} 表明，规则
$$
A_1 \mapsto T_{x|_{A_1}}(M \otimes_A A_1)
$$
是 $\Spec(A)$ 的小 \'etale 格址上的层；这里需要张量积，才能使
$A[M] \to A_1[M \otimes_A A_1]$ 成为平坦环映射。
可以选取一个忠实平坦 \'etale 环映射 $A \to A_1$，使
$x|_{A_1}$ 可提升为态射 $u_1 : \Spec(A_1) \to U$；例如参见
《叠上的层》中的引理
\ref{stacks-sheaves-lemma-surjective-flat-locally-finite-presentation}。
写成 $A_2 = A_1 \otimes_A A_1$，并置
$B_1 = B \otimes_A A_1$ 和 $B_2 = B \otimes_A A_2$。考虑图表
$$
\xymatrix{
0 \ar[r] &
T_y(M \otimes_A B) \ar[r] &
T_{y|_{B_1}}(M \otimes_A B_1) \ar[r] &
T_{y|_{B_2}}(M \otimes_A B_2) \\
0 \ar[r] &
T_x(M) \ar[r] \ar[u] &
T_{x|_{A_1}}(M \otimes_A A_1) \ar[r] \ar[u] &
T_{x|_{A_2}}(M \otimes_A A_2) \ar[u]
}
$$
两行由层条件可知都是正合的。又有
$M \otimes_A B_i = (M \otimes_A A_i) \otimes_{A_i} B_i$。
因此，只要对中间和右边的竖直箭头证明结论，所需结果便成立。
这把问题归结到下一段讨论的情形。

\medskip\noindent
假设 $x$ 是态射 $u : \Spec(A) \to U$ 的像。注意，
$T_u(M) \to T_x(M)$ 是满射，因为 $U \to \mathcal{X}$ 光滑且
可由代数空间表示；参见《可表示性判据》中的引理
\ref{criteria-lemma-representable-by-spaces-formally-smooth}
（解释见其前的讨论）以及《空间态射的进一步讨论》中的引理
\ref{spaces-more-morphisms-lemma-smooth-formally-smooth}。
置 $R = U \times_\mathcal{X} U$。回忆，我们得到代数空间中的群胚
$(U, R, s, t, c, e, i)$，并且 $\mathcal{X} = [U/R]$。
由《Artin 公理》中的引理 \ref{artin-lemma-ses-inf-and-T}，有正合列
$$
T_{e \circ u}(M) \to T_u(M) \oplus T_u(M) \to T_x(M) \to 0
$$
其中右端的零由上面已经证明。对到 $B$ 的基变换也有类似的序列。
因此，只要对 $T_u(M)$ 和 $T_{e \circ u}(M)$ 证明该引理的结论，
便得到所需结果。这把问题归结到下一段讨论的情形。

\medskip\noindent
假设 $\mathcal{X} = X$ 是 $S$ 上局部有限表示的代数空间。那么
$$
T_x(M) = \Hom_A(x^*\Omega_{X/S}, M)
$$
这由《空间态射的进一步讨论》第
\ref{spaces-more-morphisms-section-action-by-derivations} 节的讨论得出。
同理，
$$
T_y(M \otimes_A B) = \Hom_B(y^*\Omega_{X/S}, M \otimes_A B)
$$
由于 $X \to S$ 是局部有限表示的，可见 $\Omega_{X/S}$ 是局部有限表示的；
参见《空间态射的进一步讨论》中的引理
\ref{spaces-more-morphisms-lemma-finite-presentation-differentials}。
因此，$x^*\Omega_{X/S}$ 是有限表示 $A$-模。显然，
$y^*\Omega_{X/S} = x^*\Omega_{X/S} \otimes_A B$。
% P12-ERR-0087: see ../control/ERRATA_CANDIDATES.jsonl
由《代数的进一步讨论》中的引理
\ref{more-algebra-lemma-pseudo-coherence-and-base-change-ext} 即得结论。
\end{proof}

\begin{lemma}
\label{stacks-more-morphisms-lemma-local-lift-enough}
设 $\mathcal{X}$ 是概形 $S$ 上的代数叠，并且其结构态射
$\mathcal{X} \to S$ 是局部有限表示的。设
$(A' \to A, x)$ 是一个形变情形。如果存在一个忠实平坦有限表示
$A'$-代数 $B'$，以及 $\mathcal{X}$ 在 $B'$ 上的对象 $y'$，
它提升 $x|_{B' \otimes_{A'} A}$，那么存在 $A'$ 上提升 $x$ 的对象 $x'$。
\end{lemma}

\begin{proof}
令 $I = \Ker(A' \to A)$。置
$B'_1 = B' \otimes_{A'} B'$ 和
$B'_2 = B' \otimes_{A'} B' \otimes_{A'} B'$。令
$J = IB'$、$J_1 = IB'_1$、$J_2 = IB'_2$，以及
$B = B'/J$、$B_1 = B'_1/J_1$、$B_2 = B'_2/J_2$。
置 $y = x|_B$、$y_1 = x|_{B_1}$、$y_2 = x|_{B_2}$。
令 $F$ 为引理 \ref{stacks-more-morphisms-lemma-sheaf-of-infinitesimal-lifts} 中的 fppf 层
（它适用；参见引理 \ref{stacks-more-morphisms-lemma-T-quasi-coherent} 的证明中的脚注）。
于是有等化子图表
$$
\xymatrix{
F(A') \ar[r] &
F(B') \ar@<1ex>[r] \ar@<-1ex>[r] &
F(B'_1)
}
$$
另一方面，采用《Artin 公理》第 \ref{artin-section-inf} 节的术语，有
$F(B') = \text{Lift}(y, B')$、
$F(B'_1) = \text{Lift}(y_1, B'_1)$ 和
$F(B'_2) = \text{Lift}(y_2, B'_2)$。这些集合非空，并且分别是
$T_y(J)$、$T_{y_1}(J_1)$、$T_{y_2}(J_2)$ 的（典范）主齐性空间；
参见《Artin 公理》中的引理
\ref{artin-lemma-properties-lift-RS-star}。
因此，$y'$ 在 $F(B'_1)$ 中的两个像之差是一个元素
$$
\delta_1 \in T_{y_1}(J_1) = T_x(I) \otimes_A B_1
$$
所示等式由引理 \ref{stacks-more-morphisms-lemma-T-quasi-coherent} 得出；在应用该引理时取
% P12-ERR-0088: see ../control/ERRATA_CANDIDATES.jsonl
$A' \to B'_1$ 以及 $x$ 和 $y_1$。映射 $A' \to B'_1$ 的平坦性
还保证 $J_1 = I \otimes_{A'} B'_1$。对于 $B'$ 和 $B'_2$
也有类似的等式。一个略去的计算表明，$\delta_1$ 给出
{\v C}ech 复形中的一个 $1$-余圈
$$
T_x(I) \otimes_A B \to
T_x(I) \otimes_A B_1 \to
T_x(I) \otimes_A B_2 \to \ldots
$$
由《下降》中的引理 \ref{descent-lemma-standard-covering-Cech-quasi-coherent}，
该复形在正次数上无上同调，并且 $H^0 = T_x(I)$。
因此，可以选取 $T_x(I) \otimes_A B = T_y(J)$ 中边界为
$\delta_1$ 的一个元素。以该元素作用所得结果代替 $y'$，
便得到满足 $\delta_1 = 0$ 的新选择 $y'$。于是 $y'$ 在两个映射
$F(B') \to F(B'_1)$ 下映到同一元素，并且由层条件得到
% P12-ERR-0089: see ../control/ERRATA_CANDIDATES.jsonl
$F(A')$ 的一个元素。
\end{proof}

\section{形式光滑态射}
\label{stacks-more-morphisms-section-formally-smooth}

\noindent
本节引入代数叠的形式光滑态射
$\mathcal{X} \to \mathcal{Y}$ 的概念。这种态射的特征是：
当 $T$ 仿射时，$\mathcal{X}$ 的 $T$-值点可以提升到
$T$ 的无穷小增厚。主要结果是，形式光滑且局部有限表示的态射是光滑的；
参见引理 \ref{stacks-more-morphisms-lemma-smooth-formally-smooth}。事实证明，
该判据通常比雅可比判据更容易使用。

\begin{definition}
\label{stacks-more-morphisms-definition-formally-smooth}
若代数叠的态射 $f : \mathcal{X} \to \mathcal{Y}$ 作为群胚纤维化范畴中的
$1$-态射在对象上形式光滑（见《可表示性判据》第
\ref{criteria-section-formally-smooth} 节的解释），则称它是
{\it 形式光滑的}。
\end{definition}

\noindent
我们把该定义的条件翻译成当前使用的语言（参见《叠的性质》第
\ref{stacks-properties-section-conventions} 节）。设
$f : \mathcal{X} \to \mathcal{Y}$ 是代数叠的态射。考虑一个
$2$-交换的实线图表
\begin{equation}
\label{stacks-more-morphisms-equation-diagram}
\vcenter{
\xymatrix{
T \ar[r]_-x \ar[d]_i & \mathcal{X} \ar[d]^f \\
T' \ar[r]^-y \ar@{..>}[ru] & \mathcal{Y}
}
}
\end{equation}
其中 $i : T \to T'$ 是仿射概形的一阶增厚。令
$$
\gamma : y \circ i \longrightarrow f \circ x
$$
为见证该图表 $2$-交换性的 $2$-态射。（记号同《范畴论》第
\ref{categories-section-formal-cat-cat} 节和第
\ref{categories-section-2-categories} 节。）
给定 (\ref{stacks-more-morphisms-equation-diagram}) 和 $\gamma$，一个
{\it 虚线箭头}是三元组 $(x', \alpha, \beta)$，其中包含态射
$x' : T' \to \mathcal{X}$ 和 $2$-箭头
$\alpha : x' \circ i \to x$、$\beta : y \to f \circ x'$，并满足
$\gamma = (\text{id}_f \star \alpha) \circ
(\beta \star \text{id}_i)$；换言之，图表
$$
\xymatrix{
& f \circ x' \circ i \ar[rd]^{\text{id}_f \star \alpha} \\
y \circ i \ar[ru]^{\beta \star \text{id}_i} \ar[rr]^\gamma & &
f \circ x
}
$$
交换。虚线箭头的一个{\it 态射}
$(x'_1, \alpha_1, \beta_1) \to (x'_2, \alpha_2, \beta_2)$
是 $2$-箭头 $\theta : x'_1 \to x'_2$，它满足
$\alpha_1 = \alpha_2 \circ (\theta \star \text{id}_i)$ 和
$\beta_2 = (\text{id}_f \star \theta) \circ \beta_1$。

\medskip\noindent
刚才描述的虚线箭头范畴是《范畴论》中的定义
\ref{categories-definition-dotted-arrows} 的一个特例。

\begin{lemma}
\label{stacks-more-morphisms-lemma-reformulate-formal-smoothness}
代数叠的态射 $f : \mathcal{X} \to \mathcal{Y}$ 是形式光滑的
（定义 \ref{stacks-more-morphisms-definition-formally-smooth}），当且仅当对于每个图表
(\ref{stacks-more-morphisms-equation-diagram}) 和 $\gamma$，虚线箭头范畴都非空。
\end{lemma}

\begin{proof}
略去不同语言之间的翻译。
\end{proof}

\begin{lemma}
\label{stacks-more-morphisms-lemma-base-change-formal-smoothness}
代数叠的形式光滑态射沿任意代数叠态射的基变换仍然形式光滑。
\end{lemma}

\begin{proof}
这由《范畴论》中的引理
\ref{categories-lemma-cat-dotted-arrows-base-change} 和定义得出。
\end{proof}

\begin{lemma}
\label{stacks-more-morphisms-lemma-composition-formal-smoothness}
代数叠的形式光滑态射的复合仍然形式光滑。
\end{lemma}

\begin{proof}
这由《范畴论》中的引理
\ref{categories-lemma-cat-dotted-arrows-composition} 和定义得出。
\end{proof}

\begin{lemma}
\label{stacks-more-morphisms-lemma-formal-smoothness-representable}
设 $f : \mathcal{X} \to \mathcal{Y}$ 是可由代数空间表示的代数叠态射。
那么下列条件等价：
\begin{enumerate}
\item $f$ 形式光滑；
\item 对每个概形 $T$ 和态射 $T \to \mathcal{Y}$，态射
$\mathcal{X} \times_\mathcal{Y} T \to T$
作为代数空间的态射是形式光滑的。
\end{enumerate}
\end{lemma}

\begin{proof}
这由《范畴论》中的引理
\ref{categories-lemma-cat-dotted-arrows-base-change} 和定义得出。
\end{proof}

\begin{lemma}
\label{stacks-more-morphisms-lemma-lift-to-smooth}
设 $T \to T'$ 是仿射概形的一阶增厚。设 $\mathcal{X}'$ 是
$T'$ 上的代数叠，并且其结构态射
$\mathcal{X}' \to T'$ 是光滑的。设 $x : T \to \mathcal{X}'$
是 $T'$ 上的态射。那么存在 $T'$ 上的
% P12-ERR-0090: see ../control/ERRATA_CANDIDATES.jsonl
态射 $x' : T' \to \mathcal{X}'$，满足 $x'|_T = x$。
\end{lemma}

\begin{proof}
可以应用引理 \ref{stacks-more-morphisms-lemma-local-lift-enough} 的结果。因此，只须构造
一个光滑满态射 $W' \to T'$，其中 $W'$ 仿射，并使
% P12-ERR-0091: see ../control/ERRATA_CANDIDATES.jsonl
$x|_{T \times_{W'} T'}$ 可提升到 $W'$。（我们促请读者使用已经建立的
代数空间的类似结果，自行找出这一事实的证明。）选取一个概形 $U'$
以及一个满光滑态射 $U' \to \mathcal{X}'$。注意，
$U' \to T'$ 是光滑的，并且投影
$T \times_{\mathcal{X}'} U' \to T$ 是满光滑的。
选取一个仿射概形 $W$ 和一个 \'etale 态射
$W \to T \times_{\mathcal{X}'} U'$，使 $W \to T$ 是满射。
于是 $W \to T$ 是仿射概形的光滑态射。以若干主仿射开集的不交并
代替 $W$ 后，可以假设存在仿射概形的光滑态射 $W' \to T'$，使得
$W = T \times_{T'} W'$；参见《代数》中的引理
\ref{algebra-lemma-lift-smooth}。由《空间态射的进一步讨论》中的引理
\ref{spaces-more-morphisms-lemma-smooth-formally-smooth}，
可以找到 $T'$ 上的态射 $W' \to U'$，它提升给定态射
$W \to U'$。证明完毕。
\end{proof}

\noindent
下面的引理是本节的主要结果。结合《叠的极限》中的命题
\ref{stacks-limits-proposition-characterize-locally-finite-presentation}，
它表明，可以用群胚叠的 $1$-态射
$\mathcal{X} \to \mathcal{Y}$ 的“简单”性质来判定代数叠的态射
$f : \mathcal{X} \to \mathcal{Y}$ 是否光滑。

\begin{lemma}[无穷小提升判据]
\label{stacks-more-morphisms-lemma-smooth-formally-smooth}
设 $f : \mathcal{X} \to \mathcal{Y}$ 是代数叠的态射。下列条件等价：
\begin{enumerate}
\item 态射 $f$ 是光滑的。
\item 态射 $f$ 是局部有限表示且形式光滑的。
\end{enumerate}
\end{lemma}

\begin{proof}
假设 $f$ 光滑。由《叠的态射》中的引理
\ref{stacks-morphisms-lemma-smooth-locally-finite-presentation}，
$f$ 是局部有限表示的。因此，给定一个图表
(\ref{stacks-more-morphisms-equation-diagram}) 和
$\gamma : y \circ i \to f \circ x$，只须找到一个虚线箭头
（参见引理 \ref{stacks-more-morphisms-lemma-reformulate-formal-smoothness}）。
构造纤维积，得到
$$
\xymatrix{
T \ar[d] \ar[r] &
T' \times_\mathcal{Y} \mathcal{X} \ar[d] \ar[r] &
\mathcal{X} \ar[d] \\
T' \ar[r] & T' \ar[r] & \mathcal{Y}
}
$$
由此可见，只须在左方形中找到虚线箭头。由于
$T' \times_\mathcal{Y} \mathcal{X} \to T'$ 是光滑的
（《叠的态射》中的引理
\ref{stacks-morphisms-lemma-base-change-smooth}），
引理 \ref{stacks-more-morphisms-lemma-lift-to-smooth} 保证左方形中的虚线箭头存在。

\medskip\noindent
反过来，假设 $f$ 是局部有限表示且形式光滑的。选取一个概形 $U$
以及一个满光滑态射 $U \to \mathcal{X}$。那么
$a : U \to \mathcal{X}$ 和 $b : U \to \mathcal{Y}$
均可由代数空间表示且局部有限表示（使用《叠的态射》中的引理
\ref{stacks-morphisms-lemma-composition-finite-presentation}
以及上面已见的光滑态射是局部有限表示的事实）。
我们将应用《代数叠》中的引理
\ref{algebraic-lemma-representable-transformations-property-implication}
所给出的一般原理，以《空间态射的进一步讨论》中的引理
\ref{spaces-more-morphisms-lemma-smooth-formally-smooth} 的等价性为输入，
同时使用《可表示性判据》中的引理
\ref{criteria-lemma-representable-by-spaces-formally-smooth} 的翻译。
首先把它应用于 $a$，可见 $a$ 在对象上形式光滑。其次，
利用依假设 $f$ 在对象上形式光滑
（参见引理 \ref{stacks-more-morphisms-lemma-reformulate-formal-smoothness}）以及
《可表示性判据》中的引理
\ref{criteria-lemma-composition-formally-smooth}，可见
$b = f \circ a$ 在对象上形式光滑。然后再次应用该原理，得出
$b$ 光滑。由代数叠态射之光滑性的定义，这意味着 $f$ 光滑，
证明完毕。
\end{proof}

\section{吹起与平坦性}
\label{stacks-more-morphisms-section-blowup-flat}

\noindent
本节简要讨论，对于代数空间上的代数叠，可以从
《空间态射的进一步讨论》第
\ref{spaces-more-morphisms-section-blowup-flat} 节和第
\ref{spaces-more-morphisms-section-applications-flattening-by-blowing-up}
节推出哪些结论。

\begin{lemma}
\label{stacks-more-morphisms-lemma-flatten-stack}
设 $f : \mathcal{X} \to Y$ 是从代数叠到代数空间的态射。
设 $V \subset Y$ 是开子空间。假设
\begin{enumerate}
\item $Y$ 是拟紧且拟分离的；
\item $f$ 是有限型且拟分离的；
\item $V$ 是拟紧的；并且
\item $\mathcal{X}_V$ 在 $V$ 上平坦且局部有限表示。
\end{enumerate}
那么存在一个 $V$-可容许吹起 $Y' \to Y$，以及闭子叠
$\mathcal{X}' \subset \mathcal{X}_{Y'}$，满足
$\mathcal{X}'_V = \mathcal{X}_V$，使得
$\mathcal{X}' \to Y'$ 平坦且有限表示。
\end{lemma}

\begin{proof}
注意，$\mathcal{X}$ 是拟紧的。选取一个仿射概形 $U$
以及一个满光滑态射 $U \to \mathcal{X}$。令
$R = U \times_\mathcal{X} U$，于是得到 $Y$ 上代数空间中的群胚
$(U, R, s, t, c)$，并且 $\mathcal{X} = [U/R]$
（《代数叠》中的引理 \ref{algebraic-lemma-stack-presentation}）。
可以把《空间态射的进一步讨论》中的引理
\ref{spaces-more-morphisms-lemma-flat-after-blowing-up}
应用于 $U \to Y$ 和开子空间 $V \subset Y$。于是得到一个
$V$-可容许吹起 $Y' \to Y$，使得严格变换
$U' \subset U_{Y'}$ 在 $Y'$ 上平坦且有限表示。
令 $R' \subset R_{Y'}$ 为 $R$ 的严格变换。由于 $s$ 和 $t$
都是光滑的（特别地也是平坦的），由《空间上的除子》中的引理
\ref{spaces-divisors-lemma-strict-transform-flat}，有笛卡儿图表
$$
\vcenter{
\xymatrix{
R' \ar[r] \ar[d] & R_{Y'} \ar[d]^{s_{Y'}} \\
U' \ar[r] & U_{Y'}
}
}
\quad\text{以及}\quad
\vcenter{
\xymatrix{
R' \ar[r] \ar[d] & R_{Y'} \ar[d]^{t_{Y'}} \\
U' \ar[r] & U_{Y'}
}
}
$$
换言之，$U'$ 是 $U_{Y'}$ 的一个 $R_{Y'}$-不变闭子空间。
因此，由《叠的性质》中的引理
\ref{stacks-properties-lemma-substacks-presentation}，
$U'$ 定义闭子叠 $\mathcal{X}' \subset \mathcal{X}_{Y'}$。
态射 $\mathcal{X}' \to Y'$ 平坦且局部有限表示，因为
$U' \to Y'$ 具有这些性质。另一方面，已经知道
$\mathcal{X}' \to Y'$ 是拟紧且拟分离的
（由对 $f$ 的假设以及闭浸入具有这些性质）。证明完毕。
\end{proof}

\section{代数叠的 Chow 引理}
\label{stacks-more-morphisms-section-chow}

\noindent
本节讨论代数叠的 Chow 引理。

\begin{lemma}
\label{stacks-more-morphisms-lemma-finite-cover-factor}
设 $Y$ 是拟紧拟分离代数空间。设 $V \subset Y$ 是拟紧开子空间。
设 $f : \mathcal{X} \to V$ 是满、平坦且局部有限表示的态射。
那么存在有限满态射 $g : Y' \to Y$，使得
% P12-ERR-0092: see ../control/ERRATA_CANDIDATES.jsonl
$V' = g^{-1}(V) \to Y$ 在 Zariski 局部经由 $f$ 分解。
\end{lemma}

\begin{proof}
先证明 $Y$ 是概形的情形。可以选取一个概形 $U$ 以及一个满光滑态射
$U \to \mathcal{X}$。于是 $\{U \to V\}$ 是概形的 fppf 覆盖。
由《态射的进一步讨论》中的引理
\ref{more-morphisms-lemma-dominate-fppf-finite}，存在有限满态射
$V' \to V$，使 $V' \to V$ 在 Zariski 局部经由 $U$ 分解。
由《态射的进一步讨论》中的引理
\ref{more-morphisms-lemma-extend-finite-surjective-morphisms}，
可以找到有限满态射 $Y' \to Y$，它在 $V$ 上的限制如所需为
$V' \to V$。

\medskip\noindent
若 $Y$ 是代数空间，则先沿一个有限满态射 $Y' \to Y$ 作有限基变换，
其中 $Y'$ 是概形，便可看出该引理成立。参见《空间的极限》中的命题
\ref{spaces-limits-proposition-there-is-a-scheme-finite-over}。
\end{proof}

\begin{lemma}
\label{stacks-more-morphisms-lemma-make-section}
设 $f : \mathcal{X} \to Y$ 是从代数叠到代数空间的态射。
设 $V \subset Y$ 是开子空间。假设
\begin{enumerate}
\item $f$ 是分离且有限型的；
\item $Y$ 是拟紧且拟分离的；
\item $V$ 是拟紧的；并且
\item $\mathcal{X}_V$ 是 $V$ 上的胚。
\end{enumerate}
那么存在交换图表
$$
\xymatrix{
\overline{Z} \ar[rd]_{\overline{g}} &
Z \ar[l]^j \ar[d]_g \ar[r]_h & \mathcal{X} \ar[ld]^f \\
& Y
}
$$
其中 $j$ 是开浸入，$\overline{g}$ 和 $h$ 是固有的，并且
$|V|$ 包含于 $|g|$ 的像中。
\end{lemma}

\begin{proof}
假设有交换图表
$$
\xymatrix{
\mathcal{X}' \ar[d]_{f'} \ar[r] & \mathcal{X} \ar[d]^f \\
Y' \ar[r] & Y
}
$$
以及拟紧开子空间 $V' \subset Y'$，使得 $Y' \to Y$
是代数空间的固有态射，$\mathcal{X}' \to \mathcal{X}$
是代数叠的固有态射，$V' \subset Y'$ 满射到 $V$，并且
$\mathcal{X}'_{V'}$ 是 $V'$ 上的胚。那么只须对二元组
$(f' : \mathcal{X}' \to Y', V')$ 证明该引理。略去一些细节。

\medskip\noindent
证明的总体策略。我们将反复应用上面的注记，把问题归结到
$f$ 的像是开集且 $f$ 在该开集上有截面的情形。每一步都很直接，
但步骤相当多，因而证明稍显繁复。

\medskip\noindent
使用《空间的极限》中的命题
\ref{spaces-limits-proposition-there-is-a-scheme-finite-over}，
把问题归结到 $Y$ 是概形的情形。（取一个有限满态射
$Y' \to Y$，其中 $Y'$ 是概形。置
$\mathcal{X}' = \mathcal{X}_{Y'}$，并应用证明开头的注记。）

\medskip\noindent
使用引理 \ref{stacks-more-morphisms-lemma-flatten-stack}
（并使用《叠的态射》中的引理
\ref{stacks-morphisms-lemma-gerbe-fppf} 说明胚是平坦且局部有限表示的），
把问题归结到 $f$ 平坦且有限表示的情形。

\medskip\noindent
由于 $f$ 平坦且局部有限表示，可见 $|f|$ 的像是开集 $W \subset Y$。
又因为 $\mathcal{X}$ 拟紧（这是因为 $f$ 有限型且 $Y$ 拟紧），
所以 $W$ 拟紧。由引理 \ref{stacks-more-morphisms-lemma-finite-cover-factor}，
可以找到有限满态射 $g : Y' \to Y$，使
$g^{-1}(W) \to Y$ 在 Zariski 局部经由
$\mathcal{X} \to Y$ 分解。以 $Y'$ 代替 $Y$，并以
$\mathcal{X} \times_Y Y'$ 代替 $\mathcal{X}$ 后，
归结到下一段所描述的情形。

\medskip\noindent
假设存在 $n \geq 0$、拟紧开集 $W_i \subset Y$
（$i = 1, \ldots, n$），以及态射 $x_i : W_i \to \mathcal{X}$，使得
(a) $f \circ x_i = \text{id}_{W_i}$，
(b) $W = \bigcup_{i = 1, \ldots, n} W_i$ 包含 $V$，并且
(c) $W$ 是 $|f|$ 的像。
对 $n$ 作归纳。起始情形为 $n = 0$：这蕴含 $V = \emptyset$，
此时可取 $\overline{Z} = \emptyset$。
若 $n > 0$，则对 $i = 1, \ldots, n$，考虑约化闭子概形
$Y_i$，其底拓扑空间为 $Y \setminus W_i$。考虑有限态射
$$
Y' = Y \amalg \coprod\nolimits_{i = 1, \ldots, n} Y_i \longrightarrow Y
$$
以及拟紧开集
$$
V' = (W_1 \cap \ldots \cap W_n \cap V) \amalg
\coprod_{i = 1, \ldots, n} (V \cap Y_i).
$$
由证明开头的注记，只要能对二元组
$$
(\mathcal{X} \to Y, W_1 \cap \ldots \cap W_n \cap V)
\quad\text{以及}\quad
(\mathcal{X} \times_Y Y_i \to Y_i, V \cap Y_i),\quad
i = 1, \ldots, n
$$
证明该引理，所需结果便成立。这里使用集合论等式
$V = (W_1 \cap \ldots \cap W_n \cap V) \cup
\bigcup\nolimits_{i = 1, \ldots n} (V \cap Y_i)$。
归纳假设适用于上述第二类二元组。因此归结到下一段所描述的情形。

\medskip\noindent
假设存在 $n \geq 0$、拟紧开集 $W_i \subset Y$
（$i = 1, \ldots, n$），以及态射 $x_i : W_i \to \mathcal{X}$，使得
(a) $f \circ x_i = \text{id}_{W_i}$，
(b) $W = \bigcup_{i = 1, \ldots, n} W_i$ 包含 $V$，
(c) $W$ 是 $|f|$ 的像，并且
(d) $V \subset W_1 \cap \ldots \cap W_n$。
态射
$$
T_{ij} = \mathit{Isom}_\mathcal{X}(x_i|_{W_i \cap W_j \cap V},
x_j|_{W_i \cap W_j \cap V}) \longrightarrow W_i \cap W_j \cap V
$$
是满、平坦且局部有限表示的（《叠的态射》中的引理
\ref{stacks-morphisms-lemma-gerbe-isom-fppf}）。
对每个拟紧开集 $W_i \cap W_j \cap V$ 和态射
$T_{ij} \to W_i \cap W_j \cap V$ 应用引理
\ref{stacks-more-morphisms-lemma-finite-cover-factor}，得到有限满态射
$Y'_{ij} \to Y$。以所有 $Y'_{ij}$ 在 $Y$ 上的纤维积代替 $Y$ 后，
归结到下一段所描述的情形。

\medskip\noindent
假设存在 $n \geq 0$、拟紧开集 $W_i \subset Y$
（$i = 1, \ldots, n$），以及态射 $x_i : W_i \to \mathcal{X}$，使得
(a) $f \circ x_i = \text{id}_{W_i}$，
(b) $W = \bigcup_{i = 1, \ldots, n} W_i$ 包含 $V$，
(c) $W$ 是 $|f|$ 的像，
(d) $V \subset W_1 \cap \ldots \cap W_n$，并且
(e) $x_i$ 与 $x_j$ 在 $W_i \cap W_j \cap V$ 上 Zariski 局部同构。
任取 $y \in V$。假设能找到拟紧开邻域
$y \in V_y \subset V$，使得该引理对二元组
$(\mathcal{X} \to Y, V_y)$ 成立，设其解为
$\overline{Z}_y, Z_y, \overline{g}_y, g_y, h_y$。
由于 $V$ 拟紧，可以找到有限多个点 $y_1, \ldots, y_m$，使
$V = V_{y_1} \cup \ldots \cup V_{y_m}$。于是置
$$
\overline{Z} = \coprod \overline{Z}_{y_j},\quad
Z = \coprod Z_{y_j},\quad
\overline{g} = \coprod \overline{g}_{y_j},\quad
g = \coprod g_{y_j},\quad
h = \coprod h_{y_j}
$$
便给出该引理的一个解。给定 $y$，由条件 (e)，可以选取拟紧开邻域
$y \in V_y \subset V$，以及同构
$\varphi_i : x_1|_{V_y} \to x_i|_{V_y}$（$i = 2, \ldots, n$）。
置 $\varphi_{ij} = \varphi_j \circ \varphi_i^{-1}$。
这把问题带到下一段所描述的情形。

\medskip\noindent
假设存在 $n \geq 0$、拟紧开集 $W_i \subset Y$
（$i = 1, \ldots, n$），以及态射 $x_i : W_i \to \mathcal{X}$，使得
(a) $f \circ x_i = \text{id}_{W_i}$，
(b) $W = \bigcup_{i = 1, \ldots, n} W_i$ 包含 $V$，
(c) $W$ 是 $|f|$ 的像，
(d) $V \subset W_1 \cap \ldots \cap W_n$，并且
(f) 存在同构 $\varphi_{ij} : x_i|_V \to x_j|_V$，满足
$\varphi_{jk} \circ \varphi_{ij} = \varphi_{ik}$。
态射
$$
I_{ij} = \mathit{Isom}_\mathcal{X}(x_i|_{W_i \cap W_j},
x_j|_{W_i \cap W_j}) \longrightarrow W_i \cap W_j
$$
是固有的，因为 $f$ 分离（《叠的态射》中的引理
\ref{stacks-morphisms-lemma-separated-implies-isom}）。
注意，$\varphi_{ij}$ 定义 $I_{ij} \to W_i \cap W_j$ 在 $V$ 上的截面
$V \to I_{ij}$。由《空间态射的进一步讨论》中的引理
\ref{spaces-more-morphisms-lemma-get-section-after-blowup}，
可以找到 $V$-可容许吹起 $p_{ij} : Y_{ij} \to Y$，使
% P12-ERR-0093: see ../control/ERRATA_CANDIDATES.jsonl
$s_{ij}$ 延拓到 $p_{ij}^{-1}(W_i \cap W_j)$。
以所有 $Y_{ij}$ 在 $Y$ 上的纤维积代替 $Y$ 后，
到达下一段所描述的情形。

\medskip\noindent
假设存在 $n \geq 0$、拟紧开集 $W_i \subset Y$
（$i = 1, \ldots, n$），以及态射 $x_i : W_i \to \mathcal{X}$，使得
(a) $f \circ x_i = \text{id}_{W_i}$，
(b) $W = \bigcup_{i = 1, \ldots, n} W_i$ 包含 $V$，
(c) $W$ 是 $|f|$ 的像，
(d) $V \subset W_1 \cap \ldots \cap W_n$，并且
(g) 存在同构
$\varphi_{ij} : x_i|_{W_i \cap W_j} \to x_j|_{W_i \cap W_j}$，满足
$$
\varphi_{jk}|_V \circ \varphi_{ij}|_V = \varphi_{ik}|_V.
$$
必要时再次以一个 $V$-可容许吹起代替 $Y$，可以假设 $V$ 在 $Y$
中稠密且概形论稠密，从而在 $Y$ 的任意包含 $V$ 的开子空间中也如此。
作此代替后，由《空间的态射》中的引理
\ref{spaces-morphisms-lemma-equality-of-morphisms} 以及
% P12-ERR-0094: see ../control/ERRATA_CANDIDATES.jsonl
$I_{ik} \to W_i \cap W_j$ 固有（因而分离）这一事实，得出
$$
\varphi_{jk}|_{W_i \cap W_j \cap W_k} \circ
\varphi_{ij}|_{W_i \cap W_j \cap W_k} =
\varphi_{ik}|_{W_i \cap W_j \cap W_k}
$$
当然，这意味着 $(x_i, \varphi_{ij})$ 是一个
% P12-ERR-0095: see ../control/ERRATA_CANDIDATES.jsonl
下降数据；由于 $\mathcal{X}$ 是叠，得到态射
$x : W \to \mathcal{X}$，它在 $W_i$ 上与 $x_i$ 一致。
由于 $x$ 是分离态射 $\mathcal{X} \to W$ 的一个截面，
可见 $x$ 是固有的（《叠的态射》中的引理
\ref{stacks-morphisms-lemma-section-immersion}）。
因此，该引理现在对
$\overline{Z} = Y$、$Z = W$、$\overline{g} = \text{id}_Y$、
$g = \text{id}_W$、$h = x$ 成立。
\end{proof}

\begin{theorem}[Chow 引理]
\label{stacks-more-morphisms-theorem-chow-finite-type}
\begin{reference}
这是 Ofer Gabber 的结果，参见
\cite[Theorem 1.1]{olsson_proper}
\end{reference}
设 $f : \mathcal{X} \to Y$ 是从代数叠到代数空间的态射。假设
\begin{enumerate}
\item $Y$ 是拟紧且拟分离的；
\item
% P12-ERR-0096: see ../control/ERRATA_CANDIDATES.jsonl
$f$ 是分离有限型的。
\end{enumerate}
那么存在交换图表
$$
\xymatrix{
\mathcal{X} \ar[rd] & X \ar[l] \ar[d] \ar[r] & \overline{X} \ar[ld] \\
& Y
}
$$
其中 $X \to \mathcal{X}$ 是固有满态射，
$X \to \overline{X}$ 是开浸入，并且
% P12-ERR-0097: see ../control/ERRATA_CANDIDATES.jsonl
$\overline{X} \to Y$ 是代数空间的固有态射。
\end{theorem}

\begin{proof}
粗略的想法是利用 $\mathcal{X}$ 有一个稠密开子叠为胚
（《叠的态射》中的命题
\ref{stacks-morphisms-proposition-open-stratum}），
并诉诸引理 \ref{stacks-more-morphisms-lemma-make-section}。
该方法不能直接奏效的原因是，这个开子叠可能并不拟紧，
从而会遇到技术困难。因此先作一个（标准的）归约，化到 Noether 情形。

\medskip\noindent
首先选取闭浸入
$\mathcal{X} \to \mathcal{X}'$，其中 $\mathcal{X}'$
是 $Y$ 上分离且有限型的代数叠。参见《叠的极限》中的引理
\ref{stacks-limits-lemma-separated-closed-in-finite-presentation}。
显然，只须对 $\mathcal{X}'$ 证明该定理；因此可以假设
$\mathcal{X} \to Y$ 是分离且有限表示的。

\medskip\noindent
假设 $\mathcal{X} \to Y$ 是分离且有限表示的。
由《空间的极限》中的命题
\ref{spaces-limits-proposition-approximate}，可以把
$Y = \lim Y_i$ 写成一个由 Noether 代数空间组成、转移态射仿射的系统的
有向极限。由《叠的极限》中的引理
\ref{stacks-limits-lemma-descend-a-stack-down}，存在某个 $i$
以及从代数叠到 $Y_i$ 的有限表示态射
$\mathcal{X}_i \to Y_i$，使得
$\mathcal{X} = Y \times_{Y_i} \mathcal{X}_i$。
增大 $i$ 后，可以假设 $\mathcal{X}_i \to Y_i$ 分离；
参见《叠的极限》中的引理
\ref{stacks-limits-lemma-eventually-separated}。
于是只须对 $\mathcal{X}_i \to Y_i$ 证明该定理。
这把问题归结到下一段讨论的情形。

\medskip\noindent
假设 $Y$ 是 Noether 的。可以用其约化代替
$\mathcal{X}$（《叠的性质》中的定义
\ref{stacks-properties-definition-reduced-induced-stack}）。
这把问题归结到下一段讨论的情形。

\medskip\noindent
假设 $Y$ 是 Noether 的且 $\mathcal{X}$ 是约化的。
由于 $\mathcal{X} \to Y$ 分离且 $Y$ 拟分离，可见
$\mathcal{X}$ 作为代数叠是拟分离的。因此惰性态射
$\mathcal{I}_\mathcal{X} \to \mathcal{X}$ 是拟紧的。
于是，由《叠的态射》中的命题
\ref{stacks-morphisms-proposition-open-stratum}，
存在一个作为胚的稠密开子叠
$\mathcal{V} \subset \mathcal{X}$。令
$\mathcal{V} \to V$ 为把 $\mathcal{V}$ 表示成代数空间 $V$
上的胚的态射。关于 $\mathcal{V} \to V$ 的构造，参见
《叠的态射》中的引理
\ref{stacks-morphisms-lemma-gerbe-over-iso-classes}。
特别地，该构造表明态射 $\mathcal{V} \to Y$ 分解为
$\mathcal{V} \to V \to Y$。图示为
$$
\xymatrix{
\mathcal{V} \ar[r] \ar[d] & \mathcal{X} \ar[d] \\
V \ar[r] & Y
}
$$
由于态射 $\mathcal{V} \to V$ 是满、平坦且有限表示的
（《叠的态射》中的引理
\ref{stacks-morphisms-lemma-gerbe-fppf}），并且
$\mathcal{V} \to Y$ 是局部有限表示的，可知
$V \to Y$ 是局部有限表示的（《叠的态射》中的引理
\ref{stacks-morphisms-lemma-flat-finite-presentation-permanence}）。
注意，$\mathcal{V} \to V$ 是万有同胚
（《叠的态射》中的引理
\ref{stacks-morphisms-lemma-gerbe-bijection-points}）。
由于 $\mathcal{V}$ 拟紧（参见《叠的态射》中的引理
\ref{stacks-morphisms-lemma-locally-closed-in-noetherian}），
可见 $V$ 拟紧。最后，由于 $\mathcal{V} \to Y$ 分离，
$V \to Y$ 也分离；这是把《叠的态射》中的引理
\ref{stacks-morphisms-lemma-check-separated-on-ui-cover}
应用于 $\mathcal{V} \to V \to Y$ 所得的
（其假设如上所见均满足）。

\medskip\noindent
以上说明，《空间的极限》中的引理
\ref{spaces-limits-lemma-embedding-into-affine-over-qs}
的假设适用于态射 $V \to Y$。因此可以找到稠密开子空间
$V' \subset V$ 以及 $Y$ 上的浸入
$V' \to \mathbf{P}^n_Y$。显然，可以用 $V'$ 代替 $V$，并用
$V'$ 在 $\mathcal{V}$ 中的逆像代替 $\mathcal{V}$
（回忆，如上所见 $|\mathcal{V}| = |V|$）。
于是可以假设有图表
$$
\xymatrix{
\mathcal{V} \ar[rr] \ar[d] & & \mathcal{X} \ar[d] \\
V \ar[r] & \mathbf{P}^n_Y \ar[r] & Y
}
$$
其中箭头 $V \to \mathbf{P}^n_Y$ 是浸入。
令 $\mathcal{X}'$ 为态射
$$
j : \mathcal{V} \longrightarrow \mathbf{P}^n_Y \times_Y \mathcal{X}
$$
的概形论像，并令 $Y'$ 为态射
$V \to \mathbf{P}^n_Y$ 的概形论像。得到交换图表
$$
\xymatrix{
\mathcal{V} \ar[r] \ar[d] &
\mathcal{X}' \ar[r] \ar[d] &
\mathbf{P}^n_Y \times_Y \mathcal{X} \ar[d] \ar[r] &
\mathcal{X} \ar[d] \\
V \ar[r] &
Y' \ar[r] &
\mathbf{P}^n_Y \ar[r] &
Y
}
$$
（参见《叠的态射》中的引理
\ref{stacks-morphisms-lemma-factor-factor}）。
我们断言 $\mathcal{V} = V \times_{Y'} \mathcal{X}'$，并且引理
\ref{stacks-more-morphisms-lemma-make-section} 适用于态射
$\mathcal{X}' \to Y'$ 和开子空间 $V \subset Y'$。
若该断言成立，则得到
$$
\xymatrix{
\overline{X} \ar[rd]_{\overline{g}} &
X \ar[l] \ar[d]_g \ar[r]_h & \mathcal{X}' \ar[ld]^f \\
& Y'
}
$$
其中 $X \to \overline{X}$ 是开浸入，
$\overline{g}$ 和 $h$ 是固有的，并且 $|V|$ 包含于 $|g|$ 的像中。
于是复合 $X \to \mathcal{X}' \to \mathcal{X}$ 是固有的
（因为它是固有态射的复合），且其像包含 $|\mathcal{V}|$；
因此该复合是满射。类似地，$\overline{X} \to Y' \to Y$
作为固有态射的复合是固有的。

\medskip\noindent
最后一步是证明该断言。注意，$\mathcal{X}' \to Y'$
是分离且有限型的，$Y'$ 是拟紧且拟分离的，而 $V$ 是拟紧的
（我们完全略去对所有细节的检验）。其次，注意
$b : \mathcal{X}' \to \mathcal{X}$ 在 $\mathcal{V}$ 上是同构；
这由《叠的态射》中的引理
\ref{stacks-morphisms-lemma-scheme-theoretic-image-of-partial-section}
得出。特别地，$\mathcal{V}$ 被识别为 $\mathcal{X}'$ 的开子叠。
态射 $j$ 是拟紧的（源拟紧且目标拟分离），所以由《叠的态射》中的引理
\ref{stacks-morphisms-lemma-existence-plus-flat-base-change}，
$j$ 的概形论像的形成与平坦基变换交换。特别地，
$V \times_{Y'} \mathcal{X}'$ 是态射
$\mathcal{V} \to V \times_{Y'} \mathcal{X}'$ 的概形论像。
然而，由《叠的态射》中的引理
\ref{stacks-morphisms-lemma-universally-closed-permanence}，
$|\mathcal{V}| \to |V \times_{Y'} \mathcal{X}'|$ 的像是闭的
（使用如上所见 $\mathcal{V} \to V$ 是万有同胚，从而是万有闭的）。
此外，该像是稠密的（把刚才所述与《叠的态射》中的引理
\ref{stacks-morphisms-lemma-topology-scheme-theoretic-image} 结合）。
% P12-ERR-0098: see ../control/ERRATA_CANDIDATES.jsonl
因此 $|\mathcal{V}| = |V \times_{Y'} \mathcal{X}'|$。
于是 $\mathcal{V} \to V \times_{Y'} \mathcal{X}'$ 是同构，
该断言证毕。
\end{proof}

\section{Noether 赋值判据}
\label{stacks-more-morphisms-section-Noetherian-valuative-criterion}

\noindent
本节讨论 Noether 情形下，仅使用离散赋值环来表述的代数叠态射的
（精细）赋值判据。判据有许多不同变体；今后将按需在此增补。

\medskip\noindent
设 $f : \mathcal{X} \to \mathcal{Y}$ 是代数叠（或代数空间或概形）
的态射。所谓{\it 精细}赋值判据，是指给定一个具有某些性质的态射
$\mathcal{U} \to \mathcal{X}$，并且只考察如下形式的实线图表中
虚线箭头的存在性或唯一性：
$$
\xymatrix{
\Spec(K) \ar[d] \ar[r] & \mathcal{U} \ar[r] & \mathcal{X} \ar[d] \\
\Spec(A) \ar[rr] \ar@{..>}[rru] & & \mathcal{Y}
}
$$
下文用这一术语描述
% P12-ERR-0099: see ../control/ERRATA_CANDIDATES.jsonl
迄今得到的结果。

\medskip\noindent
代数叠态射的非 Noether 赋值判据：
\begin{enumerate}
\item 《叠的态射》第
\ref{stacks-morphisms-section-valuative-second} 节
（用于对角态射的分离性）；
\item 《叠的态射》第
\ref{stacks-morphisms-section-valuative-diagonal} 节
（用于分离性）；
\item 《叠的态射》第
\ref{stacks-morphisms-section-valutive-criterion} 节
（用于万有闭性）；
\item 《叠的态射》第
\ref{stacks-morphisms-section-valutive-criterion-properness} 节
（用于固有性）。
\end{enumerate}
对于代数空间，有下列赋值判据：
\begin{enumerate}
\item 《空间的态射》第
\ref{spaces-morphisms-section-valuative-criterion-universally-closed} 节
（用于万有闭性）；
\item 《空间的态射》中的引理
\ref{spaces-morphisms-lemma-refined-valuative-criterion-universally-closed}
（万有闭性的精细判据）；
\item 《空间的态射》第
\ref{spaces-morphisms-section-valuative-separatedness} 节
（用于分离性）；
\item 《空间的态射》第
\ref{spaces-morphisms-section-valuative-criterion-properness} 节
（用于固有性）；
\item 《良好空间》第
\ref{decent-spaces-section-valuative-criterion-universally-closed} 节
（用于良好空间的万有闭性）；
\item 《良好空间》中的引理
\ref{decent-spaces-lemma-re-characterize-universally-closed}
（用于代数空间之间良好态射的万有闭性）；
\item 《空间的上同调》第
\ref{spaces-cohomology-section-Noetherian-valuative-criterion} 节包含
Noether 赋值判据：
\begin{enumerate}
\item 《空间的上同调》中的引理
\ref{spaces-cohomology-lemma-check-separated-dvr}
（使用离散赋值环判定分离性）；
\item 《空间的上同调》中的引理
\ref{spaces-cohomology-lemma-check-proper-dvr}
（使用离散赋值环判定固有性）；
\item 《空间的上同调》中的注
\ref{spaces-cohomology-remark-variant}
（讨论如何归约到完备离散赋值环）。
\end{enumerate}
\item 《空间的极限》第
\ref{spaces-limits-section-Noetherian-valuative-criterion} 节讨论
Noether 赋值判据：
\begin{enumerate}
\item 《空间的极限》中的引理
\ref{spaces-limits-lemma-Noetherian-dvr-valuative-separation}
（使用离散赋值环和泛点判定分离性）；
\item 《空间的极限》中的引理
\ref{spaces-limits-lemma-Noetherian-dvr-valuative-proper}
（使用离散赋值环和泛点判定固有性）；
\item 《空间的极限》中的引理
\ref{spaces-limits-lemma-check-universally-closed-Noetherian}
（使用离散赋值环判定万有闭性）。
\end{enumerate}
\item 《空间的极限》第
\ref{spaces-limits-section-refined-valuative-criteria} 节讨论
精细 Noether 赋值判据：
\begin{enumerate}
\item 《空间的极限》中的引理
\ref{spaces-limits-lemma-refined-valuative-criterion-proper} 和
\ref{spaces-limits-lemma-refined-valuative-criterion-universally-closed}
（使用离散赋值环给出的固有性精细判据）；
\item 《空间的极限》中的引理
\ref{spaces-limits-lemma-refined-valuative-criterion-separated}
（使用离散赋值环给出的分离性精细判据）。
\end{enumerate}
\end{enumerate}
对于概形，有下列赋值判据：
\begin{enumerate}
\item 《概形》第
\ref{schemes-section-valuative-criterion-universal-closedness} 节
（用于万有闭性）；
\item 《概形》第 \ref{schemes-section-valuative-separatedness} 节
（用于分离性）；
\item 《态射》第 \ref{morphisms-section-valuative-criteria} 节
（用于固有性）；
\item 《态射》中的引理
\ref{morphisms-lemma-refined-valuative-criterion-universally-closed}
（万有闭性的精细判据）；
\item 《极限》第 \ref{limits-section-Noetherian-valuative-criterion} 节
讨论 Noether 赋值判据：
\begin{enumerate}
\item 《极限》中的引理
\ref{limits-lemma-Noetherian-dvr-valuative-separation}
（使用离散赋值环和泛点判定分离性）；
\item 《极限》中的引理
\ref{limits-lemma-Noetherian-dvr-valuative-proper}
（使用离散赋值环和泛点判定固有性）；
\item 《极限》中的引理
\ref{limits-lemma-check-universally-closed-Noetherian}
（使用离散赋值环判定万有闭性）。
\end{enumerate}
\item 《极限》第 \ref{limits-section-refined-valuative-criteria} 节
讨论精细 Noether 赋值判据：
\begin{enumerate}
\item 《极限》中的引理
\ref{limits-lemma-refined-valuative-criterion-proper} 和
\ref{limits-lemma-refined-valuative-criterion-universally-closed}
（使用离散赋值环给出的固有性精细判据）；
\item 《极限》中的引理
\ref{limits-lemma-refined-valuative-criterion-separated}
（使用离散赋值环给出的分离性精细判据）。
\end{enumerate}
\item 《极限》第 \ref{limits-section-nagata-valuative} 节讨论
Noether 基底上的赋值判据；在该情形下，可以得到在基底上本质有限型的
离散赋值环。
\end{enumerate}
至此结束对先前结果的列举。

\medskip\noindent
本节的许多结果都可以（而且或许应该）诉诸下面的引理来证明，
尽管我们并非始终这样做。

\begin{lemma}
\label{stacks-more-morphisms-lemma-reach-point-closure-Noetherian}
设 $f : \mathcal{X} \to \mathcal{Y}$ 是代数叠的态射。假设
% P12-ERR-0100: see ../control/ERRATA_CANDIDATES.jsonl
$f$ 有限型且 $\mathcal{Y}$ 局部 Noether。设
$y \in |\mathcal{Y}|$ 是 $|f|$ 的像的闭包中的一点。
那么存在代数叠的交换图表
$$
\xymatrix{
\Spec(K) \ar[r] \ar[d] & \mathcal{X} \ar[d]^f \\
\Spec(A) \ar[r] & \mathcal{Y}
}
$$
其中 $A$ 是离散赋值环，$K$ 是其分式域，并且
$\Spec(A)$ 的闭点映到 $y$。
\end{lemma}

\begin{proof}
选取一个仿射概形 $V$、一点 $v \in V$，以及把 $v$ 映到 $y$ 的
光滑态射 $V \to \mathcal{Y}$。映射
$|V| \to |\mathcal{Y}|$ 是开的；并且由《叠的性质》中的引理
\ref{stacks-properties-lemma-points-cartesian}，
$|\mathcal{X} \times_\mathcal{Y} V| \to |V|$ 的像是
$|f|$ 的像的逆像。因此，点 $v$ 属于
$|\mathcal{X} \times_\mathcal{Y} V| \to |V|$ 的像的闭包。
若对 $\mathcal{X} \times_\mathcal{Y} V \to V$ 和点 $v$
证明该引理，则对 $f$ 和 $y$ 也得到该引理。这样便归结到下一段
所描述的情形。

\medskip\noindent
假设有引理中那样的 $f : \mathcal{X} \to Y$ 和 $y \in |Y|$，
其中 $Y$ 是 Noether 仿射概形。由于 $f$ 拟紧，可知
$\mathcal{X}$ 拟紧。因此可以选取一个仿射概形 $W$
以及一个满光滑态射 $W \to \mathcal{X}$。于是 $|f|$ 的像
与 $|W| \to |Y|$ 的像相同。这样便归结到概形的情形，
即《极限》中的引理
\ref{limits-lemma-reach-point-closure-Noetherian}。
\end{proof}

\begin{lemma}
\label{stacks-more-morphisms-lemma-check-separated-dvr}
设 $f : \mathcal{X} \to \mathcal{Y}$ 是代数叠的态射。假设
\begin{enumerate}
\item $\mathcal{Y}$ 局部 Noether；
\item $f$ 局部有限型且拟分离；
\item 对每个交换图表
$$
\xymatrix{
\Spec(K) \ar[r]_x \ar[d]_j & \mathcal{X} \ar[d]^f \\
\Spec(A) \ar[r]^y \ar@{-->}[ru] & \mathcal{Y}
}
$$
其中 $A$ 是离散赋值环、$K$ 是其分式域，并且对任意
$2$-箭头 $\gamma : y \circ j \to f \circ x$，虚线箭头范畴
（《叠的态射》中的定义
\ref{stacks-morphisms-definition-fill-in-diagram}）
或者为空，或者是恰有一个同构类的集合胚（setoid）。
\end{enumerate}
那么 $f$ 分离。
\end{lemma}

\begin{proof}
为证明 $f$ 分离，必须证明
$\Delta : \mathcal{X} \to \mathcal{X} \times_\mathcal{Y} \mathcal{X}$
固有。已经知道 $\Delta$ 可由代数空间表示且局部有限型
（《叠的态射》中的引理
\ref{stacks-morphisms-lemma-properties-diagonal}），并且拟紧、拟分离
（由 $f$ 拟分离的定义）。选取一个概形 $U$ 以及一个满光滑态射
$U \to \mathcal{X} \times_\mathcal{Y} \mathcal{X}$。置
$$
V = \mathcal{X} \times_{\Delta, \mathcal{X} \times_\mathcal{Y} \mathcal{X}} U
$$
只须证明代数空间的态射 $V \to U$ 固有
（《叠的性质》中的引理
\ref{stacks-properties-lemma-check-property-covering}）。
注意，$U$ 局部 Noether（使用《叠的态射》中的引理
\ref{stacks-morphisms-lemma-locally-finite-type-locally-noetherian}
以及 $U \to \mathcal{Y}$ 局部有限型这一事实），而
$V \to U$ 有限型且拟分离（它是 $\Delta$ 的基变换，并使用上面列出的
$\Delta$ 的性质）。应用《空间的上同调》中的引理
\ref{spaces-cohomology-lemma-check-proper-dvr}，只须证明：
给定交换图表
$$
\xymatrix{
\Spec(K) \ar[r]_v \ar[d]_j &
V \ar[d]^g \ar[r] &
\mathcal{X} \ar[d]^\Delta \\
\Spec(A) \ar[r]^u \ar@{-->}[ru] \ar@{..>}[rru] &
U \ar[r] &
\mathcal{X} \times_\mathcal{Y} \mathcal{X}
}
$$
其中 $A$ 是离散赋值环、$K$ 是其分式域，存在唯一使图表交换的
虚线箭头。由《叠的态射》中的引理
\ref{stacks-morphisms-lemma-cat-dotted-arrows-base-change}，
虚线箭头与点线箭头的范畴等价。假设 (3) 蕴含存在唯一的
点线箭头（在同构意义下）；参见《叠的态射》中的引理
\ref{stacks-morphisms-lemma-helper-diagonal}。因此如所需，
存在唯一的虚线箭头。
\end{proof}

\begin{lemma}
\label{stacks-more-morphisms-lemma-refined-valuative-criterion-proper}
设 $f : \mathcal{X} \to \mathcal{Y}$ 和
$h : \mathcal{U} \to \mathcal{X}$ 是代数叠的态射。
假设 $\mathcal{Y}$ 局部 Noether，$f$ 和 $h$ 都是有限型的，
$f$ 分离，并且 $|h| : |\mathcal{U}| \to |\mathcal{X}|$ 的像
在 $|\mathcal{X}|$ 中稠密。如果给定任意 $2$-交换图表
$$
\xymatrix{
\Spec(K) \ar[r]_-u \ar[d]_j & \mathcal{U} \ar[r]_h & \mathcal{X} \ar[d]^f \\
\Spec(A) \ar[rr]^-y & & \mathcal{Y}
}
$$
其中 $A$ 是分式域为 $K$ 的离散赋值环，并给定
% P12-ERR-0101: see ../control/ERRATA_CANDIDATES.jsonl
$\gamma : y \circ j \to f \circ h \circ u$，存在域扩张 $K'/K$
和支配 $A$ 的赋值环 $A' \subset K'$，使诱导图表
$$
\xymatrix{
\Spec(K') \ar[r]_-{x'} \ar[d]_{j'} & \mathcal{X} \ar[d]^f \\
\Spec(A') \ar[r]^-{y'} \ar@{..>}[ru] & \mathcal{Y}
}
$$
对于诱导的 $2$-箭头
$\gamma' : y' \circ j' \to f \circ x'$ 的虚线箭头范畴非空
（《叠的态射》中的定义
\ref{stacks-morphisms-definition-fill-in-diagram}），那么 $f$ 固有。
\end{lemma}

\begin{proof}
只须证明 $f$ 万有闭。设 $V \to \mathcal{Y}$ 是光滑态射，
其中 $V$ 是仿射概形。由《叠的性质》中的引理
\ref{stacks-properties-lemma-points-cartesian}，
$|\mathcal{U} \times_\mathcal{Y} V| \to
|\mathcal{X} \times_\mathcal{Y} V|$ 的像 $I$
是 $|h|$ 的像的逆像。由于
$|\mathcal{X} \times_\mathcal{Y} V| \to |\mathcal{X}|$ 是开映射
（《叠的态射》中的引理
\ref{stacks-morphisms-lemma-fppf-open}），可知 $I$ 在
$|\mathcal{X} \times_\mathcal{Y} V|$ 中稠密。
此外，由于虚线箭头范畴在基变换下表现良好
（《叠的态射》中的引理
\ref{stacks-morphisms-lemma-cat-dotted-arrows-base-change}），
关于（扩张后）虚线箭头存在性的假设被态射
$\mathcal{U} \times_\mathcal{Y} V \to
\mathcal{X} \times_\mathcal{Y} V \to V$ 继承。
因此，该引理的假设对态射
$\mathcal{U} \times_\mathcal{Y} V \to
\mathcal{X} \times_\mathcal{Y} V \to V$ 成立。于是可以假设
$\mathcal{Y}$ 是仿射概形。

\medskip\noindent
假设 $\mathcal{Y} = Y$ 是仿射概形。（从现在起，不必再操心
$2$-箭头 $\gamma$ 和 $\gamma'$；参见《叠的态射》中的引理
\ref{stacks-morphisms-lemma-cat-dotted-arrows-independent}。）
于是 $\mathcal{U}$ 拟紧。选取一个仿射概形 $U$ 以及一个满光滑态射
$U \to \mathcal{U}$。可以并且确实用 $U$ 代替 $\mathcal{U}$。
因此，可以假设 $\mathcal{U}$ 是仿射概形。

\medskip\noindent
假设 $\mathcal{Y} = Y$ 和 $\mathcal{U} = U$ 都是仿射概形。
由 Chow 引理（定理 \ref{stacks-more-morphisms-theorem-chow-finite-type}），
可以选取一个满固有态射 $X \to \mathcal{X}$，其中 $X$ 是代数空间。
下文将使用 $X \to Y$ 作为分离态射的复合是分离的这一事实。
考虑代数空间 $W = X \times_\mathcal{X} U$。投影态射
$W \to X$ 是有限型的。可以用 $W \to X$ 的概形论像代替 $X$，
从而假设 $|W|$ 在 $|X|$ 中的像在 $|X|$ 中稠密
（这里使用 $|h|$ 的像在 $|\mathcal{X}|$ 中稠密，所以作此代替后，
态射 $X \to \mathcal{X}$ 仍然是满射）。我们断言，对于每个
实线交换图表
$$
\xymatrix{
\Spec(K) \ar[r] \ar[d] & W \ar[r] & X \ar[d] \\
\Spec(A) \ar[rr] \ar@{..>}[rru] & & Y
}
$$
其中 $A$ 是分式域为 $K$ 的离散赋值环，存在使图表交换的点线箭头。
首先，只须在以一个扩张代替 $K$、并以该扩张中支配 $A$ 的一个赋值环
代替 $A$ 后证明点线箭头存在；参见《空间的态射》中的引理
\ref{spaces-morphisms-lemma-push-down-solution}。
由该引理的假设，得到扩张 $K'/K$、支配 $A$ 的赋值环
$A' \subset K'$，以及一个箭头
$\Spec(A') \to \mathcal{X}$；该箭头提升复合
$\Spec(A') \to \Spec(A) \to Y$，并与复合
$\Spec(K') \to \Spec(K) \to W \to X$ 相容。
由于 $X \to \mathcal{X}$ 固有，可以使用固有性的赋值判据
（《叠的态射》中的引理
\ref{stacks-morphisms-lemma-criterion-proper}），
找到扩张 $K''/K'$、支配 $A'$ 的赋值环 $A'' \subset K''$，
以及态射 $\Spec(A'') \to X$；该态射提升复合
$\Spec(A'') \to \Spec(A') \to \mathcal{X}$，并与复合
$\Spec(K'') \to \Spec(K') \to \Spec(K) \to X$ 相容。
于是
% P12-ERR-0102: see ../control/ERRATA_CANDIDATES.jsonl
$K''/K$、$A'' \subset K''$ 和态射 $\Spec(A'') \to X$
是上述问题的一个解，故断言成立。

\medskip\noindent
由代数空间情形的该引理（《空间的极限》中的引理
\ref{spaces-limits-lemma-refined-valuative-criterion-proper}），
该断言蕴含态射 $X \to Y$ 固有。由《叠的态射》中的引理
\ref{stacks-morphisms-lemma-image-proper-is-proper}，
得出 $\mathcal{X} \to Y$ 如所需固有。
\end{proof}

\begin{lemma}
\label{stacks-more-morphisms-lemma-refined-valuative-criterion-separated}
设 $f : \mathcal{X} \to \mathcal{Y}$ 和
$h : \mathcal{U} \to \mathcal{X}$ 是代数叠的态射。假设
$\mathcal{Y}$ 局部 Noether，$f$ 局部有限型且拟分离，
$h$ 有限型，并且
$|h| : |\mathcal{U}| \to |\mathcal{X}|$ 的像在
$|\mathcal{X}|$ 中稠密。如果对任意给定的 $2$-交换图表
$$
\xymatrix{
\Spec(K) \ar[r]_-u \ar[d]_j & \mathcal{U} \ar[r]_h & \mathcal{X} \ar[d]^f \\
\Spec(A) \ar[rr]^-y \ar@{..>}[rru] & & \mathcal{Y}
}
$$
其中 $A$ 是分式域为 $K$ 的离散赋值环，并且
$\gamma : y \circ j \to f \circ h \circ u$，虚线箭头范畴
或者为空，或者是恰有一个同构类的集合胚（setoid），那么 $f$ 分离。
\end{lemma}

\begin{proof}
我们必须证明 $\Delta$ 是固有态射。先假设 $\Delta$ 分离。
那么可以将引理 \ref{stacks-more-morphisms-lemma-refined-valuative-criterion-proper}
应用于态射 $\mathcal{U} \to \mathcal{X}$ 和
$\Delta : \mathcal{X} \to
\mathcal{X} \times_\mathcal{Y} \mathcal{X}$。
注意，由于 $f$ 拟分离，$\Delta$ 拟紧。当然，$\Delta$ 局部有限型
（这对任意对角态射都成立，参见《叠的态射》中的引理
\ref{stacks-morphisms-lemma-properties-diagonal}）。最后，假设给定一个
$2$-交换图表
$$
\xymatrix{
\Spec(K) \ar[r]_-u \ar[d]_j &
\mathcal{U} \ar[r]_h &
\mathcal{X} \ar[d]^\Delta \\
\Spec(A) \ar[rr]^-y \ar@{..>}[rru] & &
\mathcal{X} \times_\mathcal{Y} \mathcal{X}
}
$$
其中 $A$ 是分式域为 $K$ 的离散赋值环，并且
$\gamma : y \circ j \to \Delta \circ h \circ u$。
由《叠的态射》中的引理
\ref{stacks-morphisms-lemma-helper-diagonal} 以及本引理中的假设，
可知存在唯一的点线箭头。这证明引理
\ref{stacks-more-morphisms-lemma-refined-valuative-criterion-proper} 的最后一个假设成立，
从而得到结论。

\medskip\noindent
在一般情形下，只须证明 $\Delta$ 分离，因为这样便回到前一种情形。
事实上，我们断言本引理的假设对
$$
\mathcal{U} \to \mathcal{X}
\quad\text{和}\quad
\Delta :
\mathcal{X} \to
\mathcal{X} \times_\mathcal{Y} \mathcal{X}
$$
成立。确实，由于 $\Delta$ 可由代数空间表示，上一段那种图表的
虚线箭头范畴是一个集合胚（例如参见《叠的态射》中的引理
\ref{stacks-morphisms-lemma-cat-dotted-arrows}）。上一段的论证表明，
这些范畴或者为空，或者恰有一个同构类。因此 $\Delta$ 分离。
\end{proof}

\begin{lemma}
\label{stacks-more-morphisms-lemma-valuative-criterion-universally-closed}
设 $f : \mathcal{X} \to \mathcal{Y}$ 是代数叠的态射。
假设 $\mathcal{Y}$ 局部 Noether，并且 $f$ 有限型。如果对任意给定的
$2$-交换图表
$$
\xymatrix{
\Spec(K) \ar[r]_-x \ar[d]_j & \mathcal{X} \ar[d]^f \\
\Spec(A) \ar[r]^-y & \mathcal{Y}
}
$$
其中 $A$ 是分式域为 $K$ 的离散赋值环，
% P12-ERR-0103: see ../control/ERRATA_CANDIDATES.jsonl
并且 $\gamma : y \circ j \to f \circ x$，存在域扩张 $K'/K$
以及支配 $A$ 的赋值环 $A' \subset K'$，使得诱导图表
$$
\xymatrix{
\Spec(K') \ar[r]_-{x'} \ar[d]_{j'} & \mathcal{X} \ar[d]^f \\
\Spec(A') \ar[r]^-{y'} \ar@{..>}[ru] & \mathcal{Y}
}
$$
连同诱导 $2$-箭头 $\gamma' : y' \circ j' \to f \circ x'$ 的
虚线箭头范畴非空（《叠的态射》中的定义
\ref{stacks-morphisms-definition-fill-in-diagram}），那么 $f$ 万有闭。
\end{lemma}

\begin{proof}
设 $V \to \mathcal{Y}$ 是光滑态射，其中 $V$ 是仿射概形。
虚线箭头范畴在基变换下表现良好（《叠的态射》中的引理
\ref{stacks-morphisms-lemma-cat-dotted-arrows-base-change}）。
因此，关于（扩张后）虚线箭头存在性的假设被态射
$\mathcal{X} \times_\mathcal{Y} V \to V$ 继承。
所以本引理的假设对态射
$\mathcal{X} \times_\mathcal{Y} V \to V$ 成立。
于是可以假设 $\mathcal{Y}$ 是仿射概形。

\medskip\noindent
假设 $\mathcal{Y} = Y$ 是 Noether 仿射概形。（从现在起，
不必再操心 $2$-箭头 $\gamma$ 和 $\gamma'$；参见《叠的态射》中的引理
\ref{stacks-morphisms-lemma-cat-dotted-arrows-independent}。）
为证明 $f$ 万有闭，由《叠的极限》中的引理
\ref{stacks-limits-lemma-test-universally-closed}，只须证明对所有 $n$，
映射 $|\mathcal{X} \times \mathbf{A}^n| \to
|Y \times \mathbf{A}^n|$ 都是闭映射。由于本引理的假设被乘积态射
$\mathcal{X} \times \mathbf{A}^n \to Y \times \mathbf{A}^n$
继承（略去细节），问题约化为证明
$|\mathcal{X}| \to |Y|$ 是闭映射。

\medskip\noindent
假设 $Y$ 是 Noether 仿射概形。设
$T \subset |\mathcal{X}|$ 是闭子集。我们必须证明 $T$ 在
$|Y|$ 中的像是闭的。可以用 $T$ 上的约化诱导闭子空间结构
代替 $\mathcal{X}$；略去验证关于虚线箭头存在性的性质在该代替下
保持不变。因此，问题约化为证明
$|\mathcal{X}| \to |Y|$ 的像是闭的。

\medskip\noindent
设 $y \in |Y|$ 是 $|\mathcal{X}| \to |Y|$ 的像之闭包中的一点。
由引理 \ref{stacks-more-morphisms-lemma-reach-point-closure-Noetherian}，可以选取交换图表
$$
\xymatrix{
\Spec(K) \ar[r] \ar[d] & \mathcal{X} \ar[d]^f \\
\Spec(A) \ar[r] & Y
}
$$
其中 $A$ 是离散赋值环，$K$ 是其分式域，并且
$\Spec(A)$ 的闭点映到 $y$。由本引理的假设立即可知，
$y$ 位于 $|\mathcal{X}| \to |Y|$ 的像中，证明完成。
\end{proof}

\section{模空间}
\label{stacks-more-morphisms-section-moduli-spaces}

\noindent
本节讨论从代数叠到代数空间的态射
$f : \mathcal{X} \to Y$。在适当假设下，称 $Y$ 是
$\mathcal{X}$ 的一个{\it 模空间}。如果
$\mathcal{X} = [U/R]$ 是一个表现，那么可得到一个
$R$-不变态射 $U \to Y$，并且（在适当假设下）$Y$ 是群胚
$(U, R, s, t, c)$ 的一个{\it 商}。关于不同类型商的讨论，
可从《群胚的商》第
\ref{groupoids-quotients-section-introduction} 节开始查阅。

\begin{definition}
\label{stacks-more-morphisms-definition-categorical-quotient}
设 $\mathcal{X}$ 是代数叠。设
$f : \mathcal{X} \to Y$ 是到代数空间 $Y$ 的态射。
\begin{enumerate}
\item 如果任意到代数空间 $W$ 的态射
$\mathcal{X} \to W$ 都唯一地经由 $f$ 分解，则称 $f$ 是一个
{\it 范畴模空间}。
\item 如果对代数空间的任意平坦态射 $Y' \to Y$，基变换
$f' : Y' \times_Y \mathcal{X} \to Y'$ 都是范畴模空间，则称 $f$
是一个{\it 一致范畴模空间}。
\end{enumerate}
设 $\mathcal{C}$ 是代数空间范畴的一个全子范畴。
\begin{enumerate}
\item[(3)] 如果 $Y \in \Ob(\mathcal{C})$，并且对任意
$W \in \Ob(\mathcal{C})$，每个态射 $\mathcal{X} \to W$
都唯一地经由 $f$ 分解，则称 $f$ 是一个
{\it $\mathcal{C}$ 中的范畴模空间}。
% P12-ERR-0104: see ../control/ERRATA_CANDIDATES.jsonl
\item[(4)] 如果 $Y \in \Ob(\mathcal{C})$，并且对
$\mathcal{C}$ 中的每个平坦态射 $Y' \to Y$，基变换
$f' : Y' \times_Y \mathcal{X} \to Y'$ 都是
$\mathcal{C}$ 中的范畴模空间，则称其为一个
{\it $\mathcal{C}$ 中的一致范畴模空间}。
\end{enumerate}
\end{definition}

\noindent
由 Yoneda 引理，范畴模空间若存在便是唯一的。下面将它与此前为商
引入的语言相对应。

\begin{lemma}
\label{stacks-more-morphisms-lemma-quotient-compare}
设 $(U, R, s, t, c)$ 是代数空间中的群胚，其中
$s, t : R \to U$ 平坦且局部有限表现。考虑代数叠
$\mathcal{X} = [U/R]$。给定代数空间 $Y$，态射
$f : \mathcal{X} \to Y$ 与 $R$-不变态射
$\phi : U \to Y$ 之间存在一个 $1$-到-$1$ 对应。
\end{lemma}

\begin{proof}
《可表示性判据》中的定理
\ref{criteria-theorem-flat-groupoid-gives-algebraic-stack}
告诉我们 $\mathcal{X}$ 是代数叠。给定态射
$f : \mathcal{X} \to Y$，令 $\phi : U \to Y$ 为复合
$U \to \mathcal{X} \to Y$。由于
$R = U \times_\mathcal{X} U$（《空间中的群胚》中的引理
\ref{spaces-groupoids-lemma-quotient-stack-2-cartesian}），
立即可知 $\phi$ 是 $R$-不变的。反之，如果
$\phi : U \to Y$ 是到代数空间的 $R$-不变态射，那么由
《空间中的群胚》中的引理
\ref{spaces-groupoids-lemma-quotient-stack-2-coequalizer}，
可得到态射 $f : \mathcal{X} \to Y$。也可以利用
《空间中的群胚》中的引理
\ref{spaces-groupoids-lemma-quotient-stack-objects}
对商叠的显式描述，从 $\phi$ 构造 $f$。
\end{proof}

\begin{lemma}
\label{stacks-more-morphisms-lemma-categorical-quotient-compare}
% P12-ERR-0105: see ../control/ERRATA_CANDIDATES.jsonl
采用引理 \ref{stacks-more-morphisms-lemma-quotient-compare} 中的假设和记号。
那么 $f$ 是（一致）范畴模空间，当且仅当 $\phi$ 是（一致）范畴商。
对于全子范畴中的模空间也有类似结论。
\end{lemma}

\begin{proof}
由引理 \ref{stacks-more-morphisms-lemma-quotient-compare} 中建立的 $1$-到-$1$ 对应，立即可知
$f$ 是范畴模空间，当且仅当 $\phi$ 是范畴商
（《群胚的商》中的定义
\ref{groupoids-quotients-definition-categorical}）。
如果 $Y' \to Y$ 是态射，那么
$U' = Y' \times_Y U \to Y' \times_Y \mathcal{X} = \mathcal{X}'$
作为 $U \to \mathcal{X}$ 的基变换，是满、平坦且局部有限表现的态射
（《可表示性判据》中的引理
\ref{criteria-lemma-flat-quotient-flat-presentation}）。
又由纤维积的传递性，$R' = Y' \times_Y R$ 等于
$U' \times_{\mathcal{X}'} U'$。因此
$\mathcal{X}' = [U'/R']$；参见《代数叠》中的备注
\ref{algebraic-remark-flat-fp-presentation}。所以，把我们的情形基变换到
$Y'$ 后，仍得到本引理陈述中的情形。由此立即可知，$f$ 是一致范畴
模空间，当且仅当 $\phi$ 是一致范畴商。
\end{proof}

\begin{lemma}
\label{stacks-more-morphisms-lemma-check-uniform-categorical-quotient-on-affines}
设 $f : \mathcal{X} \to Y$ 是从代数叠到代数空间的态射。
如果对每个仿射概形 $Y'$ 和平坦态射 $Y' \to Y$，基变换
$f' : Y' \times_Y \mathcal{X} \to Y'$ 都是范畴模空间，
那么 $f$ 是一致范畴模空间。
\end{lemma}

\begin{proof}
选取一个 \'etale 覆盖 $\{Y_i \to Y\}$，其中 $Y_i$ 是仿射概形。
% P12-ERR-0106: see ../control/ERRATA_CANDIDATES.jsonl
对每个 $i$ 和 $j$，选取一个仿射开覆盖
$Y_i \times_Y Y_j = \bigcup Y_{ijk}$。置
$\mathcal{X}_i = Y_i \times_Y \mathcal{X}$ 和
$\mathcal{X}_{ijk} = Y_{ijk} \times_Y \mathcal{X}$。
设 $g : \mathcal{X} \to W$ 是到代数空间的态射。考虑图表
$$
\xymatrix{
\mathcal{X}_i \ar[r] \ar[d] & \mathcal{X} \ar[d] \ar[r]_g & W \\
Y_i \ar[r] \ar@{..>}[rru] & Y
}
$$
$\mathcal{X}_i \to Y_i$ 是范畴模空间这一假设给出唯一的
点线箭头 $h_i : Y_i \to W$。$\mathcal{X}_{ijk} \to Y_{ijk}$
是范畴模空间这一假设说明
% P12-ERR-0107: see ../control/ERRATA_CANDIDATES.jsonl
$h_i$ 和 $h_j$ 在 $Y_{ijk}$ 上的限制相等。因此，$h_i$ 和 $h_j$
在 $Y_i \times_Y Y_j$ 上相同。由于
$Y = \coprod Y_i / \coprod Y_i \times_Y Y_j$
（由《空间》第 \ref{spaces-section-presentations} 节），可知存在唯一的
态射 $Y \to W$，使得 $g$ 经由它分解。因此 $f$ 是范畴模空间。
相同论证可在平坦基变换后应用，故 $f$ 是一致范畴模空间。
\end{proof}

\section{Keel--Mori 定理}
\label{stacks-more-morphisms-section-Keel-Mori}

\noindent
本节开始在代数叠的框架下讨论 Keel 和 Mori 的定理。
关于相关文献的讨论，请参见《文献指南》中的第
\ref{guide-subsection-coarse-moduli-spaces} 小节。

\begin{definition}
\label{stacks-more-morphisms-definition-well-nigh-affine}
设 $\mathcal{X}$ 是代数叠。如果存在仿射概形 $U$ 和一个满、平坦、
有限且有限表现的态射 $U \to \mathcal{X}$，则称
$\mathcal{X}$ 是{\it 近乎仿射的}。
\end{definition}

\noindent
我们给这个性质取了一个多少有些荒唐的名字，因为并不打算过多使用它。

\begin{lemma}
\label{stacks-more-morphisms-lemma-well-nigh-affine}
设 $\mathcal{X}$ 是代数叠。下列条件等价：
\begin{enumerate}
\item $\mathcal{X}$ 近乎仿射；
\item 存在群胚概形 $(U, R, s, t, c)$，其中 $U$ 和 $R$ 仿射，
$s, t : R \to U$ 有限局部自由，并且 $\mathcal{X} = [U/R]$。
\end{enumerate}
若这些条件成立，则 $\mathcal{X}$ 拟紧、拟-DM 且分离。
\end{lemma}

\begin{proof}
假设 $\mathcal{X}$ 近乎仿射。选取仿射概形 $U$ 和一个满、平坦、
有限且有限表现的态射 $U \to \mathcal{X}$。置
$R = U \times_\mathcal{X} U$。于是得到代数空间中的群胚
$(U, R, s, t, c)$ 和同构 $[U/R] \to \mathcal{X}$；参见
《代数叠》中的引理 \ref{algebraic-lemma-map-space-into-stack}
及备注 \ref{algebraic-remark-flat-fp-presentation}。由于
$s, t : R \to U$ 是平坦、有限且有限表现的态射
% P12-ERR-0108: see ../control/ERRATA_CANDIDATES.jsonl
（作为 $U \to \mathcal{X})$ 的基变换），可知
$s, t$ 有限局部自由（《态射》中的引理
\ref{morphisms-lemma-finite-flat}）。这蕴含 $R$ 仿射
（因为有限态射是仿射的），故 (2) 成立。

\medskip\noindent
% P12-ERR-0109: see ../control/ERRATA_CANDIDATES.jsonl
假设有一个群胚概形 $(U, R, s, t, c)$，其中 $U$ 和 $R$ 仿射，
并且 $s, t : R \to U$ 有限局部自由。置
$\mathcal{X} = [U/R]$。由《可表示性判据》中的定理
\ref{criteria-theorem-flat-groupoid-gives-algebraic-stack}，
$\mathcal{X}$ 是代数叠（严格来说，这里并不需要这一点，但无论怎样
强调其正确性都不为过）。由《可表示性判据》中的引理
\ref{criteria-lemma-flat-quotient-flat-presentation}，态射
$U \to \mathcal{X}$ 满、平坦且局部有限表现。因此，根据《叠的性质》
中的引理 \ref{stacks-properties-lemma-check-property-covering}，
可以通过检验投影 $U \times_\mathcal{X} U \to U$ 是否有限，来检验
$U \to \mathcal{X}$ 是否有限。由《空间中的群胚》中的引理
\ref{spaces-groupoids-lemma-quotient-stack-2-cartesian}，
$U \times_\mathcal{X} U = R$，故确实如此。因此
$\mathcal{X}$ 近乎仿射。

\medskip\noindent
最后证明末一断言。由《叠的性质》中的引理
\ref{stacks-properties-lemma-quasi-compact-stack}，
$\mathcal{X}$ 拟紧。由《叠的态射》中的引理
\ref{stacks-morphisms-lemma-properties-diagonal-from-presentation}，
$\mathcal{X} = [U/R]$ 拟-DM 且分离。
\end{proof}

\begin{lemma}
\label{stacks-more-morphisms-lemma-affine-over-well-nigh-affine}
设代数叠 $\mathcal{X}$ 近乎仿射。
\begin{enumerate}
\item 如果 $\mathcal{X}$ 是代数空间，那么它是仿射的。
\item 如果 $\mathcal{X}' \to \mathcal{X}$ 是代数叠的仿射态射，
那么 $\mathcal{X}'$ 近乎仿射。
\end{enumerate}
\end{lemma}

\begin{proof}
% P12-ERR-0110: see ../control/ERRATA_CANDIDATES.jsonl
(1) 立即由《空间的极限》中的引理
\ref{spaces-limits-lemma-affine} 得出。不过，这未免大材小用，因为
把引理 \ref{stacks-more-morphisms-lemma-well-nigh-affine} 与《群胚》中的命题
\ref{groupoids-proposition-finite-flat-equivalence} 结合也可得到 (1)。

\medskip\noindent
为证明 (2)，选取仿射概形 $U$ 和一个满、平坦、有限且有限表现的态射
$U \to \mathcal{X}$。那么
$U' = \mathcal{X}' \times_\mathcal{X} U$ 到 $U$ 有仿射态射
（《叠的态射》中的引理
\ref{stacks-morphisms-lemma-base-change-affine}）。因此 $U'$ 是仿射概形。
当然，$U' \to \mathcal{X}'$ 作为 $U \to \mathcal{X}$ 的基变换，
是满、平坦、有限且有限表现的。
\end{proof}

\begin{lemma}
\label{stacks-more-morphisms-lemma-well-nigh-affine-moduli-space}
设代数叠 $\mathcal{X}$ 近乎仿射。在仿射概形范畴中存在一致范畴模空间
$$
f : \mathcal{X} \longrightarrow M
$$
而且 $f$ 分离、拟紧，并且是万有同胚。
\end{lemma}

\begin{proof}
写作 $\mathcal{X} = [U/R]$，其中 $(U, R, s, t, c)$
如引理 \ref{stacks-more-morphisms-lemma-well-nigh-affine} 所述。设 $C$ 是 $U$ 上
$R$-不变函数组成的环；参见《群胚》第
\ref{groupoids-section-finite-flat} 节。置 $M = \Spec(C)$。
由引理 \ref{stacks-more-morphisms-lemma-quotient-compare}，$R$-不变态射 $U \to M$
对应于态射 $f : \mathcal{X} \to M$。《概形》中的引理
\ref{schemes-lemma-morphism-into-affine} 给出的到仿射概形态射的刻画
立即保证
% P12-ERR-0111: see ../control/ERRATA_CANDIDATES.jsonl
$\phi : U \to M$ 是仿射概形范畴中的范畴商。因此 $f$ 是仿射概形
范畴中的范畴模空间（引理
\ref{stacks-more-morphisms-lemma-categorical-quotient-compare}）。

\medskip\noindent
由引理 \ref{stacks-more-morphisms-lemma-well-nigh-affine}，$\mathcal{X}$ 分离，故由
《叠的态射》中的引理
\ref{stacks-morphisms-lemma-compose-after-separated}，$f$ 分离。

\medskip\noindent
由于 $U \to \mathcal{X}$ 满，并且 $U \to M$ 拟紧，由《叠的态射》
中的引理 \ref{stacks-morphisms-lemma-surjection-from-quasi-compact}，
可知 $f$ 拟紧。

\medskip\noindent
由《群胚》中的引理 \ref{groupoids-lemma-integral-over-invariants}，复合
$$
U \to \mathcal{X} \to M
$$
是仿射概形的整态射。特别地，它万有闭
（《态射》中的引理
\ref{morphisms-lemma-integral-universally-closed}）。由于
$U \to \mathcal{X}$ 满，可知 $\mathcal{X} \to M$ 万有闭
（《叠的态射》中的引理
\ref{stacks-morphisms-lemma-image-proper-is-proper}）。为得出
$\mathcal{X} \to M$ 是万有同胚，只须证明它万有双射，即它满且万有单射。

\medskip\noindent
由《叠的态射》中的引理
\ref{stacks-morphisms-lemma-points-presentation}，有
$|\mathcal{X}| = |U|/|R|$。因此，由《群胚》中的引理
\ref{groupoids-lemma-points}，$|f|$ 满，甚至是双射。

\medskip\noindent
设 $C \to C'$ 是环映射。设 $(U', R', s', t', c')$ 是
$(U, R, s, t, c)$ 沿 $M' = \Spec(C') \to M$ 的基变换。
置 $\mathcal{X}' = [U'/R']$。由《群胚的商》中的引理
\ref{groupoids-quotients-lemma-base-change-quotient-stack}，有
$M' \times_M \mathcal{X} = \mathcal{X}'$。设 $C^1$ 是 $U'$ 上
$R'$-不变函数组成的环。置 $M^1 = \Spec(C^1)$，并考虑图表
$$
\xymatrix{
\mathcal{X}' \ar[d]^{f'} \ar[r] & \mathcal{X} \ar[dd]^f \\
M^1 \ar[d] \\
M' \ar[r] & M
}
$$
由《群胚》中的引理 \ref{groupoids-lemma-invariants-base-change}
和《代数》中的引理 \ref{algebra-lemma-universally-bijective}，态射
$M^1 \to M'$ 是同胚。另一方面，把上一段应用于
$(U', R', s', t', c')$，可知 $|f'|$ 是双射。我们得出，沿任意仿射
概形作基变换后，$f$ 都诱导点集之间的双射。因此，由《叠的态射》中的
引理 \ref{stacks-morphisms-lemma-universally-injective-local}，
$f$ 万有单射。

\medskip\noindent
最后，还须证明 $f$ 是仿射概形范畴中的
% P12-ERR-0112: see ../control/ERRATA_CANDIDATES.jsonl
一致模空间。这由以上讨论以及下述事实得出：如果环映射
$C \to C'$ 平坦，那么由《群胚》中的引理
\ref{groupoids-lemma-invariants-base-change}，$C' \to C^1$ 是同构。
\end{proof}

\begin{lemma}
\label{stacks-more-morphisms-lemma-well-nigh-affine-moduli-space-etale}
设 $h : \mathcal{X}' \to \mathcal{X}$ 是代数叠的态射。假设
$\mathcal{X}'$ 和 $\mathcal{X}$ 近乎仿射，$h$ 是 \'etale 的，
并且 $h$ 在自同构群上诱导同构（《叠的态射》中的备注
\ref{stacks-morphisms-remark-identify-automorphism-groups}）。
那么存在笛卡儿图表
$$
\xymatrix{
\mathcal{X}' \ar[d] \ar[r] & \mathcal{X} \ar[d] \\
M' \ar[r] & M
}
$$
其中 $M' \to M$ 是 \'etale 的，竖直箭头是引理
\ref{stacks-more-morphisms-lemma-well-nigh-affine-moduli-space} 中构造的模空间。
\end{lemma}

\begin{proof}
由《叠的态射》中的诸引理
\ref{stacks-morphisms-lemma-aut-iso-unramified} 和
\ref{stacks-morphisms-lemma-stabilizer-preserving}，注意到 $h$ 可由
代数空间表示。选取仿射概形 $U$ 和一个满、平坦、有限且有限表现的态射
$U \to \mathcal{X}$。于是
$U' = \mathcal{X}' \times_\mathcal{X} U$ 是代数空间，并且有一个有限
（特别地，仿射）态射 $U' \to \mathcal{X}'$。由引理
\ref{stacks-more-morphisms-lemma-affine-over-well-nigh-affine}，可知 $U'$ 仿射。
置 $R = U \times_\mathcal{X} U$ 和
$R' = U' \times_{\mathcal{X}'} U'$，便得到群胚
$(U, R, s, t, c)$ 和 $(U', R', s', t', c')$，使得
$\mathcal{X} = [U/R]$ 且 $\mathcal{X}' = [U'/R']$；参见引理
\ref{stacks-more-morphisms-lemma-well-nigh-affine} 的证明。
% P12-ERR-0113: see ../control/ERRATA_CANDIDATES.jsonl
可见下列图表
$$
\xymatrix{
R' \ar[d]_{s'} \ar[r]_f & R \ar[d]^s \\
U' \ar[r]^f & U
}
\quad
\quad
\xymatrix{
R' \ar[d]_{t'} \ar[r]_f & R \ar[d]^t \\
U' \ar[r]^f & U
}
\quad
\quad
\xymatrix{
G' \ar[d] \ar[r]_f & G \ar[d] \\
U' \ar[r]^f & U
}
$$
都是笛卡儿的，其中 $G$ 和 $G'$ 是稳定子群概形。前两个图表的断言
由纤维积的传递性得出；最后一个图表的断言则是因为它是同构
$\mathcal{I}_{\mathcal{X}'} \to
\mathcal{X}' \times_\mathcal{X} \mathcal{I}_\mathcal{X}$ 的拉回
（由已经使用过的《叠的态射》中的引理
\ref{stacks-morphisms-lemma-aut-iso-unramified}）。回顾引理
\ref{stacks-more-morphisms-lemma-well-nigh-affine-moduli-space} 中的构造：$M$，相应地，\ $M'$，
分别是 $U$ 上 $R$-不变函数环的谱和相应地，\ $U'$ 上
$R'$-不变函数环的谱。
因此，由《群胚》中的引理 \ref{groupoids-lemma-etale}，
$M' \to M$ 是 \'etale 的，并且 $U' = M' \times_M U$。
由此可得
% P12-ERR-0114: see ../control/ERRATA_CANDIDATES.jsonl
$R' = M' \times_M U$；换言之，群胚 $(U', R', s', t', c')$ 是
$(U, R, s, t, c)$ 沿 $M' \to M$ 的基变换。由《群胚的商》中的引理
\ref{groupoids-quotients-lemma-base-change-quotient-stack}，
这蕴含本引理中的图表为笛卡儿图表。
\end{proof}

\begin{lemma}
\label{stacks-more-morphisms-lemma-moduli-space-finite-affine}
设代数叠 $\mathcal{X}$ 近乎仿射。引理
\ref{stacks-more-morphisms-lemma-well-nigh-affine-moduli-space} 中的态射
$$
f : \mathcal{X} \longrightarrow M
$$
是一致范畴模空间。
\end{lemma}

\begin{proof}
我们已经知道
% P12-ERR-0115: see ../control/ERRATA_CANDIDATES.jsonl
$M$ 是仿射概形范畴中的一致范畴模空间。由引理
\ref{stacks-more-morphisms-lemma-check-uniform-categorical-quotient-on-affines}，只须证明：
对于仿射概形的任意平坦态射 $M' \to M$，基变换
$f' : M' \times_M \mathcal{X} \to M'$ 是范畴模空间。
由引理 \ref{stacks-more-morphisms-lemma-affine-over-well-nigh-affine}，注意到
$\mathcal{X}' = M' \times_M \mathcal{X}$ 近乎仿射。
% P12-ERR-0116: see ../control/ERRATA_CANDIDATES.jsonl
如此，以 $\mathcal{X}'$ 代替 $\mathcal{X}$ 并以 $M'$ 代替 $M$ 后，
问题约化为证明 $f$ 是范畴模空间。

\medskip\noindent
设 $g : \mathcal{X} \to Y$ 是态射，其中 $Y$ 是代数空间。必须证明
$g = h \circ f$，其中 $h : M \to Y$ 是唯一的态射。

\medskip\noindent
唯一性。假设有两个态射 $h_i : M \to Y$，满足
$g = h_1 \circ f = h_2 \circ f$。设 $M' \subset M$ 是 $h_1$ 和
$h_2$ 的等化子。那么 $M' \to M$ 是单态射，并且
$f : \mathcal{X} \to M$ 经由 $M'$ 分解。因此 $M' \to M$
是万有同胚。由此可知 $M'$ 仿射（《态射》中的引理
\ref{morphisms-lemma-universal-homeomorphism}）。但由于
$f : \mathcal{X} \to M$ 是仿射概形范畴中的范畴模空间，
于是 $M' = M$。

\medskip\noindent
存在性。下面将证明：对每个 $p \in M$，存在笛卡儿方形
$$
\xymatrix{
\mathcal{X}' \ar[r] \ar[d] & \mathcal{X} \ar[d] \\
M' \ar[r] & M
}
$$
其中 $M' \to M$ 是仿射概形之间的 \'etale 态射，$p$ 位于其像中，
并且复合 $\mathcal{X}' \to \mathcal{X} \to Y$ 经由 $M'$ 分解。
这意味着可以在 $M$ 上 \'etale 局部地构造映射 $h : M \to Y$。
由于 $Y$ 是 \'etale 拓扑的层，再由以上证明的唯一性，这便足够了
（略去一个小细节）。

\medskip\noindent
设 $y \in |Y|$ 是 $p$ 的像。设
$(V, v) \to (Y, y)$ 是 \'etale 态射，其中 $V$ 仿射。考虑
$\mathcal{X}' = V \times_Y \mathcal{X}$。注意，
$\mathcal{X}' \to \mathcal{X}$ 作为 $V \to Y$ 的基变换，是分离且
\'etale 的。此外，$\mathcal{X}' \to \mathcal{X}$ 在自同构群上诱导
同构（《叠的态射》中的备注
\ref{stacks-morphisms-remark-identify-automorphism-groups}），因为
$V \to Y$ 具有这一性质；参见《叠的态射》中的引理
\ref{stacks-morphisms-lemma-base-change-stabilizer-preserving}。
选取如引理 \ref{stacks-more-morphisms-lemma-well-nigh-affine} 中的表现
$\mathcal{X} = [U/R]$。置
$U' = \mathcal{X}' \times_\mathcal{X} U = V \times_Y U$，并选取
$u' \in U'$，使它映到 $p$ 和 $v$（这是可行的，由《空间的性质》中的
引理 \ref{spaces-properties-lemma-points-cartesian}）。由于
$U' \to U$ 分离且 \'etale，由《态射进阶》中的引理
\ref{more-morphisms-lemma-separated-locally-quasi-finite-over-affine}，
可知 $U'$ 的任意有限点集都包含在一个仿射开集中。另一方面，态射
$U' \to \mathcal{X}'$ 满、有限、平坦且局部有限表现。置
$R' = U' \times_{\mathcal{X}'} U'$，可知
$s', t' : R' \to U'$ 有限局部自由。由《群胚》中的引理
\ref{groupoids-lemma-find-invariant-affine}，存在包含 $u'$ 的
$R'$-不变仿射开子概形 $U'' \subset U'$。设
$\mathcal{X}'' \subset \mathcal{X}'$ 是相应的开子叠。那么
$\mathcal{X}''$ 近乎仿射。由引理
\ref{stacks-more-morphisms-lemma-well-nigh-affine-moduli-space-etale}，得到笛卡儿方形
$$
\xymatrix{
\mathcal{X}'' \ar[r] \ar[d] & \mathcal{X} \ar[d] \\
M'' \ar[r] & M
}
$$
其中 $M'' \to M$ 是 \'etale 的。由于
$\mathcal{X}'' \to M''$ 是仿射概形范畴中的范畴模空间，得到态射
$M'' \to V$，使复合 $\mathcal{X}'' \to \mathcal{X}' \to V$
等于复合 $\mathcal{X}'' \to M'' \to V$。这证明了断言并完成证明。
\end{proof}

\begin{lemma}
\label{stacks-more-morphisms-lemma-etale-separated-over-well-nigh-affine}
设 $h : \mathcal{X}' \to \mathcal{X}$ 是代数叠的态射。假设
$\mathcal{X}$ 近乎仿射，$h$ 是 \'etale 的，$h$ 分离，并且 $h$ 在自同构群
上诱导同构（《叠的态射》中的备注
\ref{stacks-morphisms-remark-identify-automorphism-groups}）。
那么存在笛卡儿图表
$$
\xymatrix{
\mathcal{X}' \ar[d] \ar[r] & \mathcal{X} \ar[d] \\
M' \ar[r] & M
}
$$
其中 $M' \to M$ 是概形之间分离的 \'etale 态射，并且
$\mathcal{X} \to M$ 是引理
\ref{stacks-more-morphisms-lemma-well-nigh-affine-moduli-space} 中构造的模空间。
\end{lemma}

\begin{proof}
选取仿射概形 $U$ 和一个满、平坦、有限且局部有限表现的态射
$U \to \mathcal{X}$。由于 $h$ 可由代数空间表示
（《叠的态射》中的诸引理
\ref{stacks-morphisms-lemma-aut-iso-unramified} 和
\ref{stacks-morphisms-lemma-stabilizer-preserving}），可知
$U' = \mathcal{X}' \times_\mathcal{X} U$ 是代数空间。由于
$U' \to U$ 分离且 \'etale，可知 $U'$ 是概形，并且 $U'$ 的任意有限
点集都包含在一个仿射开集中；参见《空间的态射》中的引理
\ref{spaces-morphisms-lemma-locally-quasi-finite-separated-representable}
和《态射进阶》中的引理
\ref{more-morphisms-lemma-separated-locally-quasi-finite-over-affine}。
置 $R' = U' \times_{\mathcal{X}'} U'$，可知
$s', t' : R' \to U'$ 有限局部自由。由《群胚》中的引理
\ref{groupoids-lemma-find-invariant-affine}，存在由 $R'$-不变仿射开
子概形 $U'_i \subset U'$ 组成的开覆盖
$U' = \bigcup U'_i$。设 $\mathcal{X}'_i \subset \mathcal{X}'$
是相应的开子叠。它们近乎仿射，因为
$U'_i \to \mathcal{X}'_i$ 满、平坦、有限且有限表现。由引理
\ref{stacks-more-morphisms-lemma-well-nigh-affine-moduli-space-etale}，得到笛卡儿图表
$$
\xymatrix{
\mathcal{X}'_i \ar[r] \ar[d] & \mathcal{X} \ar[d] \\
M'_i \ar[r] & M
}
$$
其中 $M'_i \to M$ 是仿射概形之间的 \'etale 态射，而竖直箭头如引理
\ref{stacks-more-morphisms-lemma-well-nigh-affine-moduli-space} 所述。注意
$\mathcal{X}'_{ij} = \mathcal{X}'_i \cap \mathcal{X}'_j$
% P12-ERR-0117: see ../control/ERRATA_CANDIDATES.jsonl
是 $\mathcal{X}'_i$ 和 $\mathcal{X}'_j$ 的开子空间。因此得到相应的
开子概形 $V_{ij} \subset M'_i$ 和 $V_{ji} \subset M'_j$。
由引理 \ref{stacks-more-morphisms-lemma-moduli-space-finite-affine} 的结论，可知
$\mathcal{X}'_{ij} \to V_{ij}$ 和
$\mathcal{X}'_{ji} \to V_{ji}$ 都是范畴模空间！因此得到唯一同构
$\varphi_{ij} : V_{ij} \to V_{ji}$，使得
$$
\xymatrix{
\mathcal{X}'_i \ar[d] & &
\mathcal{X}'_i \cap \mathcal{X}'_j \ar[rr] \ar[ll] \ar[ld] \ar[rd] & &
\mathcal{X}'_j \ar[d] \\
M'_i &
V_{ij} \ar[l] \ar[rr]^{\varphi_{ij}} & &
V_{ji} \ar[r] &
M'_j
}
$$
交换。这些同构满足《概形》第
% P12-ERR-0118: see ../control/ERRATA_CANDIDATES.jsonl
\ref{schemes-section-glueing-schemes} 节的余圈条件；这一点可由一个计算
（以及再次应用前一引理）得出，我们略去该计算。因此，可以把这些
% P12-ERR-0119: see ../control/ERRATA_CANDIDATES.jsonl
仿射概形黏合成概形 $M'$；参见《概形》中的引理
\ref{schemes-lemma-glue}。把 $M'_i$ 与它在 $M'$ 中的像等同起来。
我们断言存在态射 $\mathcal{X}' \to M'$，使它嵌入笛卡儿图表
$$
\xymatrix{
\mathcal{X}'_i \ar[r] \ar[d] & \mathcal{X}' \ar[d] \\
M'_i \ar[r] & M'
}
$$
这由《概形》中的引理 \ref{schemes-lemma-glue} 对到黏合概形 $M'$ 的
态射之描述以及下述事实显然可见：给出态射
% P12-ERR-0120: see ../control/ERRATA_CANDIDATES.jsonl
$\mathcal{X}' \to M'$ 等价于对任意态射 $T \to \mathcal{X}'$ 给定态射
$T \to M'$。类似地，存在态射 $M' \to M$，其在 $M'_i$ 上的限制是
给定态射 $M'_i \to M$。态射 $M' \to M$ 是 \'etale 的
（因为它在一个 \'etale 覆盖的成员上是 \'etale 的），而纤维积性质
成立，因为可以在（仿射）开覆盖 $M' = \bigcup M'_i$ 的成员上检验它。
最后，$M' \to M$ 分离，因为复合
$U' \to \mathcal{X}' \to M'$ 满且万有闭，所以可以应用《态射》中的
引理 \ref{morphisms-lemma-image-universally-closed-separated}。
\end{proof}

\begin{lemma}
\label{stacks-more-morphisms-lemma-etale-local-finite-inertia}
设 $\mathcal{X}$ 是代数叠。假设
$\mathcal{I}_\mathcal{X} \to \mathcal{X}$ 有限。那么存在集合 $I$，
并且对每个 $i \in I$，存在代数叠的态射
$$
g_i : \mathcal{X}_i \longrightarrow \mathcal{X}
$$
具有下列性质：
\begin{enumerate}
\item $|\mathcal{X}| = \bigcup |g_i|(|\mathcal{X}_i|)$；
\item $\mathcal{X}_i$ 近乎仿射；
\item $\mathcal{I}_{\mathcal{X}_i} \to
\mathcal{X}_i \times_\mathcal{X} \mathcal{I}_\mathcal{X}$
是同构；
\item $g_i : \mathcal{X}_i \to \mathcal{X}$ 可由代数空间表示、
分离且 \'etale。
\end{enumerate}
\end{lemma}

\begin{proof}
对任意 $x \in |\mathcal{X}|$，可以像《叠的态射》中的引理
\ref{stacks-morphisms-lemma-etale-local-quasi-DM-at-x} 那样选取
$g : \mathcal{U} \to \mathcal{X}$、$\mathcal{U} = [U/R]$ 和 $u$。
于是，由《叠的态射》中的引理
\ref{stacks-morphisms-lemma-stabilizer-preserving-unramified}，
存在包含 $u$ 的开子叠 $\mathcal{U}' \subset \mathcal{U}$，使得
$\mathcal{I}_{\mathcal{U}'} \to
\mathcal{U}' \times_\mathcal{X} \mathcal{I}_\mathcal{X}$
是同构。设 $U' \subset U$ 是与开子叠 $\mathcal{U}'$ 对应的
$R$-不变开集。设 $u' \in U'$ 是 $U'$ 中映到 $u$ 的一点。
注意，由于 $s : R \to U$ 有限，$t(s^{-1}(\{u'\}))$ 是有限集。
由《性质》中的引理 \ref{properties-lemma-ample-finite-set-in-affine}
和《群胚》中的引理 \ref{groupoids-lemma-find-invariant-affine}，
可以找到包含 $u'$ 的 $R$-不变仿射开集 $U'' \subset U'$。
设 $R''$ 是 $R$ 在 $U''$ 上的限制。那么
$\mathcal{U}'' = [U''/R'']$ 是 $\mathcal{U}'$ 中包含 $u$ 的开子叠，
它近乎仿射，并且
$\mathcal{I}_{\mathcal{U}''} \to
\mathcal{U}'' \times_\mathcal{X} \mathcal{I}_\mathcal{X}$
是同构，而 $\mathcal{U}'' \to \mathcal{X}$
% P12-ERR-0121: see ../control/ERRATA_CANDIDATES.jsonl
可由代数空间表示且 \'etale。最后，$\mathcal{U}'' \to \mathcal{X}$
分离，因为 $\mathcal{U}''$ 分离（引理
\ref{stacks-more-morphisms-lemma-well-nigh-affine}），
% P12-ERR-0122: see ../control/ERRATA_CANDIDATES.jsonl
$\mathcal{X}$ 的对角态射分离（《叠的态射》中的引理
\ref{stacks-morphisms-lemma-diagonal-diagonal}），并且分离性由
《叠的态射》中的引理
\ref{stacks-morphisms-lemma-compose-after-separated} 得出。
由于点 $x \in |\mathcal{X}|$ 是任意的，证明完成。
\end{proof}

\begin{theorem}[Keel-Mori]
\label{stacks-more-morphisms-theorem-keel-mori}
设 $\mathcal{X}$ 是代数叠。假设
$\mathcal{I}_\mathcal{X} \to \mathcal{X}$ 有限。那么存在一致范畴模空间
$$
f : \mathcal{X} \longrightarrow M
$$
并且 $f$ 分离、拟紧且为万有同胚。
\end{theorem}

\begin{proof}
选取集合 $I$\footnote{仍在留意集合论问题的读者应当确认 $I$
不会过大。}，并对每个 $i \in I$ 选取如引理
\ref{stacks-more-morphisms-lemma-etale-local-finite-inertia} 中的代数叠态射
$g_i : \mathcal{X}_i \to \mathcal{X}$；我们将不再特别说明地使用该引理
所列的全部性质。令
$$
f_i : \mathcal{X}_i \to M_i
$$
如引理 \ref{stacks-more-morphisms-lemma-well-nigh-affine-moduli-space} 所述。对
$i, j \in I$，考虑叠
$$
\mathcal{X}_{ij} = \mathcal{X}_i \times_{g_i, \mathcal{X}, g_j} \mathcal{X}_j
$$
由《叠的态射》中的引理
\ref{stacks-morphisms-lemma-base-change-separated}，投影
$\mathcal{X}_{ij} \to \mathcal{X}_i$ 和
$\mathcal{X}_{ij} \to \mathcal{X}_j$ 分离；由《叠的态射》中的引理
\ref{stacks-morphisms-lemma-base-change-etale}，它们是 \'etale 的；
它们还在自同构群上诱导同构（如《叠的态射》中的备注
\ref{stacks-morphisms-remark-identify-automorphism-groups}），这一点由
《叠的态射》中的引理
\ref{stacks-morphisms-lemma-base-change-stabilizer-preserving} 得出。因此，可以应用
引理 \ref{stacks-more-morphisms-lemma-etale-separated-over-well-nigh-affine} 得到交换图表
$$
\xymatrix{
\mathcal{X}_i \ar[d]_{f_i} &
\mathcal{X}_{ij} \ar[d]_{f_{ij}} \ar[l] \ar[r] &
\mathcal{X}_j \ar[d]_{f_j} \\
M_i &
M_{ij} \ar[l] \ar[r] &
M_j
}
$$
其中两个方形都是笛卡儿方形，而 $M_{ij} \to M_i$ 和
$M_{ij} \to M_j$ 是概形之间分离的 \'etale 态射；这里还由引理
\ref{stacks-more-morphisms-lemma-moduli-space-finite-affine} 使用了
% P12-ERR-0123: see ../control/ERRATA_CANDIDATES.jsonl
$f_i$ 是一致范畴商这一点。断言：
$$
\coprod M_{ij} \longrightarrow \coprod M_i \times \coprod M_i
$$
是 \'etale 等价关系。

\medskip\noindent
证明该断言。置 $R = \coprod M_{ij}$ 和 $U = \coprod M_i$。
已经看到 $t : R \to U$ 和 $s : R \to U$ 是 \'etale 的。
我们构造与
$\text{pr}_{13} : U \times U \times U \to U \times U$ 相容的态射
$c : R \times_{s, U, t} R \to R$。具体而言，对
$i, j, k \in I$，考虑
$$
\mathcal{X}_{ijk} =
\mathcal{X}_i \times_{g_i, \mathcal{X}, g_j} \mathcal{X}_j
\times_{g_j, \mathcal{X}, g_k} \mathcal{X}_k =
\mathcal{X}_{ij} \times_{\mathcal{X}_j} \mathcal{X}_{jk}
$$
完全沿用上一段的论证，可知
$M_{ijk} = M_{ij} \times_{M_j} M_{jk}$ 是
$\mathcal{X}_{ijk}$ 的范畴模空间。特别地，由投影
$\mathcal{X}_{ijk} \to \mathcal{X}_{ik}$ 得到典范态射
$M_{ijk} = M_{ij} \times_{M_j} M_{jk} \to M_{ik}$。
把这些态射合在一起，便得到态射 $c$。类似地，构造与
$\Delta : U \to U \times U$ 相容的态射 $e : U \to R$，以及与对换
$U \times U \to U \times U$ 相容的态射 $i : R \to R$。
设 $k$ 是代数闭域。那么
$$
\Mor(\Spec(k), \mathcal{X}_i) \to \Mor(\Spec(k), M_i) = M_i(k)
$$
在同构类上是双射，并且沿任意态射
% P12-ERR-0124: see ../control/ERRATA_CANDIDATES.jsonl
$M' \to M$ 作基变换后仍然如此。这由 $f_i$ 的选取以及《叠的态射》
中的诸引理 \ref{stacks-morphisms-lemma-universally-injective} 和
\ref{stacks-morphisms-lemma-universally-injective-point} 得出。
由 $2$-纤维积的构造，图表
$$
\xymatrix{
\Mor(\Spec(k), \mathcal{X}_{ij}) \ar[d] \ar[r] &
\Mor(\Spec(k), \mathcal{X}_j) \ar[d] \\
\Mor(\Spec(k), \mathcal{X}_i) \ar[r] &
\Mor(\Spec(k), \mathcal{X})
}
$$
是范畴的纤维积。由 $g_i$ 的选取，该图表中的函子在自同构群上诱导
双射。因此，该图表在同构类集合上诱导纤维积图表！于是
$$
R(k) = U(k) \times_{|\Mor(\Spec(k), \mathcal{X})|} U(k)
$$
其中 $|\Mor(\Spec(k), \mathcal{X})|$ 表示同构类集合。特别地，
对任意代数闭域 $k$，
% P12-ERR-0125: see ../control/ERRATA_CANDIDATES.jsonl
在 $k$-值点上的映射是等价关系。由《群胚》中的引理
\ref{groupoids-lemma-etale-equivalence-relation}，断言得证。

\medskip\noindent
设 $M = U/R$ 为上述 \'etale 等价关系的商代数空间；参见《空间》中的
定理 \ref{spaces-theorem-presentation}。存在典范态射
$f : \mathcal{X} \to M$，它嵌入交换图表
\begin{equation}
\label{stacks-more-morphisms-equation-fundamental-diagram}
\xymatrix{
\mathcal{X}_i \ar[r]_{g_i} \ar[d]_{f_i} & \mathcal{X} \ar[d]^f \\
M_i \ar[r] & M
}
\end{equation}
具体而言，这样的态射 $f$ 由对任意概形 $T$ 都与基变换相容的函子
$$
f : \Mor(T, \mathcal{X}) \longrightarrow \Mor(T, M)
$$
给出。设 $a : T \to \mathcal{X}$ 是左端的一个对象。得到 \'etale 覆盖
$\{T_i \to T\}$，其中
$T_i = \mathcal{X}_i \times_\mathcal{X} T$，以及态射
$a_i : T_i \to \mathcal{X}_i$。于是得到
$b_i = f_i \circ a_i : T_i \to M_i$。由于
$T_i \times_T T_j = \mathcal{X}_{ij} \times_\mathcal{X} T$，
还得到态射 $a_{ij} : T_i \times_T T_j \to \mathcal{X}_{ij}$。
置 $b_{ij} = f_{ij} \circ a_{ij}$，可知 $b_i \times b_j$
经由单态射 $M_{ij} \to M_i \times M_j$ 分解。因此，诸态射
$$
T_i \xrightarrow{b_i} M_i \to M
$$
在 $T_i \times_T T_j$ 上相同。由于 $M$ 是 \'etale 拓扑的层，
这些态射黏合成唯一态射 $b = f(a) : T \to M$。我们略去验证该构造
与基变换相容，也略去验证图表
(\ref{stacks-more-morphisms-equation-fundamental-diagram}) 交换。

\medskip\noindent
断言：图表 (\ref{stacks-more-morphisms-equation-fundamental-diagram}) 是笛卡儿的。
为看出这一点，考察诱导态射
$$
h_i : \mathcal{X}_i \longrightarrow M_i \times_M \mathcal{X}
$$
这是 $\mathcal{X}$ 上叠的 \'etale 态射，所以 $h_i$ 是 \'etale 的
（《叠的态射》中的引理
\ref{stacks-morphisms-lemma-etale-permanence}）。由于 $g_i$ 分离，
可知 $h_i$ 分离（使用《叠的态射》中的引理
\ref{stacks-morphisms-lemma-compose-after-separated} 以及以上看到的
$\mathcal{X}$ 的对角态射分离这一事实）。态射 $h_i$ 在自同构群上诱导
同构（《叠的态射》中的备注
\ref{stacks-morphisms-remark-identify-automorphism-groups}），
因为 $g_i$ 具有这一性质。对代数闭域 $k$，图表
$$
\xymatrix{
\Mor(\Spec(k), M_i \times_M \mathcal{X}) \ar[r] \ar[d] &
\Mor(\Spec(k), \mathcal{X}) \ar[d] \\
M_i(k) \ar[r] &
M(k)
}
$$
% P12-ERR-0126: see ../control/ERRATA_CANDIDATES.jsonl
是范畴的笛卡儿图表，并且上方箭头在自同构群上诱导双射。另一方面，
由以上所述，有
$$
M(k) = U(k)/R(k) = U(k)/
U(k) \times_{|\Mor(\Spec(k), \mathcal{X})|} U(k) =
|\Mor(\Spec(k), \mathcal{X})|
$$
因此，上方笛卡儿图表的右竖直箭头在同构类上是双射。可知
$|\Mor(\Spec(k), M_i \times_M \mathcal{X})| \to M_i(k)$ 是双射。
总结：
% P12-ERR-0127: see ../control/ERRATA_CANDIDATES.jsonl
$h_i$ 分离、\'etale，在自同构群上诱导同构（如《叠的态射》中的备注
\ref{stacks-morphisms-remark-identify-automorphism-groups}），并且在代数闭域
上的纤维范畴上诱导等价。因此，由《叠的态射》中的引理
\ref{stacks-morphisms-lemma-etale-iso}，它是同构。

\medskip\noindent
特别地，由该断言得到：存在满 \'etale 态射 $U \to M$，使得 $f$ 的
基变换分离、拟紧且为万有同胚。由此可知 $f$ 分离、拟紧且为万有同胚。
参见《叠的态射》中的诸引理
\ref{stacks-morphisms-lemma-check-separated-covering}、
\ref{stacks-morphisms-lemma-check-quasi-compact-covering} 和
% P12-ERR-0128: see ../control/ERRATA_CANDIDATES.jsonl
\ref{stacks-morphisms-lemma-check-universal-homeomorphism-covering}

\medskip\noindent
为完成证明，必须证明 $f : \mathcal{X} \to M$ 是一致范畴模空间。
为此，只须证明：给定代数空间的平坦态射 $M' \to M$，基变换
$$
M' \times_M \mathcal{X} \longrightarrow M'
$$
是范畴模空间。为此考虑态射
$$
\theta : M' \times_M \mathcal{X} \longrightarrow E
$$
其中 $E$ 是代数空间。对每个 $i$，已知 $f_i$ 是一致范畴模空间。
因此得到
$$
\xymatrix{
M' \times_M \mathcal{X}_i \ar[d] \ar[r] &
M' \times_M \mathcal{X} \ar[d]^\theta \\
M' \times_M M_i \ar[r]^{\psi_i} &
E
}
$$
由于 $\{M' \times_M M_i \to M'\}$ 是 \'etale 覆盖，为得到所需态射
$\psi : M' \to E$，只须证明 $\psi_i$ 和 $\psi_j$ 在
$M' \times_M M_i \times_M M_j = M' \times_M M_{ij}$ 上相同。
这容易由下述事实得出：
% P12-ERR-0129: see ../control/ERRATA_CANDIDATES.jsonl
$f_{ij} : \mathcal{X}_{ij} =
\mathcal{X}_i \times_\mathcal{X} \mathcal{X}_j \to M_{ij}$ 是一致范畴商；
略去细节。最后，再证明 $\psi$ 嵌入交换图表
$$
\xymatrix{
M' \times_M \mathcal{X} \ar[d] \ar[rd]^\theta \\
M' \ar[r]^\psi &
E
}
$$
因为“$\{M' \times_M \mathcal{X}_i \to M' \times_M \mathcal{X}\}$
是一个 \'etale 覆盖”，而诸态射 $\psi_i$ 按照构造嵌入相应的交换图表。
这完成了 Keel--Mori 定理的证明。
\end{proof}

\noindent
下一个引理强调该构造的 \'etale 局部性质。

\begin{lemma}
\label{stacks-more-morphisms-lemma-etale-separated-over-keel-mori}
设 $h : \mathcal{X}' \to \mathcal{X}$ 是代数叠的态射。假设
\begin{enumerate}
\item $\mathcal{I}_\mathcal{X} \to \mathcal{X}$ 有限；
\item $h$ 是 \'etale 且分离的，并且在自同构群上诱导同构
（《叠的态射》中的备注
\ref{stacks-morphisms-remark-identify-automorphism-groups}）。
\end{enumerate}
那么存在笛卡儿图表
$$
\xymatrix{
\mathcal{X}' \ar[d] \ar[r] &
\mathcal{X} \ar[d] \\
M' \ar[r] &
M
}
$$
其中 $M' \to M$ 是代数空间之间分离的 \'etale 态射，而竖直箭头是
定理 \ref{stacks-more-morphisms-theorem-keel-mori} 中构造的模空间。
\end{lemma}

\begin{proof}
由《叠的态射》中的引理
\ref{stacks-morphisms-lemma-aut-iso-unramified}，可知
$\mathcal{I}_{\mathcal{X}'} \to
\mathcal{X}' \times_\mathcal{X} \mathcal{I}_\mathcal{X}$
是同构。因此，$\mathcal{I}_{\mathcal{X}'} \to \mathcal{X}'$
作为 $\mathcal{I}_\mathcal{X} \to \mathcal{X}$ 的基变换而有限。
设 $f' : \mathcal{X}' \to M'$ 和 $f : \mathcal{X} \to M$ 如定理
\ref{stacks-more-morphisms-theorem-keel-mori} 所述。因为
% P12-ERR-0130: see ../control/ERRATA_CANDIDATES.jsonl
$f'$ 是范畴模空间，得到本引理中的交换图表。像引理
\ref{stacks-more-morphisms-lemma-etale-local-finite-inertia} 中那样选取 $I$ 和
$g'_i : \mathcal{X}'_i \to \mathcal{X}'$。注意，
$g_i = h \circ g'_i$ 是 \'etale 且分离的，并在自同构群上诱导同构
（《叠的态射》中的备注
\ref{stacks-morphisms-remark-identify-automorphism-groups}）。设
$f'_i : \mathcal{X}'_i \to M'_i$ 如引理
\ref{stacks-more-morphisms-lemma-well-nigh-affine-moduli-space} 所述。在定理
\ref{stacks-more-morphisms-theorem-keel-mori} 的证明中已经看到，图表
$$
\xymatrix{
\mathcal{X}'_i \ar[d]_{f'_i} \ar[r]_{g'_i} &
\mathcal{X}' \ar[d]^{f'} \\
M'_i \ar[r] &
M'
}
\quad\text{以及}\quad
\xymatrix{
\mathcal{X}'_i \ar[d]_{f'_i} \ar[r]_{g_i} &
\mathcal{X} \ar[d]^f \\
M'_i \ar[r] &
M
}
$$
都是笛卡儿的，并且 $M'_i \to M'$ 和 $M'_i \to M$ 是 \'etale 的
（完全以同样方式论证，这也直接由
% P12-ERR-0131: see ../control/ERRATA_CANDIDATES.jsonl
迄今所列态射 $g'_i, g_i, f', f'_i, f$ 的性质得出）。这首先蕴含
$M' \to M$ 是 \'etale 的，其次蕴含本引理中的图表是笛卡儿的。
仍须证明 $M' \to M$ 分离。为此考察图表
$$
\xymatrix{
\mathcal{X}' \ar[r] \ar[d] &
\mathcal{X}' \times_\mathcal{X} \mathcal{X}' \ar[d] \\
M' \ar[r] &
M' \times_M M'
}
$$
由于 $\mathcal{X}' \to \mathcal{X}$ 分离，上方水平箭头万有闭。
竖直箭头如定理 \ref{stacks-more-morphisms-theorem-keel-mori} 中所述
（是 $\mathcal{X} \to M$ 的平坦基变换），因而是万有同胚。
所以，下方水平箭头万有闭。这一点与它是代数空间的 \'etale 单态射
结合，证明它如所需为闭浸入。
\end{proof}

\section{模空间的性质}
\label{stacks-more-morphisms-section-properties-moduli-spaces}

\noindent
一旦证明模空间存在，就可以（而且通常很直接地）确定这些模空间的性质。

\begin{lemma}
\label{stacks-more-morphisms-lemma-keel-mori-finite-type}
设 $p : \mathcal{X} \to Y$ 是从代数叠到代数空间的态射。假设
\begin{enumerate}
\item $\mathcal{I}_\mathcal{X} \to \mathcal{X}$ 有限；
\item $Y$ 局部 Noether；
\item $p$ 局部有限型。
\end{enumerate}
设 $f : \mathcal{X} \to M$ 是定理
\ref{stacks-more-morphisms-theorem-keel-mori} 中构造的模空间。那么 $M \to Y$ 局部有限型。
\end{lemma}

\begin{proof}
由于 $f$ 是一致范畴模空间，得到态射 $M \to Y$。只须在 $M$ 和 $Y$
上 \'etale 局部地检验 $M \to Y$ 局部有限型。由于 $f$ 是一致范畴
模空间，可以先用一个在 $Y$ 上 \'etale 的仿射概形代替 $Y$。
接着，可以像引理 \ref{stacks-more-morphisms-lemma-etale-local-finite-inertia} 中那样选取
$I$ 和 $g_i : \mathcal{X}_i \to \mathcal{X}$。然后由引理
\ref{stacks-more-morphisms-lemma-etale-separated-over-keel-mori}，问题约化为
$\mathcal{X} = \mathcal{X}_i$ 的情形。换言之，可以假设
$\mathcal{X}$ 近乎仿射。在这种情形下，
$Y = \Spec(A_0)$，$\mathcal{X} = [U/R]$ 且
$U = \Spec(A)$，并且 $M = \Spec(C)$，其中 $C \subset A$ 是
$U$ 上所有 $R$-不变函数组成的集合。参见引理
\ref{stacks-more-morphisms-lemma-well-nigh-affine} 和
\ref{stacks-more-morphisms-lemma-well-nigh-affine-moduli-space}。于是 $A_0$ 是 Noether 环，
而 $A_0 \to A$ 有限型。此外，由《群胚》中的引理
\ref{groupoids-lemma-integral-over-invariants}，$A$ 在 $C$ 上整，
因而在 $C$ 上有限（因为它在 $A_0$ 上有限型）。最后应用《代数》中的
引理 \ref{algebra-lemma-Artin-Tate} 即得结论。
\end{proof}

\begin{lemma}
\label{stacks-more-morphisms-lemma-keel-mori-diagonal}
设 $\mathcal{X}$ 是代数叠。假设
$\mathcal{I}_\mathcal{X} \to \mathcal{X}$ 有限。设
$f : \mathcal{X} \to M$ 是定理
\ref{stacks-more-morphisms-theorem-keel-mori} 中构造的模空间。
\begin{enumerate}
\item 如果 $\mathcal{X}$ 拟分离，那么 $M$ 拟分离。
\item 如果 $\mathcal{X}$ 分离，那么 $M$ 分离。
% P12-ERR-0132: see ../control/ERRATA_CANDIDATES.jsonl
\item 此处还应添加更多内容，例如上述结论的相对版本。
\end{enumerate}
\end{lemma}

\begin{proof}
为证明这些断言，考虑图表
$$
\xymatrix{
\mathcal{X} \ar[d]_f \ar[r]_{\Delta_\mathcal{X}} &
\mathcal{X} \times \mathcal{X} \ar[d]^{f \times f} \\
M \ar[r]^{\Delta_M} &
M \times M
}
$$
由于 $f$ 是万有同胚，可知 $f \times f$ 是万有同胚。

\medskip\noindent
如果 $\mathcal{X}$ 分离，那么 $\Delta_\mathcal{X}$ 固有，故
$\Delta_\mathcal{X}$ 万有闭，进而 $\Delta_M$ 万有闭，所以由
《空间的态射》中的引理
\ref{spaces-morphisms-lemma-separated-diagonal-proper}，$M$ 分离。

\medskip\noindent
如果 $\mathcal{X}$ 拟分离，那么 $\Delta_\mathcal{X}$ 拟紧，
故 $\Delta_M$ 拟紧，所以 $M$ 拟分离。
\end{proof}

\begin{lemma}
\label{stacks-more-morphisms-lemma-keel-mori-proper}
设 $p : \mathcal{X} \to Y$ 是从代数叠到代数空间的态射。假设
\begin{enumerate}
\item $\mathcal{I}_\mathcal{X} \to \mathcal{X}$ 有限；
\item $p$ 固有；
\item $Y$ 局部 Noether。
\end{enumerate}
设 $f : \mathcal{X} \to M$ 是定理
\ref{stacks-more-morphisms-theorem-keel-mori} 中构造的模空间。那么 $M \to Y$ 固有。
\end{lemma}

\begin{proof}
由引理 \ref{stacks-more-morphisms-lemma-keel-mori-finite-type}，可知 $M \to Y$ 局部有限型。
由引理 \ref{stacks-more-morphisms-lemma-keel-mori-diagonal}，可知 $M \to Y$ 分离。
当然，$M \to Y$ 拟紧且万有闭，因为这些是拓扑性质，
$\mathcal{X} \to Y$ 具有这些性质，而
$\mathcal{X} \to M$ 是万有同胚。
\end{proof}

\section{叠与 fpqc 覆盖}
\label{stacks-more-morphisms-section-fpqc}

\noindent
某些代数叠满足 fpqc 下降。本节对代数空间的对应结论见《空间的性质》
第 \ref{spaces-properties-section-fpqc} 节。


\begin{proposition}
\label{stacks-more-morphisms-proposition-stack-fpqc}
\begin{reference}
Anatoly Preygel《代数叠导论》命题 3.3.6。
\end{reference}
设 $\mathcal{X}$ 是对角态射拟仿射\footnote{假设对角态射为
ind-拟仿射便已足够。}的代数叠。那么
$\mathcal{X}$ 对 fpqc 覆盖满足下降。
\end{proposition}

\begin{proof}
我们的约定是：$\mathcal{X}$ 是基概形 $S$ 上概形范畴（赋予 fppf 拓扑）
之上的群胚叠 $p : \mathcal{X} \to (\Sch/S)_{fppf}$。该陈述的含义如下：
给定 $S$ 上概形的 fpqc 覆盖
$\mathcal{U} = \{U_i \to U\}_{i \in I}$，函子
$$
\mathcal{X}_U \longrightarrow DD(\mathcal{U})
$$
是等价。左边是 $\mathcal{X}$ 在 $U$ 上的对象范畴，右边是
$\mathcal{X}$ 中相对于 $\mathcal{U}$ 的下降数据范畴。参见《叠》第
\ref{stacks-section-descent-data} 节的讨论。

\medskip\noindent
% P12-ERR-0133: see ../control/ERRATA_CANDIDATES.jsonl
全忠实性。假设有 $\mathcal{X}$ 在 $U$ 上的两个对象 $x, y$。
那么 $I = \mathit{Isom}(x, y)$ 是 $U$ 上的代数空间。因此，$I$ 在
$U_i$ 上的一族截面，若其限制在 $U_i \times_U U_j$ 上相同，
则由代数空间的对应命题来自 $U$ 上的唯一截面；参见《空间的性质》中的
命题 \ref{spaces-properties-proposition-sheaf-fpqc}。所以该函子全忠实。

\medskip\noindent
本质满射性。这里给定 $U_i$ 上的对象 $x_i$ 和
$U_i \times_U U_j$ 上的同构
$\varphi_{ij} : \text{pr}_0^*x_i \to \text{pr}_1^*x_j$，它们在
% P12-ERR-0134: see ../control/ERRATA_CANDIDATES.jsonl
$U_i \times_U U_j \times_U U_k$ 上满足余圈条件。

\medskip\noindent
设 $W$ 是仿射概形，并设 $W \to \mathcal{X}$ 是态射。对每个 $i$，
可以形成
$$
W_i = U_i \times_{x_i, \mathcal{X}} W
$$
投影 $W_i \to U_i$ 拟仿射，因为 $\mathcal{X}$ 的对角态射拟仿射。
对每对 $i, j \in I$，同构 $\varphi_{ij}$ 诱导同构
$$
W_i \times_U U_j =
(U_i \times_U U_j) \times_{x_i \circ \text{pr}_0, \mathcal{X}} W
\to
(U_i \times_U U_j) \times_{x_j \circ \text{pr}_1, \mathcal{X}} W =
U_i \times_U W_j
$$
而且，这些同构在 $U_i \times_U U_j \times_U U_k$ 上满足余圈条件。
换言之，这些同构定义了概形 $W_i/U_i$ 相对于 $\mathcal{U}$ 的下降数据。
由《下降》中的引理 \ref{descent-lemma-quasi-affine}，该下降数据有效
\footnote{在对角态射为 ind-拟仿射的情形，也可使用《群胚进阶》中的
引理 \ref{more-groupoids-lemma-ind-quasi-affine}。}。因此，存在拟仿射态射
$W' \to U$ 和交换图表
$$
\xymatrix{
W' \ar[d] &
W_i \ar[l] \ar[d] \ar[r] &
W \ar[d] \\
U &
U_i \ar[l] \ar[r]^{x_i} &
\mathcal{X}
}
$$
其中两个方形都是笛卡儿的。由于 $\{W_i \to W'\}_{i \in I}$
是 $\mathcal{U}$ 沿 $W' \to U$ 的基变换，它是 fpqc 覆盖。
由于 $W$ 对 fpqc 覆盖满足层条件，得到唯一态射 $W' \to W$，
使 $W_i \to W' \to W$ 为给定态射 $W_i \to W$。换言之，
有交换图表
$$
\xymatrix{
W_i \ar[d] \ar[r] &
W' \ar[d] \ar[r] &
W \ar[d] \\
U_i \ar[r] \ar@/_1pc/[rr]_{x_i} &
U &
\mathcal{X}
}
$$
它们与同构 $\varphi_{ij}$ 相容，并且其中的方形和矩形都是笛卡儿的。

\medskip\noindent
选取一族仿射概形 $W_\alpha$（$\alpha \in A$）以及光滑态射
$W_\alpha \to \mathcal{X}$，使得
$\coprod W_\alpha \to \mathcal{X}$ 满。通过上一段的步骤，
对每个 $\alpha$ 产生图表
$$
\xymatrix{
W_{\alpha, i} \ar[d] \ar[r] &
W_\alpha' \ar[d] \ar[r] &
W_\alpha \ar[d] \\
U_i \ar[r] \ar@/_1pc/[rr]_{x_i} &
U &
\mathcal{X}
}
$$
于是诸态射 $W_\alpha' \to U$ 光滑且联合满射。

\medskip\noindent
% P12-ERR-0135: see ../control/ERRATA_CANDIDATES.jsonl
用 $x_\alpha$ 表示 $\mathcal{X}$ 在 $W_\alpha'$ 上与
$W_\alpha' \to W_\alpha \to \mathcal{X}$ 对应的对象。由于
$\mathcal{X}$ 是 fppf 叠，并且
$\{W_\alpha' \to U\}$ 是 fppf 覆盖，只须证明在
$W_\alpha' \times_U W'_\beta$ 上存在满足余圈条件的同构
$\text{pr}_0^*x_\alpha \to \text{pr}_1^*x_\beta$。然而，拉回到
$W_{\alpha, i}$ 后，确实在
$W_{\alpha, i} \times_{U_i} W_{\beta, i} =
U_i \times_U (W_\alpha' \times_U W'_\beta)$ 上有这样的同构，
因为 $x_\alpha$ 到 $W_{\alpha, i}$ 的拉回同构于 $x_i$ 到
$W_{\alpha, i}$ 的拉回。由于
$\{U_i \times_U (W_\alpha' \times_U W'_\beta) \to
W_\alpha' \times_U W'_\beta\}_{i \in I}$ 是 fpqc 覆盖，又由于
上述图表与 $\varphi_{ij}$ 的相容性，这些同构下降到
$W_\alpha' \times_U W'_\beta$，证明完成。
\end{proof}

\section{张量函子}
\label{stacks-more-morphisms-section-tensor-functors}

\noindent
设 $f : \mathcal{Y} \to \mathcal{X}$ 是 Noether 代数叠的态射。拉回函子
$$
f^* :
\textit{Coh}(\mathcal{O}_\mathcal{X})
\longrightarrow
\textit{Coh}(\mathcal{O}_\mathcal{Y})
$$
是右正合张量函子：它可加、右正合，并与凝聚模的张量积交换。
可以问，任意右正合张量函子
$F : \textit{Coh}(\mathcal{O}_\mathcal{X}) \to
\textit{Coh}(\mathcal{O}_\mathcal{Y})$ 在多大程度上来自某个态射
$f : \mathcal{Y} \to \mathcal{X}$。关于这种类型的一个非常一般的结论，
读者可查阅 \cite{Hall-Rydh-coherent}。本节的目标是简短证明定理
\ref{stacks-more-morphisms-theorem-main}，作为这些思想的入门。

\medskip\noindent
先给出几个引理。

\begin{lemma}
\label{stacks-more-morphisms-lemma-extend}
设 $\mathcal{X}$ 和 $\mathcal{Y}$ 是 Noether 代数叠。任意右正合张量函子
$F : \textit{Coh}(\mathcal{O}_\mathcal{X}) \to
\textit{Coh}(\mathcal{O}_\mathcal{Y})$ 都唯一延拓为与所有余极限交换的
右正合张量函子
$F : \QCoh(\mathcal{O}_\mathcal{X}) \to \QCoh(\mathcal{O}_\mathcal{Y})$。
\end{lemma}

\begin{proof}
这种延拓的存在性和唯一性是一般事实；参见《范畴论》中的引理
\ref{categories-lemma-extend-functor-by-colim}。为看出该引理适用，
注意局部 Noether 代数叠上的凝聚模按定义就是有限表现模；参见
《叠的上同调》中的定义
\ref{stacks-cohomology-definition-coherent}。因此，由《叠的上同调》中的
引理 \ref{stacks-cohomology-lemma-finite-presentation-quasi-compact-colimit}，
$\mathcal{X}$ 上的凝聚模是
$\QCoh(\mathcal{O}_\mathcal{X})$ 的范畴紧对象。最后，由《叠的上同调》
中的引理 \ref{stacks-cohomology-lemma-directed-colimit-coherent}，
每个拟凝聚模都是其凝聚子模的滤过余极限。

\medskip\noindent
由于 $F$ 可加，$F$ 的延拓也可加（略去细节）。由于 $F$ 是张量函子，
而模的余极限与取张量积交换，$F$ 的延拓也是张量函子（略去细节）。

\medskip\noindent
本段证明该延拓与任意直和交换。如果
$\mathcal{F} = \bigoplus_{j \in J} \mathcal{H}_j$，其中
$\mathcal{H}_j$ 拟凝聚。

于是
$\mathcal{F} = \colim_{J' \subset J\text{，且为有限集}}
\bigoplus_{j \in J'} \mathcal{H}_j$。仍用 $F$ 表示 $F$ 的延拓，得到
\begin{align*}
F(\mathcal{F})
& =
\colim_{J' \subset J\text{，且为有限集}}
F(\bigoplus\nolimits_{j \in J'} \mathcal{H}_j) \\
& =
\colim_{J' \subset J\text{，且为有限集}}
\bigoplus\nolimits_{j \in J'} F(\mathcal{H}_j) \\
& =
\bigoplus\nolimits_{j \in J} F(\mathcal{H}_j)
\end{align*}
因此 $F$ 与任意直和交换。

\medskip\noindent
本段证明该延拓右正合。假设
$0 \to \mathcal{F} \to \mathcal{F}' \to \mathcal{F}'' \to 0$
是拟凝聚 $\mathcal{O}_\mathcal{X}$-模的短正合列。把
$\mathcal{F}' = \bigcup \mathcal{F}'_i$ 写成其凝聚子模之并
（见以上所给参考文献）。
% P12-ERR-0136: see ../control/ERRATA_CANDIDATES.jsonl
用 $\mathcal{F}''_i \subset \mathcal{F}''$ 表示 $\mathcal{F}'_i$ 的像，
并置 $\mathcal{F}_i = \mathcal{F} \cap \mathcal{F}'_i =
\Ker(\mathcal{F}'_i \to \mathcal{F}''_i)$。于是显然有
$\mathcal{F} = \bigcup \mathcal{F}_i$ 和
$\mathcal{F}'' = \bigcup \mathcal{F}''_i$，并有短正合列
$$
0 \to \mathcal{F}_i \to \mathcal{F}_i' \to \mathcal{F}_i'' \to 0
$$
由于该延拓与滤过余极限交换，有
$F(\mathcal{F}) = \colim_{i \in I} F(\mathcal{F}_i)$、
$F(\mathcal{F}') = \colim_{i \in I} F(\mathcal{F}'_i)$ 和
$F(\mathcal{F}'') = \colim_{i \in I} F(\mathcal{F}''_i)$。
% P12-ERR-0137: see ../control/ERRATA_CANDIDATES.jsonl
由于模层的滤过余极限正合，可知 $F$ 的延拓右正合。

\medskip\noindent
证明完成，因为与所有余积交换的右正合函子与所有余极限交换；参见
《范畴论》中的引理
\ref{categories-lemma-colimits-coproducts-coequalizers}。
\end{proof}

\begin{lemma}
\label{stacks-more-morphisms-lemma-affine}
设 $\mathcal{X}$ 是对角态射仿射的代数叠。设 $B$ 是环。设
$F : \QCoh(\mathcal{O}_\mathcal{X}) \to \text{Mod}_B$
是与直和交换的右正合张量函子。设
$g : U \to \mathcal{X}$ 是态射，其中 $U = \Spec(A)$ 仿射。那么
\begin{enumerate}
\item $C = F(g_{\QCoh, *}\mathcal{O}_U)$ 是交换 $B$-代数；
\item 存在环映射 $A \to C$；
\end{enumerate}
使得 $F \circ g_{\QCoh, *} : \text{Mod}_A \to \text{Mod}_B$
把 $M$ 映为视作 $B$-模的 $M \otimes_A C$。
\end{lemma}

\begin{proof}
注意 $g$ 拟紧且拟分离；参见《叠的态射》中的引理
\ref{stacks-morphisms-lemma-quasi-compact-quasi-separated-permanence}。
在《叠的上同调》中的命题
\ref{stacks-cohomology-proposition-direct-image-quasi-coherent} 中，
已经构造函子
$g_{\QCoh, *} : \QCoh(\mathcal{O}_U) \to
\QCoh(\mathcal{O}_\mathcal{X})$。由《叠的上同调》中的诸备注
\ref{stacks-cohomology-remark-direct-image-quasi-coherent-tensor} 和
\ref{stacks-cohomology-remark-QCoh-tensor}，得到乘法
$$
\mu :
g_{\QCoh, *}\mathcal{O}_U
\otimes_{\mathcal{O}_\mathcal{X}}
g_{\QCoh, *}\mathcal{O}_U
\longrightarrow
g_{\QCoh, *}\mathcal{O}_U
$$
它使 $g_{\QCoh, *}\mathcal{O}_U$ 成为交换
$\mathcal{O}_\mathcal{X}$-代数。因此
$C = F(g_{\QCoh, *}\mathcal{O}_U)$ 是 $\text{Mod}_B$ 中的交换代数对象；
换言之，$C$ 是交换 $B$-代数。注意，有映射
$\kappa : A \to
\text{End}_{\mathcal{O}_\mathcal{X}}(g_{\QCoh, *}\mathcal{O}_U)$，
使得对任意 $a \in A$，图表
% P12-ERR-0138: see ../control/ERRATA_CANDIDATES.jsonl
$$
\xymatrix{
g_{\QCoh, *}\mathcal{O}_U \otimes_{\mathcal{O}_\mathcal{X}}
g_{\QCoh, *}\mathcal{O}_U
\ar[d]_{\kappa(r) \otimes 1} \ar[rr]_-\mu & &
g_{\QCoh, *}\mathcal{O}_U \ar[d]^{\kappa(r)} \\
g_{\QCoh, *}\mathcal{O}_U \otimes_{\mathcal{O}_\mathcal{X}}
g_{\QCoh, *}\mathcal{O}_U
\ar[rr]^-\mu & &
g_{\QCoh, *}\mathcal{O}_U
}
$$
交换。由此得到映射
$\kappa' = F(\kappa) : A \to \text{End}_B(C)$，满足
$\kappa'(a)(c) c' = \kappa'(a)(cc')$；当然，这意味着
$a \mapsto \kappa'(a)(1)$ 是环映射 $A \to C$。

\medskip\noindent
态射 $g : U \to \mathcal{X}$ 仿射；参见《叠的态射》中的引理
\ref{stacks-morphisms-lemma-affine-permanence}。因此，由《叠的上同调》中的
引理 \ref{stacks-cohomology-lemma-quasi-coherent-pushforward-affine}，
$g_{\QCoh, *}$ 正合且与直和交换。所以
$F \circ g_{\QCoh, *} : \text{Mod}_A \to \text{Mod}_B$
是与直和交换的右正合函子，并且把 $A$ 映为 $C$。由《函子与态射》中的
引理 \ref{functors-lemma-functor}，函子 $F \circ g_{\QCoh, *}$
把 $A$-模 $M$ 映为视作 $B$-模的 $M \otimes_A C$。
\end{proof}

\begin{lemma}
\label{stacks-more-morphisms-lemma-universally-injective}
记号如引理 \ref{stacks-more-morphisms-lemma-affine}。假设 $\mathcal{X}$ 是 Noether 的，
并且 $g$ 满且平坦。那么 $B \to C$ 万有单射。
\end{lemma}

\begin{proof}
考虑 $\QCoh(\mathcal{O}_\mathcal{X})$ 中的自然映射
$1 : \mathcal{O}_\mathcal{X} \to g_{\QCoh, *}\mathcal{O}_U$。
拉回到 $U$ 并使用伴随，可知复合
$$
\mathcal{O}_U =
g^*\mathcal{O}_\mathcal{X} \xrightarrow{g^*1}
g^*g_{\QCoh, *}\mathcal{O}_U \to \mathcal{O}_U
$$
是 $\QCoh(\mathcal{O}_U)$ 中的恒等态射。把
$g_{\QCoh, *}\mathcal{O}_U = \colim \mathcal{F}_i$ 写成凝聚
$\mathcal{O}_\mathcal{X}$-模的滤过余极限；参见《叠的上同调》中的引理
\ref{stacks-cohomology-lemma-directed-colimit-coherent}。当 $i$ 足够大时，
映射 $1 : \mathcal{O}_\mathcal{X} \to g_{\QCoh, *}\mathcal{O}_U$
经由 $\mathcal{F}_i$ 分解；参见《叠的上同调》中的引理
\ref{stacks-cohomology-lemma-finite-presentation-quasi-compact-colimit}。
设 $s : \mathcal{O}_\mathcal{X} \to \mathcal{F}_i$ 为该分解。那么
$$
\mathcal{O}_U \xrightarrow{g^*s} g^*\mathcal{F}_i \to
g^*g_{\QCoh, *}\mathcal{O}_U \to \mathcal{O}_U
$$
是恒等态射。换言之，$s$ 拉回到 $U$ 后成为一个直和项的嵌入。
按《叠的上同调》中的引理
\ref{stacks-cohomology-lemma-internal-hom-fp-into-qcoh} 的记号，置
$\mathcal{F}_i^\vee = hom(\mathcal{F}_i, \mathcal{O}_\mathcal{X})$。
特别地，存在求值映射
$ev : \mathcal{F}_i \otimes_{\mathcal{O}_\mathcal{X}} \mathcal{F}_i^\vee
\to \mathcal{O}_\mathcal{X}$。在 $s$ 处求值定义映射
$s^\vee : \mathcal{F}_i^\vee \to \mathcal{O}_\mathcal{X}$。
与关于 $s$ 的陈述对偶，可知 $g^*(s^\vee)$ 满；关于 $hom$ 和
$\otimes$ 与限制到 $U$ 的相容性，参见
% P12-ERR-0139: see ../control/ERRATA_CANDIDATES.jsonl
《叠的上同调》第 \ref{stacks-cohomology-section-further-remarks} 节。
由于 $g$ 满且平坦，可知 $s^\vee$ 满（见同一出处）。由于 $F$ 右正合，
可知 $F(\mathcal{F}_i^\vee) \to F(\mathcal{O}_\mathcal{X}) = B$ 满。
选取 $\lambda \in F(\mathcal{F}_i^\vee)$，使它映到 $1 \in B$。
% P12-ERR-0140: see ../control/ERRATA_CANDIDATES.jsonl
用 $e = F(s)(1) \in F(\mathcal{F}_i)$ 表示 $1$ 在映射
$F(s) : B = F(\mathcal{O}_\mathcal{X}) \to F(\mathcal{F}_i)$ 下的像。
那么映射
$$
F(ev) :
F(\mathcal{F}_i) \otimes_B F(\mathcal{F}_i^\vee) =
F(\mathcal{F}_i \otimes_{\mathcal{O}_\mathcal{X}} \mathcal{F}_i^\vee)
\longrightarrow
F(\mathcal{O}_\mathcal{X}) = B
$$
按照构造把 $e \otimes \lambda$ 映到 $1$。因此映射
$B \to F(\mathcal{F}_i)$、$b \mapsto be$ 万有单射，因为它有单侧逆
$F(\mathcal{F}_i) \to B$、$\xi \mapsto F(ev)(\xi \otimes \lambda)$。
这对所有足够大的 $i$ 都成立，故得结论。
\end{proof}

\begin{lemma}
\label{stacks-more-morphisms-lemma-flat}
设 $B \to C$ 是环映射。如果
\begin{enumerate}
\item 余投影 $C \to C \otimes_B C$ 平坦；
\item $B \to C$ 万有单射，
\end{enumerate}
那么 $B \to C$ 忠实平坦。
\end{lemma}

\begin{proof}
由于 $B \to C$ 万有单射，映射
$\Spec(C) \to \Spec(B)$ 满。因此只须证明 $B \to C$ 平坦，
而这由《下降》中的定理
\ref{descent-theorem-descend-module-properties} 得出。
\end{proof}

\noindent
下面这个非常简单的 K\"unneth 版本在我们写出关于代数叠上拟凝聚模上同调
的 K\"unneth 定理一节后，
% P12-ERR-0141: see ../control/ERRATA_CANDIDATES.jsonl
应当变得过时。

\begin{lemma}
\label{stacks-more-morphisms-lemma-easy-kunneth}
设 $a : \mathcal{Y} \to \mathcal{X}$ 和
$b : \mathcal{Z} \to \mathcal{X}$ 可由概形表示、拟紧、拟分离且平坦。
那么
$a_{\QCoh, *}\mathcal{O}_\mathcal{Y}
\otimes_{\mathcal{O}_\mathcal{X}}
b_{\QCoh, *}\mathcal{O}_\mathcal{Z} =
f_{\QCoh, *}\mathcal{O}_{\mathcal{Y} \times_\mathcal{X} \mathcal{Z}}$，
其中 $f : \mathcal{Y} \times_\mathcal{X} \mathcal{Z} \to \mathcal{X}$
是显然的态射。
\end{lemma}

\begin{proof}
简记 $\mathcal{P} = \mathcal{Y} \times_\mathcal{X} \mathcal{Z}$。
由于 $a \circ \text{pr}_1 = f$ 且 $b \circ \text{pr}_2 = f$，
得到映射 $a_*\mathcal{O}_\mathcal{Y} \to f_*\mathcal{O}_\mathcal{P}$ 和
$b_*\mathcal{O}_\mathcal{Z} \to f_*\mathcal{O}_\mathcal{P}$
（使用相对拉回映射；参见《位点》第
\ref{sites-section-pullback} 节）。因此得到相对杯积
$$
\mu :
a_*\mathcal{O}_\mathcal{Y}
\otimes_{\mathcal{O}_\mathcal{X}}
b_*\mathcal{O}_\mathcal{Z} \longrightarrow
f_*\mathcal{O}_{\mathcal{Y} \times_\mathcal{X} \mathcal{Z}}
$$
应用 $Q$ 及其与张量积的相容性（《叠的上同调》中的备注
\ref{stacks-cohomology-remark-QCoh-tensor}），得到
$\QCoh(\mathcal{O}_\mathcal{X})$ 中的箭头
$Q(\mu) : a_{\QCoh, *}\mathcal{O}_\mathcal{Y}
\otimes_{\mathcal{O}_\mathcal{X}}
b_{\QCoh, *}\mathcal{O}_\mathcal{Z} \to
f_{\QCoh, *}\mathcal{O}_{\mathcal{Y} \times_\mathcal{X} \mathcal{Z}}$。
接着，选取概形 $U$ 和满光滑态射 $U \to \mathcal{X}$。
只须证明 $Q(\mu)$ 到 $U_\etale$ 的限制是同构；参见《叠的上同调》第
\ref{stacks-cohomology-section-further-remarks} 节。继而，由同一节的材料，
只须证明 $\mu$ 到 $U_\etale$ 的限制是同构
（这里使用 $\mu$ 的源和靶都是具有基变换性质的局部拟凝聚模）。
此外，可以在 \'etale 拓扑中计算直像；参见《叠的上同调》中的命题
\ref{stacks-cohomology-proposition-loc-qcoh-flat-base-change}。
最后可见
$a_*\mathcal{O}_\mathcal{Y}|_{U_\etale} =
(V \to U)_{small, *}\mathcal{O}_V$，其中
$V = U \times_\mathcal{X} \mathcal{Y}$。$b_*$ 和 $f_*$ 的情形类似。
因此，结论由概形之间平坦、拟紧、拟分离态射的 K\"unneth 公式得出；
参见《概形的导出范畴》中的引理
\ref{perfect-lemma-kunneth}。
\end{proof}

\begin{lemma}
\label{stacks-more-morphisms-lemma-fully-faithful}
设 $\mathcal{X}$ 是对角态射仿射的代数叠。设 $B$ 是环。设
$f_i : \Spec(B) \to \mathcal{X}$（$i = 1, 2$）是两个态射。设
$t : f_1^* \to f_2^*$ 是张量函子
$f_i^* : \QCoh(\mathcal{O}_\mathcal{X}) \to \text{Mod}_B$ 的同构。
那么存在诱导 $t$ 的 $2$-箭头 $f_1 \to f_2$。
\end{lemma}

\begin{proof}
选取仿射概形 $U = \Spec(A)$ 和满光滑态射
$g : U \to \mathcal{X}$；参见《叠的性质》中的引理
\ref{stacks-properties-lemma-quasi-compact-stack}。由于
$\mathcal{X}$ 的对角态射仿射，可知
$U_i = \Spec(B) \times_{f_i, \mathcal{X}, g} U$ 仿射。
设 $U_i = \Spec(C_i)$。那么 $C_i$ 是引理
\ref{stacks-more-morphisms-lemma-affine} 中利用函子 $F = f_i^*$ 构造的带有环映射
$A \to C_i$ 的 $B$-代数。因此，$t$ 诱导与环映射
$A \to C_1$ 和 $A \to C_2$ 相容的 $B$-代数同构 $C_1 \to C_2$。
换言之，我们有
% P12-ERR-0142: see ../control/ERRATA_CANDIDATES.jsonl
交换图表
$$
\xymatrix{
U_i \ar[r] \ar[d] & U \ar[d]^g \\
\Spec(B) \ar[r]^{f_i} & \mathcal{X}
}
\quad\quad
\xymatrix{
&
U_2 \ar[ld] \ar[d]^{\cong} \ar[rd] \\
\Spec(B) &
U_1 \ar[l] \ar[r] &
U
}
$$
这已经表明，$\mathcal{X}$ 在 $\Spec(B)$ 上的对象 $f_1$ 和 $f_2$
在光滑覆盖 $\{U_1 \to \Spec(B)\}$ 后同构。为证明该同构下降为
$\Spec(B)$ 上 $f_1$ 和 $f_2$ 的同构，必须证明我们的同构
（它来自以上交换图表）与 $U_1 \times_{\Spec(B)} U_1$ 上的下降数据
相容。为此，注意 $U \times_\mathcal{X} U$ 也仿射，存在态射
$g' : U \times_\mathcal{X} U \to \mathcal{X}$，并且
$$
U_i \times_{\Spec(B)} U_i =
\Spec(B) \times_{f_i, \mathcal{X}, g'}
(U \times_\mathcal{X} U)
$$
由此可知，来自同构 $C_1 \to C_2$ 的同构
$C_1 \otimes_B C_1 \to C_2 \otimes_B C_2$ 与态射
$U_i \times_{\Spec(B)} U_i \to U \times_\mathcal{X} U$ 相容。
略去一些细节。
\end{proof}

\begin{lemma}
\label{stacks-more-morphisms-lemma-main}
设 $\mathcal{X}$ 是对角态射仿射的 Noether 代数叠。设 $B$ 是环。
设 $F : \QCoh(\mathcal{O}_\mathcal{X}) \to \text{Mod}_B$
是与直和交换的右正合张量函子。那么 $F$ 来自唯一态射
$\Spec(B) \to \mathcal{X}$。
\end{lemma}

\begin{proof}
选取满光滑态射 $g : U \to \mathcal{X}$，其中
$U = \Spec(A)$ 仿射；参见《叠的性质》中的引理
\ref{stacks-properties-lemma-quasi-compact-stack}。应用引理
\ref{stacks-more-morphisms-lemma-affine}，得到有限型交换 $B$-代数
$C = F(g_{\QCoh, *}\mathcal{O}_U)$ 和环映射 $A \to C$。
由引理 \ref{stacks-more-morphisms-lemma-universally-injective}，环映射 $B \to C$ 万有单射。
考虑代数
$$
C \otimes_B C =
F(g_{\QCoh, *}\mathcal{O}_U
\otimes_{\mathcal{O}_\mathcal{X}}
g_{\QCoh, *}\mathcal{O}_U)
$$
由于 $g$ 平坦、拟紧且拟分离，由引理
\ref{stacks-more-morphisms-lemma-easy-kunneth}，下式第一个等号成立：
$$
g_{\QCoh, *}\mathcal{O}_U
\otimes_{\mathcal{O}_\mathcal{X}}
g_{\QCoh, *}\mathcal{O}_U
=
f_{\QCoh, *}\mathcal{O}_{U \times_\mathcal{X} U} =
g_{\QCoh, *}(\text{pr}_{2, *}\mathcal{O}_{U \times_\mathcal{X} U})
$$
其中 $f : U \times_\mathcal{X} U \to \mathcal{X}$ 是显然的态射，
而 $\text{pr}_2 : U \times_\mathcal{X} U \to U$ 是第二投影。
第二个等号由《叠的上同调》中的引理
\ref{stacks-cohomology-lemma-relative-leray} 及
$f = g \circ \text{pr}_2$ 得出。由于 $\mathcal{X}$ 的对角态射仿射，
可知 $U \times_\mathcal{X} U = \Spec(R)$ 仿射。用
$\text{pr}_2 : A \to R$ 把 $R$ 视为 $A$-代数。综上，得到
$$
C \otimes_B C =
F(g_{\QCoh, *}\mathcal{O}_U
\otimes_{\mathcal{O}_\mathcal{X}}
g_{\QCoh, *}\mathcal{O}_U) =
F(g_{\QCoh, *}(\text{pr}_{2, *}\mathcal{O}_{U \times_\mathcal{X} U})) =
R \otimes_A C
$$
其中最后一个等号由引理 \ref{stacks-more-morphisms-lemma-affine} 的末一断言得出。
由于 $A \to R$ 平坦（因为 $\text{pr}_2$ 作为
$U \to \mathcal{X}$ 的基变换而平坦），可知 $C \otimes_B C$ 在
$C$ 上平坦。由引理 \ref{stacks-more-morphisms-lemma-flat}，可知 $B \to C$ 忠实平坦。

\medskip\noindent
我们断言存在实线交换图表
$$
\xymatrix{
\Spec(C \otimes_B C) \ar@<1ex>[d] \ar@<-1ex>[d] \ar[r] &
U \times_\mathcal{X} U \ar@<1ex>[d] \ar@<-1ex>[d] \\
\Spec(C) \ar[d] \ar[r] &
U \ar[d] \\
\Spec(B) \ar@{..>}[r] &
\mathcal{X}
}
$$
箭头 $\Spec(C) \to U = \Spec(A)$ 来自引理
\ref{stacks-more-morphisms-lemma-affine} 陈述中的环映射 $A \to C$。类似地，箭头
$\Spec(C \otimes_B C) \to U \times_\mathcal{X} U$
来自环映射 $R \to C \otimes_B C$。为验证上方方形交换，使用引理
\ref{stacks-more-morphisms-lemma-fully-faithful}；略去细节。由命题
\ref{stacks-more-morphisms-proposition-stack-fpqc}，可知得到点线箭头
$\Spec(B) \to \mathcal{X}$。

\medskip\noindent
$F$ 是与该点线箭头的拉回相对应的函子这一陈述，也由此以及引理
\ref{stacks-more-morphisms-lemma-affine} 中的对应陈述显然可见。略去细节。
\end{proof}

\noindent
% P12-ERR-0143: see ../control/ERRATA_CANDIDATES.jsonl
对环 $B$，用 $\text{Mod}^{fg}_B$ 表示有限生成 $B$-模
（亦称有限 $B$-模）的范畴。

\begin{theorem}
\label{stacks-more-morphisms-theorem-main}
设 $\mathcal{X}$ 是对角态射仿射的 Noether 代数叠。设 $B$ 是
Noether 环。设
$F : \text{Coh}(\mathcal{O}_\mathcal{X}) \to \text{Mod}^{fg}_B$
是右正合张量函子。那么 $F$ 来自唯一态射
$\Spec(B) \to \mathcal{X}$。
\end{theorem}

\begin{proof}
由引理 \ref{stacks-more-morphisms-lemma-extend}，可以把 $F$ 唯一延拓为右正合张量函子
$F : \QCoh(\mathcal{O}_\mathcal{X}) \to \text{Mod}_B$，并且它与所有
% P12-ERR-0144: see ../control/ERRATA_CANDIDATES.jsonl
直和交换。于是可以应用引理 \ref{stacks-more-morphisms-lemma-main}。
\end{proof}
